\chapter{Statistics for engineers}
\label{vol01:ch08}

\epigraph{Statistics is the science of learning from data, and of
measuring, controlling, and communicating uncertainty.}{Marquardt's
working definition, paraphrased and adopted as a chapter motto in
Wasserman's \emph{All of Statistics} \cite[p.~vii]{text:wasserman2004}.}

\chapmeta{Half-life: foundational. Archetypes: uncertainty (deepened
from Chapter 4), failure (statistical malpractice).
Exercise target: 35. Project track: analysis with optional household
component (hybrid). Hazard class: none for analysis, low if the
chosen experiment uses household instruments. Reader path default: standard,
with sections explicitly tagged. Named cases: none in this chapter.
Replication-crisis discussion in section 8.8 stays generic per the
chapter's editorial brief.}

\noindent
A measurement has an uncertainty; a population of measurements has a
distribution; an inference has a power, a coverage, and a way of
being wrong. Chapter 4 made the first quantity operational. A
statistical report earns trust when it states the data summary, the
distributional model, the inference from sample to population, and
the analytical practices it refused.

\section{Summary statistics: mean, median, mode, variance, IQR}\pathtag{core}

A sample is a set of $n$ measurements $x_{1}, x_{2}, \ldots, x_{n}$
drawn from an underlying population by a stated sampling process.
The summary statistics compress the sample into a few numbers that
describe its centre and spread \cite[ch.~6]{text:wasserman2004}. The
compression is lossy. The discipline is to state which information
the chosen summary discards.

\subsection{Centre: mean, median, mode}

The \engterm{sample mean} is

\[
\bar{x} = \frac{1}{n}\sum_{i=1}^{n} x_{i}.
\]

The mean is the centre of mass of the sample. It uses every value.
Its strength is sensitivity: a single large reading shifts the mean
in the direction of that reading. Its weakness is the same: a single
contaminated reading shifts the mean even when the rest of the
sample agrees on a different centre.

The \engterm{sample median} is the middle value of the sorted
sample, with the convention that for even $n$ the median is the
average of the two middle values. The median is the value that
splits the sample into halves of equal count. Its strength is
resistance to outliers: fewer than half the sample can be moved to
extreme values before the median is forced to follow. Its weakness
is information: the median discards the magnitude of values away
from the centre.

The \engterm{sample mode} is the most frequent value, or for
continuous data, the location of the highest density. The mode is
useful for categorical or discrete data and for distributions with a
clear peak; for continuous engineering measurements with finite
sample size, the mode is usually estimated from a histogram or a
kernel density estimate, both of which have their own bin-width or
bandwidth choices.

For a Gaussian distribution all three coincide. For an asymmetric
distribution they separate. The Pearson rule of thumb places the
median between the mode and the mean when the distribution is
skewed: skewed right (long upper tail), the order is mode, median,
mean from left to right; skewed left, the order reverses. The
working consequence is that quoting only one number for a skewed
sample loses the asymmetry the data is trying to report.

\subsection{Spread: variance, standard deviation, range, IQR}

The \engterm{sample variance} is

\[
s^{2} = \frac{1}{n - 1}\sum_{i=1}^{n}(x_{i} - \bar{x})^{2}.
\]

The denominator $n - 1$ rather than $n$ is the Bessel correction. It
makes $s^{2}$ an unbiased estimator of the population variance
$\sigma^{2}$ when the sample mean is used in place of the unknown
population mean. The reader will see $1/n$ in some textbooks; both
estimators are in use, and the choice matters when $n$ is small. For
$n \geq 30$ the difference is below five percent and rarely
load-bearing.

The \engterm{sample standard deviation} $s = \sqrt{s^{2}}$ has the
same units as the data. It is the canonical spread statistic when the
distribution is roughly Gaussian; for heavy-tailed or strongly
skewed distributions it is misleading.

The \engterm{range} is $x_{\max} - x_{\min}$. It uses only two
numbers and is therefore both quick and statistically poor: each
additional sample can only widen it, and a single outlier dominates.
The range is acceptable for a quick sketch and for sample sizes of
two or three; beyond that it wastes data.

The \engterm{interquartile range} (IQR) is $Q_{3} - Q_{1}$, where
$Q_{1}$ is the first quartile (the value below which a quarter of
the sample lies) and $Q_{3}$ is the third. The IQR is the spread of
the middle half of the data. It resists isolated extremes better
than the range, but it is not immune to changes in the lower or
upper quarter of the sample.

A common rule for outlier flagging uses the IQR: a point is an
outlier candidate if it lies more than $1.5 \cdot \text{IQR}$ below
$Q_{1}$ or above $Q_{3}$. The rule is convention. Statistics can
flag the point; domain knowledge decides whether it is
contamination, a real tail observation, or a rare event.

\subsection{Boxplot, histogram, scatter}

Three plots cover most of what an engineer needs to see in a sample
before fitting anything.

\engterm{Boxplot}: a box from $Q_{1}$ to $Q_{3}$, a line at the
median, whiskers to $\min(x)$ and $\max(x)$ within $1.5 \cdot
\text{IQR}$ of the box edges, and dots for points beyond the
whiskers. The boxplot reports centre, spread, asymmetry, and
extreme values in a single graphic.

\engterm{Histogram}: a bar chart of binned counts. Bin width is a
working choice; too narrow and the histogram is noisy, too wide and
features are smoothed away. The Sturges rule ($k = \lceil \log_{2} n
\rceil + 1$) and the Freedman-Diaconis rule (bin width $\approx 2
\cdot \text{IQR} / n^{1/3}$) are common starting points. The
discipline is to try at least two bin widths before believing any
shape claim.

\engterm{Scatter plot}: pairs $(x_{i}, y_{i})$. The first thing the
reader does with two-variable data is plot it. Many regression
problems are diagnosed in the scatter plot before any model is fit.

\subsection{Worked example}

A reader times their morning commute on $30$ working days. The
sample is, in minutes:

\[
\begin{aligned}
&27, 28, 29, 26, 32, 30, 27, 28, 29, 28,\\
&31, 27, 26, 28, 29, 30, 28, 27, 32, 28,\\
&29, 27, 30, 28, 31, 29, 28, 27, 28, 35.
\end{aligned}
\]

Using the printed sample, the sample mean is $\bar{x} \approx 28.73
\,\text{min}$. The median, after sorting, is $28 \,\text{min}$. The
sample standard deviation is $s \approx 1.98 \,\text{min}$. The
range is $35 - 26 = 9 \,\text{min}$. Using the linear percentile
convention, the first and third quartiles are $Q_{1} = 27.25
\,\text{min}$ and $Q_{3} = 29.75 \,\text{min}$, giving an IQR of
$2.50 \,\text{min}$. A different quartile convention (for instance,
Tukey hinges) would give slightly different values; the analyst's
job is to state the convention and use it consistently.

Three observations follow. The mean is about three quarters of a
minute above the median, suggesting a small right skew driven by
the $35$-minute day. The IQR is markedly smaller than the range,
reflecting that one or two long days inflate the range without much
affecting the central half. For a Gaussian distribution,
$\text{IQR}/1.349$ estimates $\sigma$; here that comparison is only
a rough check because the sample is small and the upper tail
contains the $35$-minute day. The reader who reports only the mean
has lost the information that several working days are longer than
the central cluster.

\begin{archetype}[Uncertainty, deepened from Chapter 4]
Chapter 4 carried uncertainty as a single number $u(x)$ beside each
measurement. This chapter asks where $u(x)$ comes from when the
underlying mechanism is variation across a sample rather than a
single measurement's instrument noise. The summary statistics
$\bar{x}$ and $s$ are estimators: $\bar{x}$ estimates the population
mean $\mu$, and $s$ estimates the population standard deviation
$\sigma$. The standard uncertainty of the mean, $u(\bar{x}) =
s/\sqrt{n}$, is the bridge between Chapter 4's framework and this
chapter's. The bridge is sturdy when the sample is independent and
identically distributed; it bends when either condition fails.
\end{archetype}

\section{Distributions: normal, log-normal, Weibull, Poisson, exponential}\pathtag{core}

A sample is a finite window onto a population. The population is
modelled by a probability distribution. For continuous data, the
density $p(x)$ assigns probability to intervals through area under
the curve; for discrete data, the probability mass function assigns
probability to values. Five distributions are common enough that
the reader should recognize their mechanisms before fitting data
\cite[ch.~3-4]{text:devore2015}. We state each in the form a working
engineer needs: its shape, its parameters, the mechanism that
produces it, and the engineering quantities it commonly models.
\Cref{fig:vol01:ch08:distribution-menagerie} collects the five
together with Student's $t$, the small-sample correction the
chapter uses throughout sections 8.3 and 8.4.

\begin{figure}[p]
\centering
\begin{tikzpicture}[x=1cm, y=1cm, font=\small]

% Six-panel distribution menagerie, two rows of three.
% Each panel: thin axes, labelled tick, density curve, name + key
% parameter annotation. Independent local coordinate per panel.

% ---- Panel 1: Normal N(0, 1) ----
\begin{scope}[shift={(0,0)}]
  \node[anchor=south, font=\small\bfseries] at (2.0, 3.2) {Normal};
  \node[anchor=south, font=\scriptsize] at (2.0, 2.8)
    {$\mathcal{N}(0, 1)$};
  \draw[->, thick] (-0.1, 0) -- (4.1, 0);
  \draw[->, thick] (0, -0.1) -- (0, 2.7);
  \draw[thick, engblue, smooth]
    plot[domain=0.05:3.95, samples=60]
    (\x, {2.4*exp(-((\x-2.0)*(\x-2.0))/(2*0.55*0.55))});
  \draw[dashed, enggray] (2.0, 0) -- (2.0, 2.4);
  \node[anchor=north, font=\scriptsize] at (2.0, -0.1) {$\mu$};
\end{scope}

% ---- Panel 2: Log-normal, mu_log = 0, sigma_log = 0.6 ----
\begin{scope}[shift={(5.0,0)}]
  \node[anchor=south, font=\small\bfseries] at (2.0, 3.2) {Log-normal};
  \node[anchor=south, font=\scriptsize] at (2.0, 2.8)
    {$\mu=0,\ \sigma=0.6$};
  \draw[->, thick] (-0.1, 0) -- (4.1, 0);
  \draw[->, thick] (0, -0.1) -- (0, 2.7);
  % approximate log-normal density on [0.1, 4]
  \draw[thick, engred, smooth]
    plot[domain=0.15:3.95, samples=70]
    (\x, {2.6*exp(-pow(ln(\x*0.7),2)/(2*0.5*0.5))/(\x*0.7)});
  \node[anchor=north, font=\scriptsize] at (1.4, -0.1) {median};
\end{scope}

% ---- Panel 3: Weibull, beta = 2, eta = 2 ----
\begin{scope}[shift={(10.0,0)}]
  \node[anchor=south, font=\small\bfseries] at (2.0, 3.2) {Weibull};
  \node[anchor=south, font=\scriptsize] at (2.0, 2.8)
    {$\beta=2,\ \eta=2$};
  \draw[->, thick] (-0.1, 0) -- (4.1, 0);
  \draw[->, thick] (0, -0.1) -- (0, 2.7);
  \draw[thick, engamber, smooth]
    plot[domain=0.05:3.95, samples=70]
    (\x, {3.6*(\x/2.0)*exp(-(\x/2.0)*(\x/2.0))});
  \draw[dashed, enggray] (2.0, 0) -- (2.0, 2.0);
  \node[anchor=north, font=\scriptsize] at (2.0, -0.1) {$\eta$};
\end{scope}

% ---- Panel 4: Poisson, lambda = 4 (PMF as bars) ----
\begin{scope}[shift={(0,-4.5)}]
  \node[anchor=south, font=\small\bfseries] at (2.0, 3.2) {Poisson};
  \node[anchor=south, font=\scriptsize] at (2.0, 2.8)
    {$\lambda=4$};
  \draw[->, thick] (-0.1, 0) -- (4.1, 0);
  \draw[->, thick] (0, -0.1) -- (0, 2.7);
  % Hand-tabulated PMF for lambda=4
  \foreach \k/\p in {0/0.018, 1/0.073, 2/0.147, 3/0.195, 4/0.195,
                     5/0.156, 6/0.104, 7/0.060, 8/0.030, 9/0.013} {
    \pgfmathsetmacro\xpos{0.2 + \k*0.38}
    \pgfmathsetmacro\height{\p*12.0}
    \draw[fill=engblue!60, draw=engblue]
      (\xpos, 0) rectangle (\xpos+0.32, \height);
  }
  \node[anchor=north, font=\scriptsize] at (2.0, -0.1) {$k$ events};
\end{scope}

% ---- Panel 5: Exponential, tau = 1 ----
\begin{scope}[shift={(5.0,-4.5)}]
  \node[anchor=south, font=\small\bfseries] at (2.0, 3.2) {Exponential};
  \node[anchor=south, font=\scriptsize] at (2.0, 2.8)
    {$\tau=1$};
  \draw[->, thick] (-0.1, 0) -- (4.1, 0);
  \draw[->, thick] (0, -0.1) -- (0, 2.7);
  \draw[thick, engred, smooth]
    plot[domain=0.0:3.95, samples=60]
    (\x, {2.4*exp(-\x*1.0)});
  \node[anchor=north, font=\scriptsize] at (2.0, -0.1) {waiting time};
\end{scope}

% ---- Panel 6: Student's t with nu = 5, overlaid Gaussian ----
\begin{scope}[shift={(10.0,-4.5)}]
  \node[anchor=south, font=\small\bfseries] at (2.0, 3.2) {Student $t$};
  \node[anchor=south, font=\scriptsize] at (2.0, 2.8)
    {$\nu = 5$ vs $\mathcal{N}(0,1)$};
  \draw[->, thick] (-0.1, 0) -- (4.1, 0);
  \draw[->, thick] (0, -0.1) -- (0, 2.7);
  % Gaussian reference
  \draw[thick, engblue!50, dashed, smooth]
    plot[domain=0.1:3.9, samples=60]
    (\x, {2.4*exp(-((\x-2.0)*(\x-2.0))/(2*0.55*0.55))});
  % t-like: slightly heavier tails (1+x^2/5)^-3
  \draw[thick, engred, smooth]
    plot[domain=0.1:3.9, samples=60]
    (\x, {2.2*pow(1+((\x-2.0)*(\x-2.0))/(5*0.55*0.55), -3)});
  \draw[dashed, enggray] (2.0, 0) -- (2.0, 2.4);
  \node[anchor=north, font=\scriptsize] at (2.0, -0.1) {standardised};
\end{scope}

\end{tikzpicture}
\caption[The five common distributions plus Student's $t$]{The
distributions an engineer should recognise on sight, together with
the parameter that controls each shape. The normal density is
symmetric and centred at $\mu$; the log-normal is right-skewed and
positive; the Weibull at $\beta = 2$ rises and falls about a scale
$\eta$ that fixes the $63.2\,\%$ failure point; the Poisson is
discrete with mass concentrated near $\lambda$; the exponential is
monotone decreasing with mean $\tau$; the Student $t$ with low
degrees of freedom has heavier tails than the Gaussian it tends to
as $\nu \to \infty$. The shapes are stylised for recognition, not
plotted to a common probability scale.}
\label{fig:vol01:ch08:distribution-menagerie}
\end{figure}


\subsection{Normal (Gaussian)}

The normal distribution $\mathcal{N}(\mu, \sigma^{2})$ has density

\[
p(x) = \frac{1}{\sigma\sqrt{2\pi}} \exp\!\left(-\frac{(x - \mu)^{2}}
{2\sigma^{2}}\right).
\]

The mean is $\mu$, the variance is $\sigma^{2}$, and the
distribution is symmetric. The total probability under the curve is
unity, a result that follows from the Gaussian normalisation
integral; we quote the result without deriving it. The probability
that a single draw lies within $\pm 1\sigma$ of the mean is about
$0.683$; within $\pm 2\sigma$, about $0.954$; within $\pm 3\sigma$,
about $0.997$. These coverage values match the coverage factors
$k = 1, 2, 3$ from Chapter 4 §4.6.

The normal distribution is the limiting distribution of sums of many
small independent contributions, by the central limit theorem
(Chapter 4 §4.4). Engineering quantities that often follow the
normal distribution include the reading on a digital instrument
operating in its linear range, the displacement of a mechanical
sensor under thermal noise, the timing jitter of a synchronous
clock, and any quantity whose error is the sum of many small,
roughly equal contributions.

The normal distribution is misused when the underlying mechanism is
multiplicative rather than additive (use lognormal), bounded below
at zero (use lognormal, half-normal, or a truncated form), or
heavy-tailed (use a $t$-distribution or a Pareto). A normal fit to
heavy-tailed data underestimates the probability of extreme events,
sometimes by orders of magnitude. The cost of that underestimate
falls on whoever sized the system for the typical case.

\subsection{Log-normal}

A quantity $X$ is \engterm{log-normal} if $\ln X$ is normally
distributed. Equivalently, $X = e^{Y}$ with $Y \sim \mathcal{N}(\mu,
\sigma^{2})$. The density of $X$, for $x > 0$, is

\[
p(x) = \frac{1}{x \sigma \sqrt{2\pi}}
\exp\!\left(-\frac{(\ln x - \mu)^{2}}{2\sigma^{2}}\right).
\]

In this distribution, $\mu$ and $\sigma$ describe $\ln X$; the mean
of $X$ is $e^{\mu + \sigma^{2}/2}$ and the median of $X$ is $e^{\mu}$.
The distribution is right-skewed and bounded below at zero.

The log-normal arises from multiplicative mechanisms: a quantity is
the product of many small independent factors. Examples in
engineering include particle-size distributions in milling and
grinding, file-size distributions in storage and network systems,
positive concentration measurements in chemical processes, asset
prices or multiplicative cost factors under simple financial
models, and biological growth quantities. Many
quantities that look skewed when plotted on linear axes look
roughly Gaussian when plotted on a logarithmic axis; the engineer
who suspects a multiplicative mechanism plots on log axes early.

\subsection{Weibull}

The Weibull distribution is a flexible model for time-to-failure
data. Its survival function (probability that the lifetime exceeds
$t$) is

\[
S(t) = e^{-(t/\eta)^{\beta}}, \quad t \geq 0,
\]

with shape parameter $\beta > 0$ and scale parameter $\eta > 0$. The
shape parameter $\beta$ controls how the failure rate evolves over
time: $\beta < 1$ models a decreasing failure rate (infant
mortality, where defective units fail early and the survivors are
sturdy); $\beta = 1$ reduces to the exponential, modelling a
constant failure rate; $\beta > 1$ models an increasing failure rate
(wear-out, where surviving units accumulate damage). The scale
$\eta$ is the time at which $63.2 \,\%$ of units have failed, which
follows by setting $t = \eta$ in $S(t)$.

The Weibull is the standard reliability distribution in mechanical,
electrical, and aerospace engineering for components whose failure
mode is fatigue, wear, or stochastic damage accumulation. It is
flexible enough to fit many lifetime datasets but fitted parameters
should be checked against the underlying physical mechanism: a
Weibull with $\beta = 0.5$ that is being applied to a wear-out
process is a fitting error, not a feature of the data.
\Cref{fig:vol01:ch08:weibull-hazard-family} shows the hazard rate
$h(t) = (\beta/\eta)(t/\eta)^{\beta - 1}$ for the three regimes;
the shape of the hazard curve, read against the known failure
physics, is the first check on any Weibull fit.

\begin{figure}[ht]
\centering
\begin{tikzpicture}[x=1cm, y=1cm, font=\small]

% Weibull hazard rate h(t) = (beta/eta) (t/eta)^(beta-1) for three
% shape parameters, eta = 1. beta < 1 decreasing (infant mortality),
% beta = 1 constant (exponential), beta > 1 increasing (wear-out).
% Plot on t in [0.05, 2.2]; clamp the vertical extent for beta<1.

% axes
\draw[->, thick] (0, 0) -- (7.2, 0)
  node[below, font=\scriptsize] {age $t / \eta$};
\draw[->, thick] (0, 0) -- (0, 5.2)
  node[left, font=\scriptsize, rotate=90, yshift=0.6cm, anchor=south]
  {hazard rate $h(t) \cdot \eta$};

% x ticks
\foreach \tval/\xc in {0.5/1.5, 1.0/3.0, 1.5/4.5, 2.0/6.0} {
  \draw (\xc, 0) -- (\xc, -0.1) node[below, font=\scriptsize] {\tval};
}
% y ticks
\foreach \hval/\yc in {1/1, 2/2, 3/3, 4/4} {
  \draw (0, \yc) -- (-0.1, \yc) node[left, font=\scriptsize] {\hval};
}

% x plot unit = 3 * (t/eta); y plot unit = 1 * (h*eta), clamped at 5.

% beta = 0.5 (decreasing): h*eta = 0.5 * t^(-0.5). Clamp large values.
\draw[thick, engblue, smooth]
  plot[domain=0.10:2.2, samples=90]
  ({3*\x}, {min(5, 0.5*(\x)^(-0.5))});

% beta = 1 (constant): h*eta = 1.
\draw[thick, engorange] (0.3, 1) -- (6.6, 1);

% beta = 3 (increasing): h*eta = 3 * t^2. Clamp at 5.
\draw[thick, engred, smooth]
  plot[domain=0.10:2.2, samples=90]
  ({3*\x}, {min(5, 3*(\x)*(\x))});

% labels
\node[engblue, anchor=west, font=\scriptsize] at (0.6, 4.6)
  {$\beta = 0.5$ (infant mortality)};
\node[engorange, anchor=west, font=\scriptsize] at (4.2, 1.25)
  {$\beta = 1$ (constant)};
\node[engred, anchor=west, font=\scriptsize] at (3.6, 4.2)
  {$\beta = 3$ (wear-out)};

\end{tikzpicture}
\caption[Weibull hazard rate across shape parameters]{The Weibull
hazard rate $h(t) = (\beta/\eta)(t/\eta)^{\beta - 1}$ for three
shape parameters at fixed scale $\eta$, plotted against
normalised age $t/\eta$. With $\beta < 1$ the hazard falls with
age, the signature of infant mortality where weak units fail early
and survivors are sturdy. With $\beta = 1$ the hazard is constant,
recovering the exponential and its memorylessness. With $\beta > 1$
the hazard climbs with age, the signature of wear-out where
surviving units accumulate damage. A fitted $\beta$ that disagrees
with the known failure physics, a $\beta < 1$ fit to a fatigue
process, is a diagnostic that the sample, the fit, or the
failure-mode model is wrong. The three curves are the three
segments of the classical bathtub reliability curve.}
\label{fig:vol01:ch08:weibull-hazard-family}
\end{figure}


\subsection{Poisson}

The Poisson distribution describes counts of events that occur
independently at a constant average rate. The probability of
observing exactly $k$ events in a fixed interval, given a mean
$\lambda$, is

\[
\Pr(K = k) = \frac{\lambda^{k} e^{-\lambda}}{k!}, \quad k = 0, 1, 2,
\ldots
\]

The mean equals the variance: $\E[K] = \Var(K) = \lambda$. The
standard deviation is $\sqrt{\lambda}$. The Poisson applies when
events are rare on the scale of the observing interval, are
independent, and occur at a stable rate. Examples in engineering
include radiation counts at a detector, defect counts on a
manufactured surface, and packet arrivals at a network interface
during low-traffic intervals. Earthquake counts can be modelled as
Poisson only under a restricted background-rate model that excludes
clustering from aftershocks and other nonstationary effects.

A common diagnostic is the variance-to-mean ratio: a sample whose
ratio is much less than one is under-dispersed (the events are not
independent and the system has some self-correction); a sample
whose ratio is much greater than one is over-dispersed (the events
cluster, often because the rate itself fluctuates). For routine
defect-count data, a variance-to-mean ratio near one is a screening
sign that a Poisson model may be plausible. How far from one is
acceptable depends on sample size, process context, and the cost of
a wrong model.

\subsection{Exponential}

The exponential distribution describes the waiting time between
events in a Poisson process, or any positive random quantity with a
constant per-unit-time hazard. Its density is

\[
p(t) = \frac{1}{\tau} e^{-t/\tau}, \quad t \geq 0,
\]

with mean $\tau$ and standard deviation also $\tau$. The
distribution has the memorylessness property: the probability of a
further wait of $\Delta t$ does not depend on how long the wait has
already lasted. This is a strong assumption and it is the source of
most exponential-fit failures in engineering. Wear-out, ageing,
fatigue, and warm-up effects all violate memorylessness; use Weibull
when the hazard rate changes with age.

The exponential is the right model for inter-arrival times in a
true Poisson process, for time to the first failure of a system
whose failure mechanism is dominated by random external shocks
rather than internal ageing, and for radioactive-decay times of a
single nucleus. The reader will encounter it again in queueing
theory (Volume IX) and in reliability analysis (Volume X).

\subsection{Choosing a distribution}

The engineer chooses a distribution from three considerations:
mechanism, plot, and fit diagnostic.

The mechanism considers which generative process plausibly produced
the data. A Gaussian is plausible when the variation is the sum of
many small independent contributions. A log-normal is plausible when
the variation is the product. A Weibull is plausible when the data
are positive lifetimes with a hazard rate that may vary with age. A
Poisson is plausible when the data are counts of rare independent
events. An exponential is plausible when the data are waiting times
in a memoryless process.

The plot reveals the shape. A normal-probability plot (a Q-Q plot
against the normal distribution) makes deviations from Gaussian
visible. A log-scale histogram makes deviations from log-normal
visible. A Weibull plot, which uses a transformation that linearises
$\ln(-\ln S(t))$ against $\ln t$, makes Weibull data fall on a
straight line. \Cref{fig:vol01:ch08:qq-plot} contrasts the Q-Q
signature of a Gaussian sample with that of a heavy-tailed sample;
the engineer who learns to read this plot saves hours of mistaken
fits.

\begin{figure}[ht]
\centering
\begin{tikzpicture}[x=1cm, y=1cm, font=\small]

% Two side-by-side panels: Gaussian sample (points on reference line)
% and heavy-tailed sample (points deviating in both tails).

% ---- Panel A: Gaussian sample ----
\begin{scope}[shift={(0,0)}]
  \draw[->, thick] (-0.2, 0) -- (4.2, 0)
    node[below, font=\scriptsize] {theoretical $z$};
  \draw[->, thick] (0, -0.2) -- (0, 4.2)
    node[left, font=\scriptsize] {observed (sorted)};
  % 45-degree reference line
  \draw[dashed, enggray] (0.2, 0.2) -- (4.0, 4.0);
  % Points approximately on the line
  \foreach \x/\y in {0.4/0.45, 0.8/0.78, 1.1/1.18, 1.5/1.54, 1.8/1.85,
                     2.1/2.05, 2.4/2.45, 2.7/2.68, 3.0/2.95, 3.3/3.30,
                     3.6/3.62} {
    \fill[engblue] (\x, \y) circle (0.06);
  }
  \node[anchor=south, font=\small\bfseries] at (2.0, 4.3)
    {(a) Gaussian sample};
\end{scope}

% ---- Panel B: Heavy-tailed sample ----
\begin{scope}[shift={(5.5,0)}]
  \draw[->, thick] (-0.2, 0) -- (4.2, 0)
    node[below, font=\scriptsize] {theoretical $z$};
  \draw[->, thick] (0, -0.2) -- (0, 4.2)
    node[left, font=\scriptsize] {observed (sorted)};
  % 45-degree reference line
  \draw[dashed, enggray] (0.2, 0.2) -- (4.0, 4.0);
  % Points showing S-shape: low-end values pulled down, high-end up
  \foreach \x/\y in {0.4/0.10, 0.8/0.45, 1.1/0.90, 1.5/1.40, 1.8/1.80,
                     2.1/2.10, 2.4/2.55, 2.7/3.00, 3.0/3.45, 3.3/3.85,
                     3.6/4.10} {
    \fill[engred] (\x, \y) circle (0.06);
  }
  \node[anchor=south, font=\small\bfseries] at (2.0, 4.3)
    {(b) Heavy-tailed sample};
\end{scope}

\end{tikzpicture}
\caption[Q-Q plots against the normal reference]{Q-Q plot of two
samples against a Gaussian reference. The dashed line is the
$y = x$ identity expected when the sample comes from the reference
distribution. Panel (a) shows a $n = 11$ Gaussian sample: the
sorted observations sit close to the line across the full range.
Panel (b) shows a heavy-tailed sample of the same size: the lower
end falls below the line (observed extremes more negative than
Gaussian predicts) and the upper end rises above the line (observed
extremes more positive). The S-shaped deviation diagnoses tails
heavier than Gaussian; a mirror-image S diagnoses lighter tails.}
\label{fig:vol01:ch08:qq-plot}
\end{figure}


The fit diagnostic checks the chosen distribution against the data
quantitatively. A goodness-of-fit test (Kolmogorov-Smirnov,
Anderson-Darling, chi-square) returns a $p$-value for the null
hypothesis that the data come from the proposed distribution. If
parameters were fitted from the same data, the reference
distribution for the test statistic may need adjustment or
simulation; the chi-square test in particular needs adequate
expected counts and defined bins. A $p$-value below the chosen
significance level rejects the model; a $p$-value above fails to
reject it and does not prove the model. The reader should not rely
on the diagnostic alone: a tiny sample produces a non-rejecting test
even for badly mis-specified distributions, and a huge sample
produces a rejecting test even for trivial deviations that do not
change any engineering decision.

\section{Hypothesis testing in engineering context}\pathtag{core}

A hypothesis test asks whether a sample is consistent with a
specified claim about the population it came from
\cite[ch.~10]{text:wasserman2004}. The framework is precise. Its
misuse is also precise, and is the subject of section 8.7.

\subsection{Null and alternative}

The \engterm{null hypothesis} $H_{0}$ is a specific statement about
the population: a parameter equals a value, two populations have
equal means, a process is on target, the rate is unchanged. The
\engterm{alternative hypothesis} $H_{1}$ is the complement, or a
specific direction of deviation. The test is set up to reject $H_{0}$
in favour of $H_{1}$ when the sample is sufficiently inconsistent
with $H_{0}$.

The test produces a \engterm{test statistic}, a function of the
sample whose distribution under $H_{0}$ is known. The observed value
of the test statistic is compared to that distribution. The
\engterm{$p$-value} is the probability, under $H_{0}$, of observing a
test statistic at least as extreme as the one observed. A $p$-value
below a chosen significance level $\alpha$ rejects $H_{0}$.

The convention $\alpha = 0.05$ is a working default in many
engineering and scientific fields. It is convention. The appropriate
$\alpha$ depends on how the null is framed and on the relative cost
of false positives and false negatives. A safety-critical test first
states which error would allow an unsafe decision, then allocates
$\alpha$ and power to control that error.

\subsection{Two errors}

A hypothesis test can be wrong in two ways. A \engterm{Type I
error} (false positive) is rejecting $H_{0}$ when it is true; its
probability is $\alpha$ by construction. A \engterm{Type II error}
(false negative) is failing to reject $H_{0}$ when $H_{1}$ is true;
its probability is $\beta$, and depends on the true effect size, the
sample size, and $\alpha$. The \engterm{power} of the test is
$1 - \beta$: the probability of rejecting $H_{0}$ when $H_{1}$ is
true at a specified effect size.

The engineer designs the test by choosing $\alpha$ from the cost of
a Type I error and choosing $n$ to achieve a target power $1 - \beta$
at the smallest effect size that matters. Section 8.6 takes up
sample size and power as a planning quantity.
\Cref{fig:vol01:ch08:type-i-type-ii} shows the geometry: $\alpha$
is the area of the null tail past the critical value $c$, and
$\beta$ is the area of the alternative density to the left of $c$.
Moving $c$ trades the two areas against each other at fixed $n$;
increasing $n$ narrows both densities and shrinks both areas at
once.

\begin{figure}[ht]
\centering
\begin{tikzpicture}[x=1cm, y=1cm, font=\small]

% Two overlapping Gaussians: null centred at 0, alternative centred at
% 1.6. Shade the alpha-tail under the null (right of critical value c)
% and the beta-region under the alternative (left of c).

\def\xc{2.4}     % critical value on the horizontal scale

% axes
\draw[->, thick] (-2.5, 0) -- (5.5, 0) node[below, font=\scriptsize] {test statistic};
\draw[->, thick] (-2.5, -0.05) -- (-2.5, 3.2);

% Null density: N(0, 1) (centred at horizontal coordinate 0)
\draw[thick, engblue, smooth]
  plot[domain=-2.5:4.0, samples=80]
  (\x, {2.6*exp(-(\x*\x)/(2*1.0*1.0))});

% Alternative density: N(1.6, 1)
\draw[thick, engred, smooth]
  plot[domain=-1.0:5.0, samples=80]
  (\x, {2.6*exp(-((\x-1.6)*(\x-1.6))/(2*1.0*1.0))});

% Type I shaded region: right tail of null beyond critical value
\fill[engblue!35, opacity=0.55]
  plot[domain=\xc:4.0, samples=40]
  (\x, {2.6*exp(-(\x*\x)/(2*1.0*1.0))}) -- (\xc, 0) -- cycle;

% Type II shaded region: left of c under the alternative
\fill[engred!35, opacity=0.55]
  plot[domain=-1.0:\xc, samples=40]
  (\x, {2.6*exp(-((\x-1.6)*(\x-1.6))/(2*1.0*1.0))}) -- (\xc, 0) -- (-1.0, 0) -- cycle;

% Critical value line
\draw[dashed, thick] (\xc, 0) -- (\xc, 3.0);
\node[anchor=south, font=\scriptsize] at (\xc, 3.0) {critical value $c$};

% Centre marks
\draw[dashed, enggray] (0, 0) -- (0, 2.6);
\draw[dashed, enggray] (1.6, 0) -- (1.6, 2.6);
\node[anchor=north, font=\scriptsize] at (0, -0.05) {$\mu_{0}$ (null)};
\node[anchor=north, font=\scriptsize] at (1.6, -0.05) {$\mu_{1}$ (alt)};

% Labels for the two error regions
\node[font=\scriptsize, engblue!60!black, anchor=west]
  at (3.2, 0.6) {$\alpha$: Type~I};
\node[font=\scriptsize, engred!60!black, anchor=east]
  at (2.2, 1.2) {$\beta$: Type~II};

\end{tikzpicture}
\caption[Type I and Type II errors as tail areas]{The geometry of
the two errors a hypothesis test commits. The null distribution
(blue) is centred at $\mu_{0}$; the alternative (red) at $\mu_{1}$.
A test rejects the null when the observed statistic exceeds the
critical value $c$. The shaded blue tail is $\alpha$, the
false-positive rate the test controls by construction. The shaded
red region is $\beta$, the false-negative probability the test must
also be designed against by choosing $n$ large enough that the
alternative density at $c$ leaves only a small remainder to its
left. The power of the test is $1 - \beta$, the unshaded area of
the red density to the right of $c$.}
\label{fig:vol01:ch08:type-i-type-ii}
\end{figure}


\subsection{Common tests}

Three tests cover most engineering work.

\engterm{One-sample t-test}. The sample comes from a population with
unknown mean $\mu$ and unknown standard deviation. The null
hypothesis $H_{0}: \mu = \mu_{0}$ is tested against $H_{1}: \mu \neq
\mu_{0}$ (two-sided) or $H_{1}: \mu > \mu_{0}$ (one-sided). The test
statistic is

\[
t = \frac{\bar{x} - \mu_{0}}{s/\sqrt{n}},
\]

which under $H_{0}$ follows Student's $t$-distribution with $n - 1$
degrees of freedom. The reader will recognise the denominator as the
standard uncertainty of the mean from Chapter 4 §4.4.

\engterm{Two-sample t-test}. Two independent samples have means
$\bar{x}_{1}$ and $\bar{x}_{2}$, sample standard deviations $s_{1}$
and $s_{2}$, and sizes $n_{1}$ and $n_{2}$. The null hypothesis
$H_{0}: \mu_{1} = \mu_{2}$ is tested against an inequality. The test
statistic, in the Welch form that does not assume equal variances,
is

\[
t = \frac{\bar{x}_{1} - \bar{x}_{2}}
{\sqrt{s_{1}^{2}/n_{1} + s_{2}^{2}/n_{2}}}.
\]

The Welch formula uses an adjusted degrees-of-freedom calculation
(the Welch-Satterthwaite formula); statistical software returns it
automatically. The two-sample test is the workhorse for comparing a
new process to an existing one, a treatment to a control, or a
specimen batch to a reference batch.

\engterm{Chi-square goodness-of-fit}. Observed counts $O_{i}$ in
$k$ categories are compared to expected counts $E_{i}$ under a
hypothesised distribution. The test statistic is

\[
\chi^{2} = \sum_{i=1}^{k} \frac{(O_{i} - E_{i})^{2}}{E_{i}},
\]

which under the null follows a chi-square distribution with degrees
of freedom equal to $k$ minus one minus the number of fitted
parameters. The chi-square test is the standard tool for asking
whether a categorical sample (defect-mode counts, traffic-direction
counts, fault-class counts) is consistent with a specified
distribution.

Other tests follow the same pattern: write the null and alternative,
compute a statistic, compare to its known null distribution. The
NIST/SEMATECH e-Handbook \cite{web:nist-eshandbook} catalogues the
common tests with worked examples. Wasserman gives the
distributional theory \cite[ch.~10]{text:wasserman2004}.

A working laboratory's choice of software package matters less than
the discipline of pre-specifying the test, but the choices have
characteristic conventions. In R, the working command is
\texttt{t.test(x, y, alternative = "two.sided", var.equal = FALSE,
conf.level = 0.95)} which defaults to the Welch correction and
returns the test statistic, degrees of freedom (Welch-Satterthwaite
adjusted), $p$-value, and confidence interval on the difference of
means. In Python's SciPy, \texttt{scipy.stats.ttest\_ind(x, y,
equal\_var=False)} returns the test statistic and two-sided
$p$-value; the confidence interval requires \texttt{statsmodels} or
manual calculation. In MATLAB, \texttt{ttest2(x, y, 'Vartype',
'unequal')} returns the test decision, $p$-value, and confidence
interval. Each implementation has subtly different defaults: SciPy
defaults to equal-variance (\texttt{equal\_var=True}) prior to
version 1.6, an opt-in Welch convention thereafter; R's
\texttt{t.test} defaults to Welch from the beginning; MATLAB's
\texttt{ttest2} defaults to equal-variance. A laboratory pipeline
that switches between packages must check the variance assumption
at every interface or it will see test results change without
explanation.

For the chi-square goodness-of-fit test, the working alternative
when expected counts are small ($E_{i} < 5$) is Fisher's exact test
for $2 \times 2$ tables or the Monte Carlo $p$-value approximation
for larger tables. R's \texttt{chisq.test(simulate.p.value = TRUE,
B = 10000)} runs the Monte Carlo version; SciPy's
\texttt{scipy.stats.chi2\_contingency(observed, lambda\_ = "log-
likelihood")} offers the log-likelihood (G-test) alternative. The
log-likelihood ratio test is preferred over Pearson chi-square
when the table is sparse, the cell counts are non-integer (a
weighted sample), or the null is composite (parameters fitted from
the same data); the asymptotic distributions of the two statistics
are identical, but the small-sample behaviour differs.

\subsection{What a $p$-value is and is not}

The $p$-value is the probability, under $H_{0}$ and the specified
alternative, of a test statistic at least as extreme as the one
observed. It answers one narrow question: how often the null model
would produce a statistic this extreme. The truth of $H_{0}$, the
probability of replication, the effect size, and the engineering
importance require other evidence.

These misreadings are common and load-bearing. Section 8.7 develops
them as failure modes; for now, the working engineer reads a
$p$-value as: \emph{if the population were as the null says, we would
see data this extreme this often by chance}. A small $p$-value is
evidence that the data are unlikely under the null; whether that
evidence is decisive depends on the prior plausibility of the null,
the effect size, and the cost of the two errors.

\subsection{Engineering examples}

Three examples ground the framework.

A process-control engineer wants to know whether the mean fill
weight of a production line equals the nominal $500 \,\si{\gram}$.
A sample of $n = 25$ bottles has $\bar{x} = 498.7 \,\si{\gram}$ and
$s = 3.2 \,\si{\gram}$. The one-sample $t$-statistic is $t = (498.7
- 500)/(3.2/5) = -2.03$; with $24$ degrees of freedom, the
two-sided $p$-value is about $0.054$. At $\alpha = 0.05$ the engineer
fails to reject the null; at $\alpha = 0.10$ they reject it. The
engineering question is which threshold the regulatory or contractual
context demands.

A quality engineer compares two suppliers' bearing diameters. Sample
$1$ has $n_{1} = 30$, $\bar{x}_{1} = 10.005 \,\si{\milli\meter}$,
$s_{1} = 0.012 \,\si{\milli\meter}$; sample $2$ has $n_{2} = 30$,
$\bar{x}_{2} = 10.015 \,\si{\milli\meter}$, $s_{2} = 0.014
\,\si{\milli\meter}$. The Welch $t$-statistic is $t \approx -2.97$,
$p \approx 0.004$. The engineer rejects the equal-means null at any
of the standard thresholds. The follow-up question is whether the
$0.010 \,\si{\milli\meter}$ difference matters in the assembly's
tolerance; statistical significance and engineering relevance are
distinct.

A reliability engineer counts $42$ defects on a wafer that nominally
produces a Poisson-mean of $\lambda = 30$ defects. The
chi-square-equivalent test (or, equivalently, a one-sample test for
the Poisson mean) gives $p \approx 0.04$. The engineer notes the
disagreement with nominal but also asks whether the wafer condition
or the inspection procedure has changed since the nominal was set.

\paragraph{Worked output: the bearing-supplier two-sample test in
R and Python.} The supplier-comparison case above is the
working example for showing what the standard statistical
software produces. Running the analysis in R:

\begin{verbatim}
> sample1 <- c(...)  # 30 observations from supplier 1
> sample2 <- c(...)  # 30 observations from supplier 2
> t.test(sample1, sample2, alternative = "two.sided",
+        var.equal = FALSE, conf.level = 0.95)

        Welch Two Sample t-test

data:  sample1 and sample2
t = -2.9716, df = 56.71, p-value = 0.004296
alternative hypothesis: true difference in means is not equal to 0
95 percent confidence interval:
 -0.016745 -0.003255
sample estimates:
mean of x mean of y
   10.005    10.015
\end{verbatim}

\noindent
The $\hat{R}$ output's six pieces of information are the standard
working pieces: the test statistic ($t = -2.97$), the
Welch-Satterthwaite-adjusted degrees of freedom (df $= 56.7$,
not the equal-variance $58$), the two-sided $p$-value ($0.004$),
the alternative hypothesis statement, the $95\%$ confidence
interval on the difference of means ($-0.0167$ to $-0.0033\,\si{
\milli\meter}$, which excludes zero), and the sample means.
The reader rejects the equal-means null at $\alpha = 0.05$ and
extracts an effect-size estimate of $0.010\,\si{\milli\meter}$
with a CI of $\pm 0.007\,\si{\milli\meter}$ at $k = 2$.

In Python with SciPy:

\begin{verbatim}
>>> from scipy import stats
>>> result = stats.ttest_ind(sample1, sample2, equal_var=False)
>>> result
TtestResult(statistic=-2.9716, pvalue=0.004296, df=56.71)
>>> ci = result.confidence_interval(confidence_level=0.95)
>>> ci
ConfidenceInterval(low=-0.016745, high=-0.003255)
\end{verbatim}

\noindent
SciPy version $1.10$ and later returns the confidence interval
as a method on the result object; earlier versions require manual
computation through the test statistic and degrees of freedom. The
\texttt{equal\_var=False} argument is the Welch correction; the
default \texttt{equal\_var=True} is the equal-variance assumption
that should not be relied on without checking. The MATLAB
equivalent is \texttt{[h, p, ci, stats] = ttest2(sample1,
sample2, 'Vartype', 'unequal')}.

The same analysis in Stata: \texttt{ttest sample1 == sample2,
unpaired unequal}. In SAS: \texttt{proc ttest; class supplier;
var diameter; run;}. The cross-tool difference is in the default
variance assumption and the default-output formatting; the
numerical results agree to floating-point precision because all
five tools implement the same Welch-Satterthwaite formula.

\paragraph{Power analysis: how many observations were needed?}
Before running the test, the engineer should have determined the
sample size required to detect the smallest engineering-meaningful
difference. The standard tool is the \texttt{pwr} package in R:

\begin{verbatim}
> library(pwr)
> pwr.t.test(d = 0.77, power = 0.80, sig.level = 0.05,
+            type = "two.sample", alternative = "two.sided")

     Two-sample t test power calculation

              n = 27.4
              d = 0.77
      sig.level = 0.05
          power = 0.80
    alternative = two.sided

NOTE: n is number in *each* group
\end{verbatim}

\noindent
Cohen's $d$ of $0.77$ (the standardised effect size,
$d = |\bar{x}_{1} - \bar{x}_{2}|/s_{\text{pooled}}$) at $\alpha
= 0.05$ and power $0.80$ requires $n = 28$ per group. The
engineer's $n = 30$ per group is adequate for this effect size;
to detect a smaller effect of $d = 0.3$ would require $n = 175$
per group. The Python equivalent uses
\texttt{statsmodels.stats.power.TTestIndPower().solve\_power()};
the Stata equivalent is \texttt{power twomeans 10.005 10.015,
sd1(0.012) sd2(0.014) power(0.80)}.

The discipline is to compute the required sample size before
the data are collected, choose $n$ accordingly, and report both
the achieved power for the chosen effect size and the rationale
for the chosen effect size. The chapter's pre-registration
principle (section 8.7) applies: the sample-size choice and its
effect-size justification are part of the test specification,
not an afterthought to the result.

\subsection{Effect sizes: Cohen's \texorpdfstring{$d$}{d}, Hedges'
\texorpdfstring{$g$}{g}, Glass's \texorpdfstring{$\Delta$}{Delta},
and variance-explained measures}

A $p$-value tells the reader whether the data are unusual under the
null. The effect size tells the reader whether the difference is
large enough to matter. The two answer disjoint questions: a small
effect can produce a tiny $p$-value when $n$ is enormous, and a
large effect can produce $p > 0.05$ when $n$ is tiny. Reporting one
without the other is a recurring failure mode the rest of the
chapter develops as ``statistical significance without engineering
relevance.'' The 7th edition of the APA Publication Manual
\cite{text:ch08t3b-apa-2020} requires effect sizes with confidence
intervals alongside every reported $p$-value; the ASA's 2019
follow-up to its 2016 statement
\cite{paper:ch08t3-wasserstein-2019} extends the requirement to all
quantitative research. The engineering preference, codified in
internal quality reports and in the IEEE and ASME submission
guidelines, is the same.

\paragraph{Cohen's $d$ (two-sample standardised mean difference).}
Jacob Cohen's 1988 monograph \cite{text:ch08t3b-cohen-1988} defined

\[
d = \frac{\bar{x}_{1} - \bar{x}_{2}}{s_{\text{pooled}}},
\quad
s_{\text{pooled}} = \sqrt{\frac{(n_{1} - 1) s_{1}^{2} + (n_{2} - 1)
s_{2}^{2}}{n_{1} + n_{2} - 2}}.
\]

The denominator pools the within-group variances, weighting by
degrees of freedom. The numerator is the raw mean difference. The
ratio is unitless: a Cohen's $d$ of $0.5$ means the two group means
differ by half a within-group standard deviation, regardless of
whether the underlying quantity is bearing diameter in
millimetres or fill weight in grams.

Cohen's working benchmarks, derived from a survey of the
behavioural-sciences literature available to him in the late 1980s,
are $d = 0.2$ small, $d = 0.5$ medium, $d = 0.8$ large. The
benchmarks are field-dependent and the engineer should treat them
as orientation, not as decision thresholds: a $d = 0.2$ in
high-precision optics may be a major engineering event, while a
$d = 0.8$ in a noisy biomechanical measurement may not change
clinical practice.

\emph{Worked engineering case.} The bearing-supplier comparison in
section 8.3 had $\bar{x}_{1} = 10.005$, $\bar{x}_{2} = 10.015$
$\si{\milli\meter}$, $s_{1} = 0.012$, $s_{2} = 0.014$
$\si{\milli\meter}$, $n_{1} = n_{2} = 30$. The pooled standard
deviation is

\[
s_{\text{pooled}} = \sqrt{\frac{29 \cdot 0.012^{2} + 29 \cdot
0.014^{2}}{58}} = \sqrt{\frac{0.00418 + 0.00568}{58}}
= \sqrt{1.70 \times 10^{-4}}
\approx 0.0130 \,\si{\milli\meter}.
\]

\noindent
The standardised effect size is $d = |10.005 - 10.015|/0.0130
\approx 0.77$. By Cohen's benchmarks this is a medium-to-large
effect; the corresponding power calculation in section 8.6 used
exactly this value to argue that $n = 28$ per group would suffice
for $\alpha = 0.05$ and power $0.80$. The engineer who reports the
$p = 0.004$ alongside $d = 0.77$ is reporting both whether the
difference is statistically detectable and how large it is in
operational units.

\paragraph{Hedges' $g$ (small-sample bias correction).}
Cohen's $d$ is biased upward in small samples because $s_{\text{
pooled}}$ underestimates $\sigma_{\text{pooled}}$ when $n_{1} +
n_{2}$ is small. Larry Hedges' 1981 correction
\cite{paper:ch08t3b-hedges-1981} is

\[
g = d \cdot \left(1 - \frac{3}{4(n_{1} + n_{2}) - 9}\right).
\]

The multiplier approaches $1$ as the combined sample grows; at
$n_{1} + n_{2} = 60$ (the bearing case) the multiplier is $1 -
3/231 \approx 0.987$, so $g = 0.77 \times 0.987 \approx 0.76$, an
imperceptible adjustment. At $n_{1} + n_{2} = 10$ the multiplier is
$1 - 3/31 \approx 0.903$, a noticeable correction. The discipline:
report Cohen's $d$ for routine cases with $n_{1} + n_{2} \geq 40$,
report Hedges' $g$ for small-sample work, and state which was
computed. The R package \texttt{effsize} returns both via
\texttt{cohen.d(x, y, hedges.correction = TRUE)}; Python's
\texttt{pingouin.compute\_effsize(x, y, eftype = "hedges")} does
the same.

\paragraph{Glass's $\Delta$ (control-group-only standardiser).}
When one group is the reference (a control supplier, a baseline
process, a placebo arm) and the other is a treatment, the
within-treatment variance is sometimes inflated by the treatment
itself. Gene Glass's 1976 effect size
\cite{paper:ch08t3b-glass-1976} uses only the control-group
standard deviation as the standardiser:

\[
\Delta = \frac{\bar{x}_{\text{treatment}} - \bar{x}_{\text{control}}}{s_{\text{control}}}.
\]

The use case is when the engineer trusts the control group's
variance as a stable estimate of natural process scatter and treats
the treatment group's variance as a possibly inflated quantity that
should not contribute to the standardiser. A worked instance: a
supplier-qualification programme where supplier A is the
$25$-year incumbent (control, $s_{A} = 0.012$) and supplier B is a
new entrant whose process is still stabilising ($s_{B} = 0.020$).
Pooling the two variances would inflate $s_{\text{pooled}}$ and
deflate Cohen's $d$; Glass's $\Delta = 0.010/0.012 = 0.83$
preserves the incumbent's variance as the reference scale. Report
$\Delta$ alongside $d$ when the variances differ by more than a
factor of $1.5$ and the engineering question concerns deviation
from the control.

\paragraph{Variance-explained measures: $\eta^{2}$, $\omega^{2}$,
$f^{2}$.} For ANOVA-like settings with $k > 2$ groups, the
standardised difference of means generalises to a variance-
explained measure. The eta-squared $\eta^{2}$ is the ratio of
between-group sum of squares to total sum of squares:

\[
\eta^{2} = \frac{\text{SS}_{\text{between}}}{\text{SS}_{\text{total}}}
= 1 - \frac{\text{SS}_{\text{within}}}{\text{SS}_{\text{total}}}.
\]

The interpretation is the fraction of total variance attributable
to group membership. Like Cohen's $d$, $\eta^{2}$ is biased upward
in small samples; Hays' omega-squared
\cite{text:ch08t3b-hays-1994}

\[
\omega^{2} = \frac{\text{SS}_{\text{between}} - (k - 1)
\text{MS}_{\text{within}}}{\text{SS}_{\text{total}} +
\text{MS}_{\text{within}}}
\]

corrects the bias and is preferred for small-sample reporting.
Cohen's $f^{2}$ is the related quantity for regression models:
$f^{2} = R^{2}/(1 - R^{2})$, the ratio of explained to unexplained
variance. Cohen's benchmarks for $f^{2}$ are $0.02$ small, $0.15$
medium, $0.35$ large; for $\eta^{2}$ the benchmarks are $0.01$,
$0.06$, $0.14$. The R \texttt{effectsize} package returns
$\eta^{2}$, $\omega^{2}$, partial $\eta^{2}$, and Cohen's $f$ from
an \texttt{aov} object; Python's \texttt{pingouin.anova} returns
the same with the \texttt{effsize = "n2"} argument.

\paragraph{Confidence intervals on effect sizes.} The point
estimate of an effect size, like the point estimate of a mean, is
incomplete without an uncertainty band. Closed-form CIs are
available for $d$ (the non-central $t$ distribution), $\eta^{2}$
(the non-central $F$ distribution), and $\omega^{2}$ (asymptotic
bootstrap). For the bearing case at $d = 0.77$, $n_{1} = n_{2} = 30$,
the $95 \,\%$ CI is approximately $[0.24, 1.30]$ via the
non-central $t$ method. R: \texttt{cohen.d(x, y, conf.level =
0.95)} from the \texttt{effsize} package. Python:
\texttt{pingouin.compute\_effsize\_from\_t(t, nx, ny,
eftype = "cohen")} with bootstrap CI from
\texttt{pingouin.compute\_bootci}. The engineer who reports $d =
0.77$, $95 \,\%$ CI $[0.24, 1.30]$, $p = 0.004$ has reported all
three pieces a colleague needs to judge both the size and the
reliability of the result.

\paragraph{Reporting template.} The post-ASA-2016 engineering
report carries, for each comparison: the point estimate of the
effect size, its $95 \,\%$ confidence interval, the test statistic,
the $p$-value, and the sample size. The four-number minimum
($d$, $\text{CI}_{d}$, $p$, $n$) lets a downstream reader recompute
the test, judge practical relevance, and audit the sample-size
adequacy. The recurring failure mode is the report that lists only
$p$: a colleague who reads it can neither size the effect nor
judge whether the test was adequately powered for the effect that
mattered.

\section{Confidence vs prediction intervals}\pathtag{core}

Chapter 4 §4.6 introduced the confidence interval as an engineering
tool: a range of values, computed from a sample, that is intended to
contain the true value of a measured quantity with a stated coverage
probability. The frequentist definition is procedural: the interval-
computing method, in repeated application, contains the true
parameter the stated fraction of the time
\cite[ch.~6]{text:wasserman2004}.

This chapter draws a distinction Chapter 4 did not need: the
\engterm{confidence interval} is for a parameter (the mean, the
slope of a regression, the proportion of defects). The
\engterm{prediction interval} is for the next observation. The
prediction interval is wider, sometimes much wider. A reader who
uses a confidence interval for a next-unit decision will usually
understate operational scatter.

\subsection{The confidence interval for a mean}

For a sample of $n$ observations from a distribution with unknown
mean $\mu$ and unknown standard deviation, the $100(1 - \alpha) \,\%$
confidence interval for $\mu$ is

\[
\bar{x} \pm t_{1-\alpha/2, n-1} \cdot \frac{s}{\sqrt{n}},
\]

where $t_{1-\alpha/2, n-1}$ is the positive critical value of
Student's $t$-distribution with $n - 1$ degrees of freedom. The
half-width $t_{1-\alpha/2, n-1} \cdot s/\sqrt{n}$ is the standard uncertainty of
the mean times a coverage factor. For $n \geq 30$ and $\alpha =
0.05$, the critical value is close to $1.96$; for $n = 10$ it is
about $2.26$; for $n = 5$ it is about $2.78$. The reader sees the
small-sample correction directly in the wider interval.

The width of the confidence interval shrinks as $1/\sqrt{n}$. To
halve the width takes four times the data; to reduce by a factor of
ten takes a hundred times. This is the same $1/\sqrt{n}$ wall
Chapter 4 §4.7 named.

\subsection{The prediction interval for the next observation}

The prediction interval answers a different question. Given the
observed sample, what range will the next single observation fall
into with stated probability? A new draw has the population's
inherent variability $\sigma$ added to the uncertainty of the
estimated mean. For a single new observation from a Gaussian
population, the $100(1 - \alpha) \,\%$ prediction interval is

\[
\bar{x} \pm t_{1-\alpha/2, n-1} \cdot s \cdot \sqrt{1 + \frac{1}{n}}.
\]

The factor $\sqrt{1 + 1/n}$ replaces the $1/\sqrt{n}$ of the
confidence interval. As $n \to \infty$, the factor approaches one,
so the prediction interval tends to $\bar{x} \pm t_{\alpha/2} \cdot
s$, which is the inherent population scatter. The confidence
interval, in the same limit, shrinks to zero width: the mean is
known exactly. The prediction interval cannot shrink below the
population's variability no matter how much data the engineer
collects.

\subsection{Why the distinction is operational}

A laboratory quality control engineer reports a $95 \,\%$ confidence
interval for the mean of a key process variable. A production
operator reads the report and assumes that $95 \,\%$ of new units
will fall within that interval. The operator is wrong, often by a
factor of three or more. The interval the operator wanted is the
prediction interval, which is wider.

A medical device engineer reports a $95 \,\%$ confidence interval
for the mean response time of a sensor. A clinical user reads the
report and assumes $95 \,\%$ of new readings will fall within that
interval. Same error.

The discipline is to ask, before constructing the interval, what
question the interval is meant to answer. If the question is about
the parameter (the mean of the process, the calibration coefficient,
the average rate), the answer is a confidence interval. If the
question is about the next observation (what will the next reading
be, what range should we tolerate, what specification should we
write), the answer is a prediction interval.

A third interval, the \engterm{tolerance interval}, contains a
specified proportion of the population with a specified confidence.
Tolerance intervals are the right tool for setting acceptance bounds
that should contain, say, $99 \,\%$ of all production units with
$95 \,\%$ confidence. The Devore text \cite[ch.~7]{text:devore2015}
gives the engineering tables.

\subsection{Tolerance intervals: the working procedure}

The frequentist tolerance interval for a normal population
$\mathcal{N}(\mu, \sigma^{2})$ takes the form $\bar{x} \pm k \cdot
s$ where $k$ is the \engterm{tolerance factor} that depends on
the sample size $n$, the population coverage $p$ (the fraction
of the population the interval is meant to contain), and the
confidence level $1 - \gamma$ (the long-run probability that
the constructed interval actually contains that fraction).
The exact two-sided tolerance factor is computed from the
noncentral chi-square distribution; the Howe approximation
\cite{paper:ch08t3p-howe-1969} gives the working closed form

\[
k \approx z_{p^{*}} \sqrt{\frac{(n - 1)(1 + 1/n)}{\chi^{2}_{\gamma, n-1}}},
\]

where $z_{p^{*}}$ is the normal quantile at $p^{*} = (1 + p)/2$
for a two-sided interval, and $\chi^{2}_{\gamma, n-1}$ is the
lower $\gamma$-quantile of the chi-square distribution with
$n - 1$ degrees of freedom. The NIST/SEMATECH e-Handbook
\cite{web:ch08t3p-nist-eshandbook-7263} tabulates Howe-formula
factors out to $n = 1000$; for the working ranges $n = 10$ to
$n = 100$ at $p = 0.95$ or $p = 0.99$ at $\gamma = 0.05$, the
formula is accurate to within $0.1\%$ of the exact value.

\paragraph{Worked example: acceptance bounds on shaft diameter.}
A precision-shaft supplier measures a production run of
$n = 30$ shafts and reports $\bar{x} = 15.012\,\si{\milli\meter}$
with $s = 0.008\,\si{\milli\meter}$. The customer requires that
the supplier's shipped lots contain $99\%$ of units within an
acceptance window with $95\%$ confidence: $p = 0.99$, $\gamma
= 0.05$, $n = 30$. The Howe factor: $z_{0.995} = 2.576$,
$\chi^{2}_{0.05, 29} = 17.708$, giving
$k = 2.576\sqrt{29 \times (1 + 1/30)/17.708} = 2.576 \times
\sqrt{29 \times 1.033/17.708} = 2.576 \times 1.301 = 3.351$.
The acceptance window is $\bar{x} \pm 3.351 s = 15.012 \pm
0.027 = (14.985, 15.039)\,\si{\milli\meter}$. Compare to the
$95\%$ confidence interval for the mean ($\bar{x} \pm 0.003$),
the $95\%$ prediction interval for the next unit ($\bar{x} \pm
0.016$), and the unadjusted $\pm 3s$ heuristic ($\bar{x} \pm
0.024$). The tolerance interval is wider than the prediction
interval by the factor that accounts for the bound on the
population proportion under sampling uncertainty.

\paragraph{The same calculation in software.} In R, the
\texttt{tolerance} package \cite{web:ch08t3p-r-tolerance} returns
the exact factor: \texttt{normtol.int(x, alpha = 0.05, P = 0.99,
side = 2)} returns a tibble with \texttt{2-sided.lower = 14.985,
2-sided.upper = 15.039}. In Python with the
\texttt{statsmodels} package: \texttt{statsmodels.stats.tolerance.
tolerance\_interval(x, alpha=0.05, P=0.99, sides=2)}. In SAS:
\texttt{proc capability; var diameter; intervals normal tol;}.
In Minitab: \texttt{Stat $\to$ Quality Tools $\to$ Tolerance
Intervals}. The exact noncentral chi-square calculation differs
from the Howe approximation at $n < 10$ by up to $3\%$; the
working choice for $n \geq 30$ is the Howe formula; for $n < 10$,
software with the exact procedure.

\paragraph{One-sided tolerance intervals.} For an upper-bound
specification (e.g., maximum allowable surface roughness, maximum
contaminant concentration), the one-sided tolerance bound is
$\bar{x} + k_{1} s$ where the one-sided factor $k_{1}$ uses the
$p$-quantile (not $p^{*}$) and a different denominator. The
working closed form is the noncentral-$t$ approximation
\[
k_{1} \approx \frac{z_{p} + \sqrt{z_{p}^{2} - a b}}{a},
\quad
a = 1 - \frac{z_{\gamma}^{2}}{2(n - 1)},
\quad
b = z_{p}^{2} - \frac{z_{\gamma}^{2}}{n},
\]
where $z_{p}$ is the normal quantile at the coverage $p$ and
$z_{\gamma}$ the quantile at the confidence $1 - \gamma$. ISO
16269-6:2014 \cite{std:ch08t3p-iso-16269-6} gives the
one-sided factors with engineering examples.

\textit{Worked example: an upper bound on surface roughness.} A
machining process is specified so that no more than $5 \,\%$ of
parts exceed a roughness ceiling, and the customer wants that
statement backed at $95 \,\%$ confidence: $p = 0.95$,
$\gamma = 0.05$. A sample of $n = 40$ parts gives
$\bar{x} = 0.82\,\si{\micro\meter}$ (arithmetic-average roughness
$R_{a}$) with $s = 0.11\,\si{\micro\meter}$. The one-sided factor
uses $z_{0.95} = 1.645$ for both the coverage and the confidence
(here $p = 1 - \gamma$ coincidentally), giving
$a = 1 - 1.645^{2}/(2 \cdot 39) \approx 0.9653$,
$b = 1.645^{2} - 1.645^{2}/40 \approx 2.638$, and
\[
k_{1} \approx \frac{1.645 + \sqrt{1.645^{2} - 0.9653 \cdot 2.638}}
{0.9653} \approx \frac{1.645 + 0.399}{0.9653} \approx 2.12.
\]
The upper tolerance bound is
$\bar{x} + k_{1} s = 0.82 + 2.12 \cdot 0.11 \approx
1.05\,\si{\micro\meter}$. The engineering statement: with
$95 \,\%$ confidence, at least $95 \,\%$ of parts have $R_{a}$
below $1.05\,\si{\micro\meter}$. Compare the naive
$\bar{x} + 1.645 s = 1.00\,\si{\micro\meter}$, which ignores the
sampling uncertainty in $\bar{x}$ and $s$ and understates the
bound by $0.05\,\si{\micro\meter}$; the gap between $k_{1} = 2.12$
and the bare normal quantile $1.645$ is the price of not knowing
$\mu$ and $\sigma$ exactly at $n = 40$. As $n$ grows the factor
$k_{1}$ falls toward $z_{p} = 1.645$; at $n = 40$ the sampling
penalty is about $29 \,\%$ of the factor.

\paragraph{Nonparametric tolerance intervals.} When the
Gaussian assumption fails, the distribution-free tolerance
interval uses order statistics. For a sample of size $n$ the
interval $[x_{(r)}, x_{(n-s+1)}]$ contains at least proportion
$p$ of the population with confidence $1 - \gamma$ where
$r + s$ satisfies a binomial sum involving $n$, $p$, and
$\gamma$. The advantage: no distributional assumption. The
disadvantage: very large $n$ is required for tight intervals
($n \geq 459$ for $p = 0.99$, $\gamma = 0.05$, $r = s = 1$, two
sided). For working data sets where distributional shape is
suspect, fit a tolerance interval on the log or other
transformed scale; for genuinely heavy-tailed populations, the
nonparametric procedure is the safe default.

\subsection{The bootstrap, when the parametric form is suspect}

The intervals above assume a parametric model for the sample (most
often, Gaussian). The \engterm{bootstrap}
\cite{paper:ch08-efron-tibshirani-1993} replaces the parametric
assumption with an empirical resampling procedure. Given a sample
of size $n$, the analyst draws $B$ resamples of size $n$ with
replacement, computes the statistic of interest on each resample,
and uses the empirical distribution of the $B$ replicates as an
estimate of the sampling distribution of the statistic.
\Cref{fig:vol01:ch08:bootstrap-schematic} shows the pipeline.

\begin{figure}[ht]
\centering
\begin{tikzpicture}[x=1cm, y=1cm, font=\small,
  every node/.style={font=\scriptsize}]

% Original sample box
\draw[thick, fill=engblue!10] (0, 4) rectangle (3.2, 5.0)
  node[pos=0.5] {original sample $x_{1}, \ldots, x_{n}$};

% Down arrow into the resampling cloud
\draw[->, thick] (1.6, 4.0) -- (1.6, 3.4);
\node[anchor=west] at (1.65, 3.7) {resample with replacement};

% Three resample boxes (representative of B)
\foreach \i/\xs in {1/-2.0, 2/2.0, 3/6.0} {
  \draw[thick, fill=engamber!10] (\xs, 2.0) rectangle (\xs+3.2, 3.0)
    node[pos=0.5] {$x^{*(\i)}_{1}, \ldots, x^{*(\i)}_{n}$};
  \draw[->, thick] (1.6, 3.4) -- (\xs+1.6, 3.0);
}

% Dots between to indicate B replicates
\node at (4.6, 2.5) {$\cdots$};
\node at (8.8, 2.5) {$\cdots$};

% Statistic computation
\foreach \xs in {-2.0, 2.0, 6.0} {
  \draw[->, thick] (\xs+1.6, 2.0) -- (\xs+1.6, 1.4);
}
\foreach \i/\xs in {1/-2.0, 2/2.0, 3/6.0} {
  \draw[thick, fill=engred!10] (\xs+0.4, 0.4) rectangle (\xs+2.8, 1.4)
    node[pos=0.5] {$\hat{\theta}^{*(\i)}$};
}

% Aggregation arrow
\foreach \xs in {-2.0, 2.0, 6.0} {
  \draw[->, thick] (\xs+1.6, 0.4) -- (4.6, -0.4);
}

% Final histogram-of-replicates box
\draw[thick, fill=engblue!15] (1.8, -1.6) rectangle (7.4, -0.4)
  node[pos=0.5, align=center, font=\scriptsize]
  {empirical distribution of $\hat{\theta}^{*}$ \\
   (bootstrap CI from percentiles)};

\end{tikzpicture}
\caption[Bootstrap resampling schematic]{The nonparametric bootstrap
in one picture. From the original sample of size $n$, the analyst
draws $B$ resamples of size $n$ with replacement. Each resample
produces a replicate of the statistic $\hat{\theta}$. The empirical
distribution of the $B$ replicates approximates the sampling
distribution of $\hat{\theta}$; the $2.5$th and $97.5$th
percentiles give a $95 \,\%$ percentile bootstrap confidence
interval. Typical $B$ is $10^{3}$ to $10^{4}$ for interval
estimation, $10^{4}$ or more when the tail probability itself is of
interest.}
\label{fig:vol01:ch08:bootstrap-schematic}
\end{figure}


A $95 \,\%$ percentile-bootstrap confidence interval for the
statistic is the interval between the $2.5$th and $97.5$th
percentiles of the bootstrap distribution. The bootstrap inherits
the sample's biases (a non-representative sample yields a
non-representative bootstrap) and breaks down for statistics whose
sampling distribution depends sensitively on the tails (sample
maximum, extreme quantiles); it is the default tool for confidence
intervals on statistics that have no closed-form sampling
distribution: ratios, medians, trimmed means, regression
coefficients in non-linear models, correlations between dependent
samples. Typical $B$ is $10^{3}$ for an interval and $10^{4}$ or
more when the tail probability itself is the quantity of interest.

\paragraph{Worked example.} Applied to the commute series, the
bootstrap returns a $95 \,\%$ interval of approximately
$[28.0, 29.5] \,\text{min}$, in close agreement with the $t$-based
interval above. The agreement is expected: $n = 30$ is large enough
that the sampling distribution of the mean is approximately
Gaussian by the central limit theorem, and the bootstrap and
$t$-based intervals coincide when both apply. The bootstrap earns
its keep when, for instance, the analyst wants a confidence
interval for the sample median or for the $90$th-percentile
commute, both of which have no clean parametric form.

\begin{mastery}
\textbf{The delta method: closed-form uncertainty for a function
of an estimate.} The bootstrap resamples to find the sampling
distribution of a statistic. When the statistic is a smooth
function $g(\hat\theta)$ of an estimate $\hat\theta$ whose own
standard error is known, the \engterm{delta method} gives the
answer in closed form and connects this chapter to the
propagation-of-uncertainty rule of Chapter 4 §4.5. A first-order
Taylor expansion of $g$ about the true $\theta$ gives
\[
g(\hat\theta) \approx g(\theta) + g'(\theta)(\hat\theta - \theta),
\]
so, taking variances and using $\Var(\hat\theta)$,
\[
\Var\!\left(g(\hat\theta)\right) \approx
\left[g'(\theta)\right]^{2} \Var(\hat\theta),
\qquad
\text{se}\!\left(g(\hat\theta)\right) \approx
\left|g'(\hat\theta)\right| \cdot \text{se}(\hat\theta).
\]
This is Chapter 4's single-variable uncertainty-propagation
formula, now read as a statement about the sampling distribution
of an estimator rather than about instrument error.

\smallskip
\textit{Worked instance.} A reliability engineer estimates a
failure rate $\hat\lambda = 0.020$ per hour with standard error
$\text{se}(\hat\lambda) = 0.004$ per hour and wants the
uncertainty of the mean time between failures
$\text{MTBF} = g(\lambda) = 1/\lambda$. Here
$g'(\lambda) = -1/\lambda^{2}$, so
\[
\text{se}(\widehat{\text{MTBF}}) \approx
\frac{1}{\hat\lambda^{2}} \cdot \text{se}(\hat\lambda)
= \frac{0.004}{0.020^{2}} = 10\,\si{\hour}.
\]
The point estimate is $1/0.020 = 50\,\si{\hour}$, so the MTBF is
$50 \pm 10\,\si{\hour}$ to first order. The delta method is exact
only in the linear limit; for the strongly nonlinear reciprocal at
a small $\hat\lambda$, the sampling distribution of the MTBF is
right-skewed, and a bootstrap or a transformation to the log scale
gives a more honest interval. The working rule: the delta method
is the quick closed-form check, and the bootstrap is the arbiter
when the transform is sharply curved over the range of the
estimate's uncertainty.
\end{mastery}

\subsection{Bias-corrected percentile bootstrap: a sensor-calibration
study}

The percentile bootstrap is unbiased on the median and the mean of
a roughly symmetric sample. On skewed statistics, on small samples,
or on the upper percentiles of a heavy-tailed distribution, the
percentile interval misses the true parameter in a direction the
analyst can detect and correct. The bias-corrected and accelerated
(BCa) bootstrap, originally developed by Efron
\cite{paper:ch08-efron-bca-1987} and standardised in Davison and
Hinkley \cite{text:ch08-davison-hinkley-1997}, applies two
corrections: a \emph{bias correction} that shifts the percentile
endpoints to match the median of the bootstrap distribution to the
sample estimate, and an \emph{acceleration} that accounts for the
rate at which the standard error of the statistic changes with the
true parameter. The result is a confidence interval whose coverage
matches the nominal level to second order for a wide class of
problems, where the bare percentile interval matches only to first
order.

\textit{The engineering case.} A pressure transducer is calibrated
against a primary deadweight tester across $n = 24$ load points
covering the full operating range. The residuals, defined as the
transducer reading minus the deadweight reference, are recorded
(current as of 2025):
\[
\begin{aligned}
&-0.12, +0.08, +0.31, -0.05, +0.19, -0.22, +0.14, +0.41, \\
&+0.27, -0.03, +0.08, +0.16, +0.52, +0.11, -0.08, +0.24, \\
&+0.18, +0.06, +0.33, +0.21, +0.09, +0.04, +0.29, +0.61
\end{aligned}
\]
all in units of $\si{\kilo\pascal}$. The sample mean is
$\bar{r} \approx 0.166 \,\si{\kilo\pascal}$, median is
$\tilde{r} \approx 0.155 \,\si{\kilo\pascal}$, and sample standard
deviation is $s \approx 0.182 \,\si{\kilo\pascal}$. The
distribution is visibly right-skewed: a long upper tail and
fewer-than-Gaussian observations on the lower tail. A $t$-based
$95 \,\%$ CI for the mean residual is
$0.166 \pm 2.07 \cdot 0.182/\sqrt{24} \approx 0.166 \pm 0.077
\,\si{\kilo\pascal}$, or $[0.089, 0.243] \,\si{\kilo\pascal}$. The
calibration question is whether the lower bound clears zero, in
which case the transducer is biased high and an offset correction
is justified.

\textit{Setup of the percentile bootstrap.} Draw $B = 10{,}000$
resamples of size $n = 24$ with replacement from the residual
series. For each, compute the resample mean
$\bar{r}^{*}_{b}$. The percentile $95 \,\%$ CI is the
$[2.5\text{th}, 97.5\text{th}]$ quantile of the $10{,}000$ bootstrap
means. For this sample the percentile CI is approximately
$[0.095, 0.244] \,\si{\kilo\pascal}$, agreeing with the $t$-based
CI to within Monte Carlo noise.

\textit{Bias correction.} Let $G^{*}(\cdot)$ be the empirical CDF
of the bootstrap means. Define
\[
z_{0} = \Phi^{-1}\!\left(G^{*}(\bar{r})\right),
\]
the standard-normal quantile of the bootstrap probability that
$\bar{r}^{*} \leq \bar{r}$. A symmetric bootstrap distribution
gives $G^{*}(\bar{r}) = 0.5$, so $z_{0} = 0$. A right-skewed
bootstrap distribution has $\bar{r}^{*}$ above $\bar{r}$ more often
than below, so $G^{*}(\bar{r}) < 0.5$ and $z_{0} < 0$. For our
sample, the bootstrap distribution has $G^{*}(\bar{r}) \approx 0.47$,
giving $z_{0} \approx -0.075$.

\textit{Acceleration constant.} Compute the leave-one-out
jackknife means $\bar{r}_{(-i)}$ for $i = 1, \ldots, n$ and define
the acceleration
\[
a = \frac{\sum_{i=1}^{n} (\bar{r}_{(\cdot)} - \bar{r}_{(-i)})^{3}}
{6 \left[\sum_{i=1}^{n} (\bar{r}_{(\cdot)} - \bar{r}_{(-i)})^{2}
\right]^{3/2}},
\]
where $\bar{r}_{(\cdot)}$ is the mean of the jackknife replicates.
For the calibration sample, $a \approx 0.062$. A positive $a$
indicates that the standard error of the mean grows as the mean
itself grows, which is the signature of a right-skewed parent
distribution.

\textit{BCa endpoints.} The corrected lower and upper percentile
levels are
\[
\alpha_{1} = \Phi\!\left(z_{0} + \frac{z_{0} + z_{0.025}}
{1 - a(z_{0} + z_{0.025})}\right),
\quad
\alpha_{2} = \Phi\!\left(z_{0} + \frac{z_{0} + z_{0.975}}
{1 - a(z_{0} + z_{0.975})}\right).
\]
With $z_{0} = -0.075$, $a = 0.062$, $z_{0.025} = -1.96$,
$z_{0.975} = 1.96$, the corrected levels evaluate to
$\alpha_{1} \approx 0.0145$ and $\alpha_{2} \approx 0.962$. The BCa
$95 \,\%$ interval is the $[1.45 \text{th}, 96.2 \text{nd}]$
quantile of the bootstrap distribution: approximately
$[0.087, 0.234] \,\si{\kilo\pascal}$.

\textit{Comparison.} Three intervals agree on the central
estimate near $0.16 \,\si{\kilo\pascal}$ and on the engineering
verdict (the lower bound clears zero, the transducer is biased
high). They differ in the upper bound by roughly
$0.01 \,\si{\kilo\pascal}$, which is below the calibration
threshold for this transducer. The $t$-based interval is slightly
too wide on the upper side (treats the parent as symmetric); the
percentile interval is slightly too narrow on the lower side
(carries the bootstrap skew through unchanged); the BCa interval
corrects both directions and is the working default when the
sample shows visible skew.

\textit{When BCa pays.} On samples with $n \geq 30$ and near-
Gaussian shape, the three intervals coincide and the choice is
immaterial. On samples with $n \leq 20$, visible skew, or for
statistics other than the mean (median, ratios, regression
coefficients, $90$th percentile), the percentile and $t$-based
intervals can mis-state coverage by enough to affect a calibration
decision. The BCa interval requires the jackknife pass (linear in
$n$) plus the percentile bootstrap (linear in $B$), totalling
$O(B + n B)$ statistic evaluations at modest additional cost. The
reference implementation is the \texttt{boot} package in R
\cite{text:ch08-davison-hinkley-1997}; equivalent routines are
available in SciPy and in MATLAB's Statistics Toolbox. The
companion script
\repopath{volumes/01-quantity/ch08-statistics-for-engineers/code/bootstrap\_ci.py}
includes BCa as an option flag and runs the calibration sample
above.

\textit{Reporting.} \emph{Pressure-transducer calibration residual:
$\bar{r} = 0.166 \,\si{\kilo\pascal}$, BCa $95 \,\%$ CI
$[0.087, 0.234] \,\si{\kilo\pascal}$, $n = 24$, $B = 10{,}000$,
$z_{0} = -0.075$, $a = 0.062$.} The reader can audit every
correction by replaying the bootstrap and the jackknife on the raw
residuals.

\subsection{Worked example: the commute revisited}

Return to the $30$-day commute sample of section 8.1, with $\bar{x}
= 28.73 \,\text{min}$, $s = 1.98 \,\text{min}$, $n = 30$. The
$95 \,\%$ confidence interval for the mean is

\[
28.73 \pm 2.05 \cdot \frac{1.98}{\sqrt{30}}
= 28.73 \pm 0.74 \,\text{min},
\]

so approximately $[28.0, 29.5] \,\text{min}$. This interval comes
from a procedure with $95 \,\%$ long-run coverage for the
population mean, assuming the model conditions hold.

The $95 \,\%$ prediction interval for the next morning's commute is

\[
28.73 \pm 2.05 \cdot 1.98 \cdot \sqrt{1 + 1/30}
= 28.73 \pm 4.11 \,\text{min},
\]

so approximately $[24.6, 32.8] \,\text{min}$. The prediction
statement concerns the next commute, not the long-run average. The
prediction interval is more than five times wider than the
confidence interval.

The operational consequence is direct. The confidence interval
supports a claim about the long-run average commute. The prediction
interval supports a claim about tomorrow's commute. A schedule
promise for one future trip must use the prediction interval or a
percentile model, not the confidence interval.

\subsection{Process capability: \texorpdfstring{$C_{p}$, $C_{pk}$,
$P_{p}$, $P_{pk}$}{Cp, Cpk, Pp, Ppk} from a real sample}

A manufacturing process produces parts to a specification with a
lower specification limit (LSL) and an upper specification limit
(USL). The engineer asks: does the current process produce
parts that meet specification, and if so, with what margin? The
\engterm{process capability} indices answer that question by
relating the specification width to the process spread. ISO 22514-2
\cite{std:ch08-iso-22514-2-2017} codifies the indices for routine
reporting; Montgomery \cite[ch.~8]{text:ch08-montgomery-spc-2019}
gives the engineering interpretation.

\textit{The four indices.} Two pairs of indices serve two distinct
purposes. The $C$ indices use the within-subgroup (short-term)
standard deviation $\hat{\sigma}_{\text{within}}$ and assume the
process is in statistical control. The $P$ indices use the overall
(long-term) standard deviation $\hat{\sigma}_{\text{overall}} = s$
computed from the full sample without regard to subgrouping, and
make no assumption about control. The definitions are
\[
C_{p} = \frac{\text{USL} - \text{LSL}}{6 \hat{\sigma}_{\text{within}}},
\quad
C_{pk} = \min\!\left(
\frac{\text{USL} - \mu}{3 \hat{\sigma}_{\text{within}}},
\frac{\mu - \text{LSL}}{3 \hat{\sigma}_{\text{within}}}
\right),
\]
\[
P_{p} = \frac{\text{USL} - \text{LSL}}{6 s},
\quad
P_{pk} = \min\!\left(
\frac{\text{USL} - \bar{x}}{3 s},
\frac{\bar{x} - \text{LSL}}{3 s}
\right).
\]
The numerator $6\sigma$ is the natural process spread under a
Gaussian assumption: $\pm 3\sigma$ captures about $99.73 \,\%$ of
production. The $C_{p}$ and $P_{p}$ indices compare the
specification width to that natural spread; values above one mean
the spec is wider than the spread, values below one mean the
process cannot meet spec without rework or scrap.

\textit{The engineering case.} A precision machining process
produces shafts to a target diameter of $25.000 \,\si{\milli\meter}$
with LSL $= 24.950 \,\si{\milli\meter}$ and
USL $= 25.050 \,\si{\milli\meter}$ (a $\pm 0.050 \,\si{\milli\meter}$
tolerance band, current as of 2025). A sample of $n = 125$ shafts is
drawn across $25$ subgroups of $5$ consecutive parts each, providing
both the overall and within-subgroup variability. The sample
yields $\bar{x} = 25.012 \,\si{\milli\meter}$, $s = 0.018
\,\si{\milli\meter}$, and within-subgroup standard deviation
$\hat{\sigma}_{\text{within}} = 0.014 \,\si{\milli\meter}$
(estimated from the average range, $\bar{R}/d_{2}$ with
$d_{2} = 2.326$ for subgroup size $5$).

\textit{Computation of $C_{p}$.}
\[
C_{p} = \frac{25.050 - 24.950}{6 \cdot 0.014}
= \frac{0.100}{0.084}
\approx 1.19.
\]
The process spread is about $84 \,\%$ of the specification width,
which means the process is capable in principle (the natural spread
fits inside the spec) but with limited margin.

\textit{Computation of $C_{pk}$.}
\[
\frac{\text{USL} - \mu}{3 \hat{\sigma}_{\text{within}}}
= \frac{25.050 - 25.012}{3 \cdot 0.014}
= \frac{0.038}{0.042} \approx 0.90,
\]
\[
\frac{\mu - \text{LSL}}{3 \hat{\sigma}_{\text{within}}}
= \frac{25.012 - 24.950}{3 \cdot 0.014}
= \frac{0.062}{0.042} \approx 1.48.
\]
The minimum is on the USL side: $C_{pk} \approx 0.90$. The process
mean has drifted upward of the specification target by $0.012
\,\si{\milli\meter}$, so the working margin on the USL side is
about three sigma minus the offset.

\textit{Computation of $P_{p}$ and $P_{pk}$.} Using the overall
$s = 0.018 \,\si{\milli\meter}$:
\[
P_{p} = \frac{0.100}{6 \cdot 0.018}
= \frac{0.100}{0.108}
\approx 0.93,
\]
\[
P_{pk} = \min\!\left(\frac{0.038}{0.054},
\frac{0.062}{0.054}\right)
= \min(0.70, 1.15)
\approx 0.70.
\]
The long-term indices are systematically lower than the short-term
indices. The gap $P_{p} = 0.93 < C_{p} = 1.19$ means subgroup-to-
subgroup variation is contributing extra spread that within-
subgroup analysis misses; the process is not in tight control
across the $25$ subgroups.

\textit{Interpreting the four numbers.} An engineering reading of
the four indices proceeds in three steps.

\textit{Step 1: capability ceiling.} $C_{p} > 1$ confirms that the
process \emph{could} fit the spec if the mean were centred and the
between-subgroup variation were eliminated. $P_{p} < 1$ confirms
that the process \emph{does not currently} fit the spec under
its full operational variability. The gap between $C_{p}$ and
$P_{p}$ is the headroom available if the engineer can stabilise
the between-subgroup behaviour.

\textit{Step 2: centring.} The gap $C_{pk} = 0.90 < C_{p} = 1.19$
means the mean is off-target. The fix is a mean shift, often a
tool-offset adjustment. After centring (assuming
$\hat{\sigma}_{\text{within}}$ unchanged), $C_{pk}$ would rise to
match $C_{p}$ at $1.19$.

\textit{Step 3: long-term defect rate.} Assuming a Gaussian parent,
the fraction of parts above the USL is
\[
\Pr(X > \text{USL}) = \Pr\!\left(Z > \frac{\text{USL} - \bar{x}}
{s}\right)
= \Pr(Z > 2.11)
\approx 0.0174,
\]
or about $17{,}400$ parts per million (ppm). The fraction below
the LSL is $\Pr(Z < -3.44) \approx 2.9 \times 10^{-4}$, or about
$290$ ppm. Total expected nonconforming: about $17{,}700$ ppm.
By contrast, a centred process with $C_{p} = 1.33$ (the working
``capable'' threshold) produces about $63$ ppm; a centred process
with $C_{p} = 1.67$ produces about $0.6$ ppm.

\textit{Industry thresholds.} Working benchmarks (informal but
widely used in automotive and aerospace QC) are:
$C_{pk} \geq 1.00$ minimal acceptability,
$C_{pk} \geq 1.33$ capable for routine production,
$C_{pk} \geq 1.67$ capable with margin (six-sigma target).
The shaft process at $C_{pk} \approx 0.90$ does not meet even the
minimal acceptability threshold and would not pass a customer
audit at the $1.33$ level.

\textit{Engineering response.} Three interventions are sequenced
by cost. (i) Centre the mean (tool offset; estimated benefit
$C_{pk}$ from $0.90$ to $1.19$, defect rate from $17{,}400$ to
roughly $200$ ppm). (ii) Reduce between-subgroup variation (find
the assignable cause of the gap $P_{p} < C_{p}$; estimated benefit
$P_{pk}$ closes to $C_{pk}$). (iii) Reduce within-subgroup
variation (longer-cycle fix: machine maintenance, tooling
upgrades, fixture redesign; pushes both indices upward together).
The capability analysis is the diagnostic; the corrective sequence
is the design.

\textit{Common errors.} (i) Reporting $C_{p}$ alone without
$C_{pk}$, hiding a centring problem. (ii) Reporting only the
within-subgroup indices, hiding between-subgroup drift.
(iii) Computing capability on a sample that is not in statistical
control (large $\bar{R}$ subgroups mixed with small), producing
indices that are not interpretable. (iv) Forgetting the Gaussian
assumption: for skewed processes (machining time, surface finish,
positive-bounded measurements), the percentile-based capability
$C_{pm}$ or a transformation is required.

\textit{Reporting.} \emph{Shaft diameter, $n = 125$, $25$ subgroups
of $5$: $\bar{x} = 25.012$, $s = 0.018$, $\hat{\sigma}_{\text{within}}
= 0.014$ (all $\si{\milli\meter}$); $C_{p} = 1.19$, $C_{pk} = 0.90$,
$P_{p} = 0.93$, $P_{pk} = 0.70$; mean off-target by $+0.012
\,\si{\milli\meter}$; estimated total nonconforming $\approx 17{,}700$
ppm; not capable to $C_{pk} \geq 1.33$.} The customer can audit
every index.

\textit{Standards references.} Three published references frame
the working calculation. The AIAG Statistical Process Control
Reference Manual, 2nd edition (2005)
\cite{std:ch08t3b-aiag-spc-2005}, jointly published by Chrysler,
Ford, and General Motors via the Automotive Industry Action Group,
codifies the $C_{p}$, $C_{pk}$, $P_{p}$, $P_{pk}$ definitions used
across the automotive supply chain; chapter 4 gives the
calculation procedure and chapter 5 the working interpretation.
ISO 22514-2:2017
\cite{std:ch08-iso-22514-2-2017}, section 6.2, gives the
international equivalent and harmonises the indices' definitions
with the AIAG procedure. ISO 22514-7:2012 extends the framework
to measurement-system capability. The auto-industry working
threshold $C_{pk} \geq 1.33$ for a routine production process and
$C_{pk} \geq 1.67$ for safety-critical or supplier-qualification
work is recorded in AIAG SPC chapter 5, section 5.3; it is
working practice rather than a statistical theorem. The threshold
corresponds to a long-term defect rate of about $63$ ppm under a
Gaussian assumption with the standard $1.5\sigma$ long-term mean
shift assumed by the six-sigma framework.

\textit{Process performance vs process capability: the
operational distinction.} The $C$ indices use the within-subgroup
standard deviation, which under the Shewhart subgrouping
assumption captures only the common-cause (short-term) variation.
The $P$ indices use the overall standard deviation, which
captures both common-cause and special-cause (long-term)
variation. The gap $P_{p} < C_{p}$ is therefore the engineering
evidence that the process is being disturbed by assignable causes
not present within a single subgroup: tool wear over a shift,
operator handoff between shifts, raw-material lot variation,
seasonal temperature swings. The corrective action targets the
assignable cause; the $P$ index does not improve until the
assignable cause is removed, even when the $C$ index is
unchanged. AIAG SPC chapter 4, section 4.5 lists the diagnostic
sequence for identifying the assignable cause: stratify the data
by suspect factors (shift, operator, lot, tool), recompute the
subgroup variance for each stratum, and identify the stratum
whose within-stratum variance approaches the within-subgroup
variance of the full sample.

\textit{Skewed processes: $C_{pm}$ and the percentile-based
capability.} A Gaussian capability index uses $6 \hat\sigma$ as
the process spread. For a skewed process the long tail extends
beyond $\bar{x} + 3 \hat\sigma$, and the capability index
overstates the available margin on that side. The remedy is the
percentile-based capability $C_{pk, \text{percentile}}$, which
replaces $\bar{x} - 3 \hat\sigma$ and $\bar{x} + 3 \hat\sigma$
with the $0.135$th and $99.865$th sample percentiles of the
observed distribution. The two coincide for a Gaussian sample by
construction; for a skewed sample the percentile-based version is
the conservative report. The alternative is to transform the
data (Box-Cox or log) until the transformed sample is Gaussian,
report the capability on the transformed scale, and back-transform
the specification limits for interpretation. AIAG SPC chapter 4,
section 4.7 prescribes the Box-Cox transformation route when the
skewness exceeds $\pm 1$.

\subsection{Applied vignette: average run length and the economics
of control-chart limits}

A control chart flags a process as out of control when a plotted
statistic falls outside its control limits. The engineering
question the raw chart does not answer is how often the chart will
raise a false alarm when the process is fine, and how quickly it
will catch a real shift when the process is not. The
\engterm{average run length} (ARL) makes both answerable and turns
the choice of control-limit width from a convention into an
economic decision.

The ARL is the expected number of samples between signals. For a
Shewhart $\bar{X}$ chart with limits at $\pm L\sigma_{\bar{x}}$ and
independent samples, the probability that a single plotted point
falls outside the limits when the process is on target is
$p_{0} = 2\Phi(-L)$, and the in-control run length is geometric
with mean
\[
\text{ARL}_{0} = \frac{1}{p_{0}} = \frac{1}{2\Phi(-L)}.
\]
At the standard $L = 3$ limits, $p_{0} = 2\Phi(-3) \approx 0.0027$,
so $\text{ARL}_{0} \approx 370$: a stable process signals a false
alarm on average once every $370$ samples. Tightening to
$L = 2$ raises $p_{0}$ to $0.0455$ and drops $\text{ARL}_{0}$ to
about $22$, a false alarm roughly seventeen times more often. The
three-sigma convention is not arbitrary; it buys a long quiet
stretch between false alarms.

The out-of-control ARL measures detection speed. When the process
mean shifts by $\delta\sigma$ (in units of the single-observation
standard deviation), a subgroup of size $n$ shifts the plotted mean
by $\delta\sqrt{n}$ standard errors. The probability that a point
now falls outside the limits is
$p_{1} = \Phi(-L + \delta\sqrt{n}) + \Phi(-L - \delta\sqrt{n})$,
and $\text{ARL}_{1} = 1/p_{1}$. For a one-sigma shift
($\delta = 1$) with subgroup size $n = 4$ at $L = 3$, the shift is
$\delta\sqrt{n} = 2$ standard errors, so
$p_{1} = \Phi(-1) + \Phi(-5) \approx 0.159$ and
$\text{ARL}_{1} \approx 6.3$: the chart takes about six subgroups
to catch a one-sigma shift. Doubling the subgroup to $n = 16$
moves the shift to $4$ standard errors, raises $p_{1}$ to about
$0.84$, and drops $\text{ARL}_{1}$ to about $1.2$: near-immediate
detection, at four times the sampling cost.

The economic design of the chart trades the two ARLs against each
other and against cost. A false alarm costs an investigation that
finds nothing (operator time, a halted line, a needless teardown);
a missed shift costs whatever the out-of-control process produces
before the chart catches it (scrap, rework, a field return). The
control-limit width $L$ sets $\text{ARL}_{0}$, the subgroup size
$n$ and sampling interval set $\text{ARL}_{1}$, and the working
design minimises the expected total cost per unit time:
\[
\text{expected cost rate} \approx
\underbrace{\frac{c_{\text{FA}}}{\text{ARL}_{0} \cdot h}}
_{\text{false-alarm rate} \times \text{cost}}
+ \underbrace{c_{\text{shift}} \cdot \frac{\text{ARL}_{1} \cdot h}
{2}}_{\text{expected out-of-control exposure}}
+ \underbrace{\frac{c_{\text{sample}} \cdot n}{h}}
_{\text{sampling cost}},
\]
where $h$ is the sampling interval, $c_{\text{FA}}$ the cost of a
false alarm, $c_{\text{shift}}$ the cost per unit time of running
out of control, and $c_{\text{sample}}$ the per-observation cost.
The formula is a planning sketch, not a production optimiser; the
Duncan economic-design model and its successors, catalogued in
Montgomery \cite[ch.~10]{text:ch08-montgomery-spc-2019}, refine it.
The working point is that $L$ and $n$ are decision variables, not
constants handed down by convention: a line where a false alarm
halts a furnace worth thousands per hour wants a wider $L$ and a
longer $\text{ARL}_{0}$; a line where an undetected shift ships
defective safety parts wants a larger $n$ and a shorter
$\text{ARL}_{1}$, and pays for it in sampling. The
CUSUM and EWMA charts, which accumulate evidence across successive
samples rather than judging each one alone, reach a much shorter
$\text{ARL}_{1}$ for small shifts at the same $\text{ARL}_{0}$, and
are the working choice when the shifts that matter are smaller than
one sigma.

\section{Bayesian intuition for the practitioner}\pathtag{standard}

A frequentist reads probability as long-run frequency: the
probability of a coin landing heads is the fraction of heads in many
tosses. A Bayesian reads probability as degree of belief: the
probability of a hypothesis is how strongly the evidence supports it
\cite[ch.~11]{text:wasserman2004}. With large enough data and weak
prior information, frequentist and Bayesian analyses often give
similar numerical estimates for routine engineering questions. They
still differ in interpretation and in how they represent prior
information. Chapter 4 §4.6 said this in
passing; we now develop it operationally without computing
integrals.

\subsection{Bayes' rule on a small contingency table}

A diagnostic test for a manufacturing defect has the following
characteristics, established by extensive testing on calibration
samples:

\begin{itemize}[leftmargin=*]
  \item Sensitivity (true positive rate): $95 \,\%$. A truly
    defective unit is flagged $95 \,\%$ of the time.
  \item Specificity (true negative rate): $98 \,\%$. A truly good
    unit is cleared $98 \,\%$ of the time.
  \item Prior: $1 \,\%$ of units coming off the line are actually
    defective.
\end{itemize}

A unit tests positive. What is the probability it is actually
defective?

The intuition that says ``$95 \,\%$, since the test is $95 \,\%$
sensitive'' is wrong. The correct calculation works on a population.
Imagine $10{,}000$ units. One hundred of them are defective
($1 \,\%$ of $10{,}000$). Of those hundred, the test flags $95$ on
average ($95 \,\%$ of $100$). The remaining $9{,}900$ are good. Of
those, the test falsely flags $198$ ($2 \,\%$ of $9{,}900$). The
total positives are $95 + 198 = 293$. Of those, $95/293 \approx
0.324$ are actually defective. About one in three positive tests
identifies a real defect; about two in three are false alarms.

The result is direct: with a $1 \,\%$ defect prior, this $95 \,\%$-sensitive
and $98 \,\%$-specific test returns false alarms for about two
thirds of positive results. The intuition
that pinned the probability of true defect at $95 \,\%$ confused the
sensitivity (a property of the test, conditional on the unit being
defective) with the posterior (a property of the unit, conditional
on the test result). The two are different by a factor near three
when the prior is small.

\subsection{The general form}

Bayes' rule, in the discrete case, is

\[
\Pr(H \mid D) =
\frac{\Pr(D \mid H) \Pr(H)}{\Pr(D)},
\]

where $H$ is the hypothesis (the unit is defective), $D$ is the
data (the test was positive), $\Pr(H)$ is the prior probability of
the hypothesis, $\Pr(D \mid H)$ is the likelihood of the data given
the hypothesis, and $\Pr(D)$ is the marginal probability of the
data across all hypotheses. The denominator is the normalising
constant.

In the defect example, $\Pr(H) = 0.01$, $\Pr(D \mid H) = 0.95$,
and the marginal $\Pr(D) = \Pr(D \mid H) \Pr(H) + \Pr(D \mid \bar H)
\Pr(\bar H) = 0.95 \cdot 0.01 + 0.02 \cdot 0.99 = 0.0095 + 0.0198 =
0.0293$. The posterior is

\[
\Pr(H \mid D) = \frac{0.0095}{0.0293} \approx 0.324,
\]

agreeing with the population calculation. The contingency table
that underwrites the same calculation is the working presentation
an engineer should produce when the answer matters:

\begin{center}
\begin{tabular}{lccc}
\toprule
& Test positive & Test negative & Total \\
\midrule
Defective & $95$ & $5$ & $100$ \\
Good & $198$ & $9{,}702$ & $9{,}900$ \\
\midrule
Total & $293$ & $9{,}707$ & $10{,}000$ \\
\bottomrule
\end{tabular}
\end{center}

The table makes each rate visible: the $95$ true positives come
from $100 \cdot 0.95$; the $198$ false positives come from $9{,}900
\cdot 0.02$; the posterior is $95/293 \approx 0.324$. An engineer
who is unsure of a Bayesian calculation should write the contingency
table and check every count.

\subsection{Updating, sequentially}

Bayes' rule supports sequential updating. The posterior after
observing data $D_{1}$ becomes the prior before observing $D_{2}$,
and the posterior after $D_{1}$ and $D_{2}$ is the same whether the
data are processed jointly or one at a time. Bayesian methods are
therefore natural for online estimation, sensor fusion, and
sequential testing.

Continuing the defect example: a unit tests positive. The posterior
probability of defect is $0.324$. The unit is then re-tested with a
second test whose errors are conditionally independent of the first
test given the unit's true state. If it tests positive again, the
new prior is $0.324$ and the new posterior is

\[
\frac{0.95 \cdot 0.324}{0.95 \cdot 0.324 + 0.02 \cdot 0.676}
= \frac{0.3078}{0.3213} \approx 0.958.
\]

After a second independent positive, the probability of true defect
rises to about $96 \,\%$, which is decision-relevant. This is the
principle behind two-step screening in industrial quality control
and in medical diagnostics: a single positive on a low-prior
population is weak evidence; two independent positives are strong.

\subsection{Beta-binomial conjugate inference: pass/fail testing
under a prior}

A second working example moves from a single hypothesis to a
continuous parameter: the unknown defect rate $p$ of a manufacturing
process estimated from a small pass/fail test. The
\engterm{beta distribution} is the conjugate prior for binomial
sampling, which means the posterior is again a beta distribution
with arithmetically updated parameters. The closed-form update
makes beta-binomial inference the working tool for reliability
demonstration, acceptance sampling, and any small-sample
proportion estimate where historical data inform the prior
\cite[ch.~2]{text:ch08-gelman-bda3-2013}.

\textit{Setup.} A vendor delivers a new batch of $n = 20$ optical
connectors. Industry experience across past batches suggests a
defect rate centred near $5 \,\%$ with substantial spread; an
analyst encodes this as a $\text{Beta}(2, 38)$ prior, which has
mean $2/(2 + 38) = 0.05$ and standard deviation
$\sqrt{2 \cdot 38/(40^{2} \cdot 41)} \approx 0.034$. The vendor
sends the batch and the QC team tests all $20$ units; $k = 1$
fails. What is the posterior distribution of the batch defect
rate, and what is the posterior probability that the rate exceeds
the contract threshold of $10 \,\%$?

\textit{Conjugate update.} The binomial likelihood for $k$ failures
in $n$ trials is
$L(p) \propto p^{k} (1 - p)^{n - k}$. The beta prior is
$\pi(p) \propto p^{\alpha - 1} (1 - p)^{\beta - 1}$. The posterior
is proportional to their product:
\[
\pi(p \mid k, n) \propto p^{k + \alpha - 1} (1 - p)^{n - k + \beta - 1},
\]
which is $\text{Beta}(k + \alpha, n - k + \beta)$. The update is
arithmetic: add observed failures to the first shape parameter,
add observed successes to the second.

\textit{Computation.} With $\alpha = 2$, $\beta = 38$, $k = 1$,
$n - k = 19$, the posterior is $\text{Beta}(3, 57)$. The
posterior mean is $3/60 = 0.050$, almost unchanged from the
prior because one failure in twenty agrees with a $5 \,\%$ prior.
The posterior standard deviation is
\[
\sqrt{\frac{3 \cdot 57}{60^{2} \cdot 61}} \approx 0.028,
\]
slightly smaller than the prior's $0.034$. The data sharpened the
estimate but did not move it.

\textit{Credible interval.} A $95 \,\%$ central credible interval
runs from the $2.5$th to the $97.5$th percentile of
$\text{Beta}(3, 57)$, which are approximately $0.011$ and $0.117$.
The posterior probability that the true defect rate lies in
$[0.011, 0.117]$ is exactly $0.95$ by construction. By contrast,
a frequentist Wilson interval at the same data gives roughly
$[0.009, 0.236]$, considerably wider on the upper side because the
frequentist analysis has no prior to anchor extreme rates.

\textit{Contract decision.} The contract threshold is
$p \leq 0.10$. The posterior probability that $p > 0.10$ is
$\Pr(\text{Beta}(3, 57) > 0.10)$, evaluated as
$1 - I_{0.10}(3, 57) \approx 0.044$, where $I_{x}(a, b)$ is the
regularised incomplete beta function. The team can therefore
state that the posterior probability of out-of-spec defect rate is
about $4.4 \,\%$, well below a typical $10 \,\%$ decision
threshold; the batch is accepted.

\textit{Sequential updating.} A second batch arrives and tests
$k_{2} = 4$ failures in $n_{2} = 30$. Use the previous posterior
as the new prior: $\text{Beta}(3, 57)$ becomes
$\text{Beta}(3 + 4, 57 + 26) = \text{Beta}(7, 83)$. Posterior mean
$7/90 \approx 0.078$, posterior standard deviation $\approx
0.028$. The second batch has shifted the estimate upward by about
half a standard deviation. The posterior probability of exceeding
$10 \,\%$ rises to about $0.21$; the contract decision becomes
borderline and triggers a third-batch sampling plan.

\textit{Prior sensitivity.} The choice of prior matters when data
are sparse. A flat $\text{Beta}(1, 1)$ prior (uniform on $[0, 1]$)
after the first batch's $1$-in-$20$ gives posterior $\text{Beta}(2,
20)$, mean $0.091$, $\Pr(p > 0.10) \approx 0.41$. The flat prior
inflates the posterior because it places mass on rates the
historical experience rules out. A more informative
$\text{Beta}(5, 95)$ prior (mean $0.05$, sd $\approx 0.022$)
after $1$-in-$20$ gives $\text{Beta}(6, 114)$, mean $0.050$, sd
$\approx 0.020$, $\Pr(p > 0.10) \approx 0.003$. The reporting
discipline is to state the prior explicitly and to run the
analysis under a default-flat alternative as a sensitivity check.

\textit{When the conjugate form helps.} Beta-binomial closed-form
updates are useful for sequential acceptance sampling, where each
test result revises the posterior in one arithmetic step;
reliability demonstration with limited test units, where industry
experience can be encoded in $\alpha$ and $\beta$; and online
quality monitoring, where the running posterior is the basis for
disposition decisions. When the likelihood is not binomial or the
prior is not conjugate, the closed-form trick fails and the
analyst turns to Markov-chain Monte Carlo (MCMC). MCMC samples
from the posterior by running a Markov chain whose stationary
distribution is the posterior; standard implementations
(\textsf{Stan}, \textsf{PyMC}, \textsf{JAGS}) handle the
mechanics. The conceptual move from conjugate updating to MCMC is
the move from arithmetic to sampling, and Volume IX treats it in
detail.

\textit{Common errors.} (i) Choosing a strongly informative prior
because it confirms the conclusion the analyst wanted, then
treating the posterior as objective. (ii) Forgetting that the
posterior collapses to the prior when no data are observed; an
empty test yields the prior, not a defensible estimate. (iii)
Confusing a credible interval (a Bayesian statement about $p$
given the data) with a confidence interval (a frequentist
procedure with stated long-run coverage).

\textit{Reporting.} \emph{Batch 1: prior $\text{Beta}(2, 38)$,
data $1$/$20$, posterior $\text{Beta}(3, 57)$, posterior mean
$0.050$, $95 \,\%$ credible interval $[0.011, 0.117]$,
$\Pr(p > 0.10) = 0.044$; batch accepted under prior, data, and
threshold all stated.} The vendor and customer can audit the
prior choice as well as the data.

\subsection{Priors as engineering judgement}

The prior in a Bayesian analysis is the engineer's claim about the
plausibility of each hypothesis before the data arrive. In the
defect example, the prior of $1 \,\%$ came from production records.
In other settings the prior comes from previous studies, similar
processes, physical models, or expert judgement.

A non-informative prior tries to encode ignorance: a uniform
distribution over a parameter, a wide Gaussian, or one of several
formal choices. With enough data the posterior becomes nearly
independent of the prior, which is the operational sense in which
the frequentist and Bayesian numerical answers converge for routine
engineering questions. The Bayesian framework's distinctive
contribution is its handling of small samples, expert prior
information, and decision-making with explicit utility functions.
For working calculations, an engineer who has used both frameworks
chooses the one whose assumptions match the problem.

\subsection{When to reach for Bayesian thinking}

Three signals point to Bayesian methods:

\begin{itemize}[leftmargin=*]
  \item Strong prior information that the engineer trusts more than
    the small dataset they have just collected. A new sensor whose
    calibration is being verified by ten readings is a Bayesian
    problem; the prior comes from the manufacturer's calibration
    history.
  \item Sequential updating with online data, where new measurements
    arrive one at a time and the engineer wants to revise their
    estimate continuously.
  \item Decision-making with explicit consequences, where the cost
    of being wrong differs across hypotheses and the engineer wants
    to integrate over the posterior to find the action with minimum
    expected cost.
\end{itemize}

The framework recurs in Volume IX (systems and control) and Volume
X (failure analysis), where Bayesian network models and Bayesian
reliability updating are the standard tools.

\subsection{Working conjugate priors}

A \engterm{conjugate prior} is a prior distribution family such
that, combined with the likelihood, the posterior remains in the
same family. The combination produces closed-form posterior
parameters: no numerical integration required. The working list:

\begin{itemize}[leftmargin=*]
  \item \textit{Beta-Binomial}: prior $\text{Beta}(\alpha, \beta)$
    on a probability $\theta \in [0, 1]$, likelihood $k$ successes
    in $n$ Bernoulli trials, posterior $\text{Beta}(\alpha + k,
    \beta + n - k)$. The prior parameters $(\alpha, \beta)$ are
    interpretable as $(\alpha - 1)$ prior successes and
    $(\beta - 1)$ prior failures. The uninformative Jeffreys prior
    is $\text{Beta}(0.5, 0.5)$; the uniform prior is
    $\text{Beta}(1, 1)$. The working pass-fail quality-control
    case from the chapter's worked example is a beta-binomial
    update.
  \item \textit{Gamma-Poisson}: prior $\text{Gamma}(\alpha, \beta)$
    on a rate $\lambda > 0$, likelihood Poisson counts $\sum k_{i}
    = K$ over total exposure $T$, posterior $\text{Gamma}(\alpha +
    K, \beta + T)$. The reliability-engineering MTBF estimate from
    field-failure data is a gamma-Poisson update.
  \item \textit{Normal-Normal (known variance)}: prior $\mathcal{N}
    (\mu_{0}, \tau_{0}^{2})$ on a mean $\mu$ with known
    observation variance $\sigma^{2}$, posterior $\mathcal{N}
    (\mu_{n}, \tau_{n}^{2})$ where $1/\tau_{n}^{2} = 1/\tau_{0}^{2}
    + n/\sigma^{2}$ and $\mu_{n} = \tau_{n}^{2} (\mu_{0}/\tau_{0}^{2}
    + n\bar{x}/\sigma^{2})$. The Kalman filter (Volume IX
    chapter 7) is the sequential-application form.
  \item \textit{Normal-Inverse-Gamma (unknown mean and variance)}:
    prior on $(\mu, \sigma^{2})$ jointly. The hyperparameters
    $(\mu_{0}, \kappa_{0}, \alpha_{0}, \beta_{0})$ encode the
    prior estimates and confidence. Working choice for
    measurement-model parameter estimation when the variance is
    itself uncertain.
  \item \textit{Dirichlet-Multinomial}: prior $\text{Dir}
    (\alpha_{1}, \ldots, \alpha_{k})$ on a categorical
    distribution, posterior $\text{Dir}(\alpha_{i} + n_{i})$ after
    observing $n_{i}$ in category $i$. Used for defect-mode
    classification and any categorical-count problem.
\end{itemize}

When the likelihood is not from a conjugate-prior family, the
posterior cannot be written in closed form. The working
alternative is \engterm{Markov chain Monte Carlo} (MCMC): a
random-walk algorithm that draws samples from the posterior
proportional to its density without requiring the normalising
constant. The working implementations:

\engterm{Stan} (free, BSD-3-Clause licence, developed at Columbia
University by Andrew Gelman and collaborators since 2012): a probabilistic
programming language with a No-U-Turn Sampler (NUTS), an
adaptive variant of Hamiltonian Monte Carlo. Stan models are
written in a domain-specific language compiled to C\texttt{++};
interface libraries provide access from R (\texttt{rstan},
\texttt{cmdstanr}), Python (\texttt{cmdstanpy}, \texttt{pystan}),
Julia, MATLAB, and other languages. Stan is the working choice
for hierarchical models with continuous parameters and is the
de facto reference for Bayesian inference in academic
statistics.

\engterm{PyMC} (formerly PyMC3, then PyMC v4 and v5, MIT licence):
a Python-native probabilistic programming library that compiles
models to PyTensor (formerly Aesara) for autodifferentiation and
runs NUTS by default. PyMC trades slightly less efficient NUTS
than Stan for tighter Python integration and better discrete-
parameter handling.

\engterm{JAGS} (Just Another Gibbs Sampler, GPL): a Gibbs-sampler
implementation that handles many small models that Stan would
struggle with (discrete latent variables, models with conjugate
substructure). The working interface library is \texttt{rjags} in
R.

\engterm{BUGS} (Bayesian inference Using Gibbs Sampling): the
ancestor of all modern Bayesian software, developed at the MRC
Biostatistics Unit in Cambridge from 1989. WinBUGS and OpenBUGS
are the working implementations. BUGS introduced the
domain-specific-language approach to Bayesian model
specification that Stan and PyMC inherited.

\engterm{TensorFlow Probability and PyTorch's Pyro}: probabilistic
programming libraries built on top of automatic-differentiation
frameworks. Used when the Bayesian model is part of a
deep-learning pipeline; less efficient than Stan or PyMC for
classical Bayesian hierarchical models but seamless with neural-
network components.

Working diagnostic for MCMC: the Gelman-Rubin
$\hat{R}$ statistic (potential scale reduction factor) compares
within-chain variance to between-chain variance across multiple
independent chains; $\hat{R} < 1.01$ is the working target
acceptance threshold per the current Stan best-practice
\cite{text:ch08-gelman-bda3-2013}. The effective sample size
(ESS) $N_{\text{eff}}$ measures the equivalent number of
independent samples after accounting for chain autocorrelation;
ESS per parameter should exceed $400$ for stable posterior
percentile estimates and $1000$ for tail-probability work.

\subsection{A second beta-binomial worked example: sensor
reliability with an informative prior}

A semiconductor fabrication line ships pressure sensors. Field
experience across the prior model year says about $1 \,\%$ of
shipped sensors return as defective within the warranty period,
with the rate stable across batches. The reliability engineer
encodes that experience as a $\text{Beta}(2, 200)$ prior on the
defect rate $p$: mean $2/202 \approx 0.0099$, standard deviation
$\sqrt{2 \cdot 200/(202^{2} \cdot 203)} \approx 0.0070$. The prior
asserts $p \approx 1 \,\%$ with $\pm 0.7 \,\%$ uncertainty,
matching the long-run record. A new batch of $n = 50$ sensors goes
through a stress screen. One sensor fails ($k = 1$).

\textit{Conjugate update.} The posterior is $\text{Beta}(2 + 1,
200 + 49) = \text{Beta}(3, 249)$. Posterior mean $3/252 \approx
0.0119$, posterior standard deviation
$\sqrt{3 \cdot 249/(252^{2} \cdot 253)} \approx 0.0068$. The $95
\,\%$ central credible interval runs from the $2.5$th to the
$97.5$th percentile of $\text{Beta}(3, 249)$, approximately
$[0.0040, 0.0331]$, or $0.4 \,\%$ to $3.3 \,\%$.

\textit{Frequentist comparison.} The Wilson score interval for
$1/50$ at $95 \,\%$ is approximately $[0.0003, 0.0568]$, or
$0.03 \,\%$ to $5.7 \,\%$. The Clopper-Pearson exact interval is
$[0.0005, 0.1065]$, or $0.05 \,\%$ to $10.7 \,\%$. Both
frequentist intervals are much wider on the upper side because
they have no historical context to rule out the $5 \,\%$-or-higher
defect-rate hypothesis. The Bayesian interval is asymmetric in the
other direction because the prior anchors the bulk of plausible
defect rates near $1 \,\%$.

\textit{Engineering implication.} The contract specification is
$p \leq 2 \,\%$. The posterior probability $\Pr(p > 0.02)$ is
$1 - I_{0.02}(3, 249) \approx 0.052$ (here $I_{x}(a, b)$ is the
regularised incomplete beta function). The reliability engineer
reports: prior $\text{Beta}(2, 200)$ (justified by 2024 field
data, current as of 2025), data $1/50$, posterior $\text{Beta}(3,
249)$, $95 \,\%$ credible interval $[0.0040, 0.0331]$, $\Pr(p >
0.02) = 0.052$. The contract threshold is borderline at the
posterior; the team triggers a second batch sampling plan rather
than accepting on the first batch alone. A frequentist analysis
at the same data would have reported a $1/50 = 2 \,\%$ point
estimate with an upper $97.5$th-percentile bound near $11 \,\%$,
which both wastes the prior information and forces a more
conservative downstream decision than the data warrant.

\subsection{Gamma-Poisson worked example: a failure-rate estimate
from field exposure}

The beta-binomial pair estimates a probability from pass/fail
counts. When the data are event counts accumulated over a measured
exposure rather than successes out of a fixed number of trials, the
conjugate pair is gamma-Poisson. The unknown quantity is a rate
$\lambda$: failures per unit of time, defects per unit of area,
arrivals per unit of duration. The \engterm{gamma distribution} is
the conjugate prior for the Poisson rate, and the update is again
arithmetic \cite[ch.~2]{text:ch08-gelman-bda3-2013}.

\textit{Setup.} A pump-seal supplier tracks warranty failures
across a fleet. Historical experience across the prior model year
gives a mean failure rate near $0.10$ failures per
$10^{3}\,\si{\hour}$ of operation, with substantial batch-to-batch
spread. The reliability engineer encodes this as a
$\text{Gamma}(\alpha = 2, \beta = 20)$ prior on $\lambda$ (with
$\lambda$ in failures per $10^{3}\,\si{\hour}$), which has mean
$\alpha/\beta = 2/20 = 0.10$ and standard deviation
$\sqrt{\alpha}/\beta = \sqrt{2}/20 \approx 0.071$. A new production
lot accumulates $T = 40$ units of exposure (each unit is
$10^{3}\,\si{\hour}$ of operation) and returns $K = 6$ failures.
What is the posterior rate, and what is the posterior probability
that the rate exceeds the contract ceiling of $0.20$ failures per
$10^{3}\,\si{\hour}$?

\textit{Conjugate update.} The Poisson likelihood for $K$ events
over exposure $T$ is $L(\lambda) \propto \lambda^{K}
e^{-\lambda T}$. The gamma prior is $\pi(\lambda) \propto
\lambda^{\alpha - 1} e^{-\beta \lambda}$. Their product is
proportional to $\lambda^{K + \alpha - 1} e^{-(T + \beta)\lambda}$,
which is $\text{Gamma}(\alpha + K, \beta + T)$. Add observed
events to the shape parameter, add observed exposure to the rate
parameter.

\textit{Computation.} With $\alpha = 2$, $\beta = 20$, $K = 6$,
$T = 40$, the posterior is $\text{Gamma}(8, 60)$. The posterior
mean is $8/60 \approx 0.133$ failures per $10^{3}\,\si{\hour}$, and
the posterior standard deviation is $\sqrt{8}/60 \approx 0.047$.
The data pulled the estimate up from the prior mean of $0.10$ (the
observed rate $6/40 = 0.15$ is above the prior mean) and sharpened
it (the prior standard deviation $0.071$ fell to $0.047$).
\Cref{fig:vol01:ch08:gamma-poisson-update} shows the prior and
posterior densities.

\begin{figure}[ht]
\centering
\begin{tikzpicture}[x=1cm, y=1cm, font=\small]

% Prior and posterior gamma densities on a failure rate lambda.
% Prior Gamma(2, 1) broad; posterior Gamma(11, 5) tighter and shifted.
% Densities are drawn as hand-shaped smooth curves (schematic, not
% numerically exact) to show the sharpening-and-shifting of the update.

% Axes
\draw[->, thick] (0, 0) -- (9.6, 0)
  node[below right, font=\scriptsize] {failure rate $\lambda$ (per unit-time)};
\draw[->, thick] (0, 0) -- (0, 4.4)
  node[above, font=\scriptsize] {posterior density};

% x tick marks
\foreach \x/\lab in {0/0, 2/1, 4/2, 6/3, 8/4} {
  \draw (\x, 0.08) -- (\x, -0.08) node[below, font=\scriptsize] {\lab};
}

% Prior: broad Gamma, mode near lambda=1 (x=2), long tail
\draw[thick, engblue, smooth]
  plot coordinates {
    (0.15,0.15) (0.8,0.75) (1.5,1.35) (2.0,1.55) (2.6,1.45)
    (3.4,1.15) (4.2,0.85) (5.2,0.55) (6.4,0.32) (7.6,0.17) (8.8,0.08)
  };
\node[engblue, anchor=west, font=\scriptsize] at (4.4, 0.95)
  {prior $\text{Gamma}(2,1)$};

% Posterior: tighter, mode near lambda=2 (x=4)
\draw[thick, engred, smooth]
  plot coordinates {
    (1.2,0.05) (2.2,0.55) (3.0,1.85) (3.6,3.15) (4.0,3.65)
    (4.4,3.35) (5.0,2.35) (5.8,1.15) (6.6,0.45) (7.4,0.16) (8.2,0.05)
  };
\node[engred, anchor=west, font=\scriptsize] at (5.6, 2.55)
  {posterior $\text{Gamma}(11,5)$};

\end{tikzpicture}
\caption[Gamma-Poisson conjugate update of a failure rate]{The
gamma-Poisson conjugate update of section 8.5, drawn schematically.
The engineer starts with a broad $\text{Gamma}(2,1)$ prior on the
failure rate $\lambda$ (blue), encoding weak belief that the rate is
near one failure per unit-time. After observing nine failures over
four units of exposure, the posterior is $\text{Gamma}(2+9,\,1+4) =
\text{Gamma}(11,5)$ (red): shifted toward the data's implied rate of
$9/4 = 2.25$ and sharpened by the added information. The update is
arithmetic, add observed counts to the first parameter and observed
exposure to the second, so no numerical integration is required. The
posterior mean $11/5 = 2.2$ sits between the prior mean and the raw
data rate, weighted by the prior's and the data's precisions.}
\label{fig:vol01:ch08:gamma-poisson-update}
\end{figure}


\textit{Credible interval and contract decision.} A $95 \,\%$
central credible interval for $\text{Gamma}(8, 60)$ runs from the
$2.5$th to the $97.5$th percentile, approximately $[0.058, 0.246]$
failures per $10^{3}\,\si{\hour}$. The posterior probability that
the rate exceeds the contract ceiling of $0.20$ is
$\Pr(\text{Gamma}(8, 60) > 0.20) \approx 0.087$, so the team can
state that the ceiling is exceeded with about $9 \,\%$ posterior
probability. The batch is accepted under a typical $10 \,\%$
threshold, but the margin is thinner than the pass/fail phrasing
alone would suggest, and a purely frequentist rate estimate of
$6/40 = 0.15 \pm 0.061$ (the Poisson standard error
$\sqrt{K}/T$) would have carried no prior context to weigh against
the ceiling.

\textit{Mean-time-between-failures reading.} The rate estimate
inverts directly to a mean time between failures. The posterior
mean rate $0.133$ failures per $10^{3}\,\si{\hour}$ corresponds to
an MTBF of $1/0.133 \approx 7.5 \times 10^{3}\,\si{\hour}$; the
upper credible bound $0.246$ corresponds to an MTBF floor of
about $4.1 \times 10^{3}\,\si{\hour}$. A maintenance planner who
sizes a spares budget uses the MTBF floor, not the point estimate,
so that the spares stock is not undersized when the true rate sits
at the pessimistic end of the credible interval.

\textit{Reporting.} \emph{Prior $\text{Gamma}(2, 20)$ (2024 fleet
data, current as of 2025), exposure $T = 40$ unit-$10^{3}\,$h,
$K = 6$ failures, posterior $\text{Gamma}(8, 60)$, mean $0.133$
failures per $10^{3}\,$h, $95 \,\%$ credible interval
$[0.058, 0.246]$, $\Pr(\lambda > 0.20) = 0.087$; MTBF
$\approx 7.5 \times 10^{3}\,$h, floor $\approx 4.1 \times
10^{3}\,$h.} The reader can audit the prior, the exposure, and the
event count.

\subsection{Dirichlet-multinomial worked example: defect-mode
mix estimation}

A third conjugate pair handles categorical counts. When a
manufacturing failure is classified into one of several mutually
exclusive modes, the unknown quantity is the vector of mode
proportions $\boldsymbol\theta = (\theta_{1}, \ldots, \theta_{m})$
summing to one. The \engterm{Dirichlet distribution} is the
conjugate prior for multinomial category counts, and the update is
arithmetic in the same way the beta-binomial and gamma-Poisson
updates are \cite[ch.~3]{text:ch08-gelman-bda3-2013}.

\textit{Setup.} A returned-goods analysis classifies each field
failure of a connector into one of four modes: contact corrosion,
solder-joint crack, housing fracture, or contamination. A prior
$\text{Dir}(2, 2, 2, 2)$ encodes weak symmetry (no mode favoured,
each with an effective prior count of one observation). A batch of
$n = 60$ field returns is classified, giving observed counts
$\boldsymbol n = (31, 14, 9, 6)$ across the four modes.

\textit{Conjugate update.} The multinomial likelihood is
$L(\boldsymbol\theta) \propto \prod_{j=1}^{m} \theta_{j}^{n_{j}}$
and the Dirichlet prior is $\pi(\boldsymbol\theta) \propto
\prod_{j=1}^{m} \theta_{j}^{\alpha_{j} - 1}$. Their product is
$\text{Dir}(\alpha_{1} + n_{1}, \ldots, \alpha_{m} + n_{m})$. Add
the observed count in each category to that category's prior
parameter.

\textit{Computation.} The posterior is
$\text{Dir}(33, 16, 11, 8)$, with total concentration
$33 + 16 + 11 + 8 = 68$. The posterior mean for mode $j$ is
$(\alpha_{j} + n_{j})/68$: contact corrosion
$33/68 \approx 0.485$, solder-joint crack $16/68 \approx 0.235$,
housing fracture $11/68 \approx 0.162$, contamination
$8/68 \approx 0.118$. The posterior standard deviation for the
dominant mode is
\[
\sqrt{\frac{0.485(1 - 0.485)}{68 + 1}} \approx 0.060,
\]
using the Dirichlet marginal-variance formula
$\theta_{j}(1 - \theta_{j})/(A + 1)$ with $A = 68$ the posterior
concentration. A $95 \,\%$ credible interval on the corrosion
share is roughly $0.485 \pm 2 \cdot 0.060 = [0.365, 0.605]$.

\textit{Engineering reading.} Contact corrosion is the dominant
mode with posterior mean near $0.49$ and a credible interval that
excludes $0.35$, so corrosion accounts for something between a
third and three-fifths of field failures. The corrective-action
budget follows the posterior mean weighted by the per-mode repair
cost, not the raw counts: a rare mode with a catastrophic
consequence can outrank a common mode with a cheap fix once the
per-mode cost multiplies the proportion. The Dirichlet posterior
supplies the proportion and its uncertainty; the cost model
supplies the weighting.

\textit{Common errors.} (i) Treating the raw fractions
$31/60, 14/60, \ldots$ as certain, ignoring that $n = 60$ leaves
each proportion uncertain to several percentage points. (ii)
Comparing two modes by their point estimates without checking
whether their credible intervals overlap; solder-joint crack
($0.235$) and housing fracture ($0.162$) have overlapping
intervals and are not cleanly separated at $n = 60$. (iii)
Choosing a strongly informative Dirichlet prior that encodes the
conclusion the analyst expected.

\textit{Reporting.} \emph{Prior $\text{Dir}(2, 2, 2, 2)$, counts
$(31, 14, 9, 6)$ over $n = 60$, posterior
$\text{Dir}(33, 16, 11, 8)$; posterior mean mode mix
$(0.49, 0.24, 0.16, 0.12)$; corrosion $95 \,\%$ credible interval
$[0.37, 0.60]$.} The reader can audit the prior and the counts and
recompute every proportion.

\subsection{Sequential testing for go/no-go gates}

A working pattern in acceptance sampling is to define a stopping
rule on the posterior tail probability and to update the posterior
after each test result. The rule has the form:

\begin{itemize}[leftmargin=*]
  \item \emph{Stop and accept} when $\Pr(p > p_{\text{spec}}) <
    \alpha_{\text{accept}}$.
  \item \emph{Stop and reject} when $\Pr(p > p_{\text{spec}}) >
    \alpha_{\text{reject}}$.
  \item \emph{Continue} otherwise.
\end{itemize}

Concrete values for the sensor case: $p_{\text{spec}} = 0.02$,
$\alpha_{\text{accept}} = 0.05$ (high confidence that the rate is
below spec), $\alpha_{\text{reject}} = 0.50$ (better than even
odds that the rate exceeds spec). The stopping boundary is
asymmetric on purpose: rejection is cheap (return the batch),
acceptance has long-tail consequences (warranty exposure).

The running calculation after each test result updates the
posterior arithmetically. Test result $1$: prior $\text{Beta}(2,
200)$, result pass, posterior $\text{Beta}(2, 201)$, $\Pr(p >
0.02) \approx 0.038$. The boundary has already been crossed in
the ``accept'' direction after one pass: the prior alone, combined
with a single pass, brings the tail below $5 \,\%$. The engineer
may choose to require at least $n = 50$ tests regardless (to
catch a degenerate sensor batch where every unit fails); the
stopping rule then becomes ``stop only after $n \geq n_{\min}$
and the posterior boundary is crossed.''

A working version of this rule appeared in the United States
Department of Defense MIL-STD-105E acceptance-sampling tables
\cite{std:ch08t3b-mil-std-105e-1989}, which give attribute-by-
attribute accept/reject numbers as a function of lot size and
inspection level. The Bayesian formulation generalises the table
to arbitrary priors and to continuous stopping rules; the
MIL-STD-105E numbers are recovered as a special case under a flat
prior and fixed sample size. ANSI/ASQ Z1.4-2003
\cite{std:ch08t3b-ansi-z1-4-2003} is the modern U.S.\ equivalent
adopted after MIL-STD-105E was cancelled in 1995.

\textit{Stopping rule arithmetic, second example.} A batch of
critical-application sensors arrives with no prior history; the
engineer uses a weakly informative $\text{Beta}(1, 1)$ prior
(uniform on $[0, 1]$). Spec $p_{\text{spec}} = 0.05$,
$\alpha_{\text{accept}} = 0.05$. After $20$ pass, $0$ fail:
posterior $\text{Beta}(1, 21)$, $\Pr(p > 0.05) = 0.659$. Continue.
After $50$ pass, $0$ fail: posterior $\text{Beta}(1, 51)$,
$\Pr(p > 0.05) = 0.073$. Still continue. After $60$ pass, $0$
fail: posterior $\text{Beta}(1, 61)$, $\Pr(p > 0.05) = 0.046$.
Boundary crossed; accept the batch. The cost of the weakly
informative prior is roughly $60$ tests; the cost of the prior
$\text{Beta}(2, 200)$ was the single test of the previous example.
The trade-off the engineer manages is the credibility of the
prior against the cost of testing.

\subsection{Bayesian linear regression vs OLS}

Ordinary least squares (OLS) regression returns point estimates of
slope and intercept and frequentist standard errors. Bayesian
linear regression returns a posterior distribution over slope and
intercept, given a prior. The two coincide numerically when the
prior is uninformative; they diverge when the prior is informative
or when the sample is small relative to the prior's strength.

\textit{Setup.} A measurement system relates a sensor reading $y$
to a known reference value $x$ through the linear model
$y = \beta_{0} + \beta_{1} x + \varepsilon$, with $\varepsilon
\sim \mathcal{N}(0, \sigma^{2})$. The OLS estimates are
$\hat{\beta}_{0}, \hat{\beta}_{1}$ minimising $\sum (y_{i} -
\beta_{0} - \beta_{1} x_{i})^{2}$, with closed-form variance
estimates from the residual sum of squares.

The conjugate Bayesian formulation uses a normal-inverse-gamma
prior on $(\boldsymbol\beta, \sigma^{2})$. Let
$\boldsymbol\beta = (\beta_{0}, \beta_{1})^{T}$. The prior is
\[
\sigma^{2} \sim \text{Inv-Gamma}(a_{0}, b_{0}),
\quad
\boldsymbol\beta \mid \sigma^{2} \sim \mathcal{N}(\boldsymbol
\mu_{0}, \sigma^{2} \boldsymbol\Lambda_{0}^{-1}),
\]
where $\boldsymbol\mu_{0}$ is the prior mean of the coefficient
vector and $\boldsymbol\Lambda_{0}$ is the prior precision matrix.
With design matrix $\boldsymbol X$ and response vector
$\boldsymbol y$, the posterior is normal-inverse-gamma with
updated hyperparameters:
\[
\boldsymbol\Lambda_{n} = \boldsymbol X^{T} \boldsymbol X +
\boldsymbol\Lambda_{0},
\quad
\boldsymbol\mu_{n} = \boldsymbol\Lambda_{n}^{-1}(\boldsymbol X^{T}
\boldsymbol y + \boldsymbol\Lambda_{0} \boldsymbol\mu_{0}),
\]
\[
a_{n} = a_{0} + n/2,
\quad
b_{n} = b_{0} + \tfrac{1}{2}\left(\boldsymbol y^{T} \boldsymbol y
+ \boldsymbol\mu_{0}^{T} \boldsymbol\Lambda_{0} \boldsymbol\mu_{0}
- \boldsymbol\mu_{n}^{T} \boldsymbol\Lambda_{n} \boldsymbol
\mu_{n}\right).
\]
The posterior mean $\boldsymbol\mu_{n}$ is a precision-weighted
combination of the prior mean and the OLS estimate
$(\boldsymbol X^{T} \boldsymbol X)^{-1} \boldsymbol X^{T}
\boldsymbol y$. When $\boldsymbol\Lambda_{0} \to 0$ (flat prior),
$\boldsymbol\mu_{n}$ approaches the OLS estimate; when the prior
is strong, the posterior is pulled toward $\boldsymbol\mu_{0}$.

\textit{When Bayesian regression helps.} Three settings make the
Bayesian framework worth the additional bookkeeping:

\begin{itemize}[leftmargin=*]
  \item \textit{Informative prior on the slope.} A new sensor of
    a known family carries calibration history from previous
    units: the slope is centred near $1$ with a tight prior
    standard deviation. A small calibration dataset on the new
    unit combined with the informative prior gives a tighter
    posterior than OLS on the small dataset alone.
  \item \textit{Small samples.} OLS standard errors scale as
    $\sqrt{1/n}$ and are unbiased asymptotically; the posterior
    standard deviation under an informative prior is tighter at
    small $n$ because the prior provides additional effective
    sample size.
  \item \textit{Shrinkage to a physical baseline.} Ridge
    regression and lasso regression are Bayesian regressions
    with specific shrinkage priors (normal and Laplace,
    respectively). The penalty parameter $\lambda$ in ridge
    regression corresponds to a prior precision; the engineer
    who tunes $\lambda$ by cross-validation is tuning the
    prior strength.
\end{itemize}

\textit{When OLS is fine.} For routine calibration with $n \geq
30$ and no prior information stronger than a default-flat,
OLS and Bayesian regression give numerically identical fits.
The Bayesian framework adds value only when there is real prior
information or when the analyst wants a full posterior over
$(\beta_{0}, \beta_{1}, \sigma^{2})$ jointly for downstream
decision-making.

\subsection{Computational Bayes in practice: a Stan example}

A working Stan model for the linear regression above is short.
The Stan source file specifies the data, parameters, model, and
generated quantities:

\begin{verbatim}
// linear-regression.stan
data {
  int<lower=1> n;
  vector[n] x;
  vector[n] y;
  real beta0_prior_mean;
  real<lower=0> beta0_prior_sd;
  real beta1_prior_mean;
  real<lower=0> beta1_prior_sd;
}
parameters {
  real beta0;
  real beta1;
  real<lower=0> sigma;
}
model {
  beta0 ~ normal(beta0_prior_mean, beta0_prior_sd);
  beta1 ~ normal(beta1_prior_mean, beta1_prior_sd);
  sigma ~ exponential(1);
  y ~ normal(beta0 + beta1 * x, sigma);
}
generated quantities {
  vector[n] y_rep;
  for (i in 1:n)
    y_rep[i] = normal_rng(beta0 + beta1 * x[i], sigma);
}
\end{verbatim}

\noindent
The model is roughly fifteen lines, the same length as the OLS
implementation in base R but with explicit prior specifications.
The R interface compiles and runs the model:

\begin{verbatim}
library(cmdstanr)
mod <- cmdstan_model("linear-regression.stan")
fit <- mod$sample(
  data = list(n = 30, x = x_obs, y = y_obs,
              beta0_prior_mean = 0, beta0_prior_sd = 1,
              beta1_prior_mean = 1, beta1_prior_sd = 0.1),
  chains = 4, parallel_chains = 4,
  iter_warmup = 1000, iter_sampling = 1000,
  seed = 42)
fit$summary()
\end{verbatim}

\noindent
The four chains run in parallel; the warmup phase tunes the NUTS
sampler; the sampling phase produces $4 \times 1000 = 4000$
posterior draws of each parameter. The summary output reports the
posterior mean, posterior standard deviation, $5$th, $50$th, and
$95$th percentiles, the Gelman-Rubin $\hat{R}$ diagnostic, and the
effective sample size for each parameter.

\textit{Interpreting the trace plot.} A trace plot is a time
series of posterior draws for each parameter, with one line per
chain. The convergence criterion is that the chains should
overlap and look like white noise around a stable mean: four
caterpillars stacked on top of each other. Divergent transitions
(rare but informative warnings from NUTS) flag regions of the
posterior where the sampler had numerical trouble; persistent
divergences require model reparametrisation. The Stan manual
\cite{text:ch08t3b-stan-manual-2024} and Gelman's \emph{Bayesian
Data Analysis} third edition \cite{text:ch08-gelman-bda3-2013}
chapter 11 give the diagnostic protocol.

\textit{Convergence criteria.} Two numerical thresholds are the
working acceptance criteria, per the Stan best-practice document:

\begin{itemize}[leftmargin=*]
  \item $\hat{R} < 1.01$ for every parameter. Higher values
    indicate that the chains have not mixed; rerun with more
    warmup iterations or reparametrise the model.
  \item Effective sample size (ESS) $> 400$ for every parameter.
    Low ESS indicates high autocorrelation within chains; the
    posterior summary statistics are unreliable. Run more
    sampling iterations or thin the chains.
\end{itemize}

Posterior predictive checks complete the diagnostic. The
generated-quantities block above produces $y_{\text{rep}}$, a
posterior predictive draw of the response vector. A plot of the
observed $y$ against the median and $95 \,\%$ band of
$y_{\text{rep}}$ shows whether the model can reproduce the
observed data; systematic discrepancies flag model misspecification.

\section{Sample size and statistical power}\pathtag{standard}

A test that fails to reject the null may do so because the null is
true or because the test had insufficient power to detect a true
effect. The discipline of designing a test before running it is to
choose a sample size large enough that a meaningful effect will be
detected with high probability \cite[ch.~14]{text:devore2015}.

\subsection{The four quantities}

A sample-size calculation for a hypothesis test relates four
quantities, with one fixed by the experiment and the other three
chosen by the engineer:

\begin{itemize}[leftmargin=*]
  \item The significance level $\alpha$, the probability of a false
    positive. Typical: $0.05$.
  \item The power $1 - \beta$, the probability of detecting an
    effect of stated size when one is present. Typical: $0.80$.
  \item The effect size $\delta$, the smallest difference that
    matters for the engineering decision. Set by the application.
  \item The sample size $n$.
\end{itemize}

Fixing any three determines the fourth. The most common direction
is to fix $\alpha$, $1 - \beta$, and $\delta$, and solve for $n$.

\subsection{The two-sample t-test rule of thumb}

For a two-sample $t$-test comparing means with equal sample sizes,
under standard assumptions (Gaussian populations, equal variances,
two-sided test at $\alpha = 0.05$, power $0.80$), the required
sample size per group is approximately

\[
n \approx \frac{16 \sigma^{2}}{\delta^{2}},
\]

where $\sigma$ is the population standard deviation and $\delta$ is
the smallest mean difference the engineer wants to detect. The
factor $16$ comes from the squared sum of the critical values
$z_{\alpha/2} + z_{\beta} \approx 1.96 + 0.84 \approx 2.80$, with
$2 \cdot (2.80)^{2} \approx 15.7$. The reader applies the formula
as a planning approximation. The derivation uses the standard error
of the difference of means and the normal critical values; the
exercises ask the reader to carry that algebra for the
equal-variance case.

\begin{estimation}
\textbf{Question.} A process engineer wants to detect a $1 \,\%$
mean difference in fill weight between two production lines, with
$\alpha = 0.05$ and power $0.80$. The process standard deviation is
$\sigma = 5 \,\si{\gram}$ and the nominal mean is $500 \,\si{\gram}$.
How many units per line are required?

\smallskip
\textit{Estimate, before reading on.}

\smallskip
A $1 \,\%$ mean difference at nominal $500 \,\si{\gram}$ is $\delta
= 5 \,\si{\gram}$. The relevant ratio is $\sigma/\delta = 5/5 = 1$,
so the rule of thumb gives $n \approx 16$ per line. To detect a
$0.1 \,\%$ difference instead ($\delta = 0.5 \,\si{\gram}$,
$\sigma/\delta = 10$), $n \approx 1600$ per line. The cost of
detecting small effects scales as the inverse square of the effect,
which is the central fact of experimental design.

A $10 \,\%$ difference at the same $\sigma$ gives a formal value
$n \approx 0.16$. That number does not mean a two-sample $t$-test
can be run with one observation per group; the test still needs
enough observations to estimate the within-group variance, check
the process, and tolerate at least a small disturbance. A practical
minimum (often $n \geq 5$ to $10$ per group, with more required
when the variance is uncertain) overrides the formula in this
regime: the rule of thumb confirms that the effect is far larger
than the process scatter, not that an inadmissible test is now
admissible.
\end{estimation}

The formula is approximate and assumes Gaussian populations and
equal variances. Welch's correction, unequal sample sizes, and
non-Gaussian distributions all change the constant. Statistical
software returns the exact required $n$ for the chosen test. The
rule of thumb is for planning, not for the final calculation.
\Cref{fig:vol01:ch08:effect-size-vs-n} plots the rule of thumb
against the standardised effect size $d = \delta/\sigma$ (Cohen's
$d$), highlighting the inverse-square scaling: each halving of $d$
quadruples $n$.

\begin{figure}[ht]
\centering
\begin{tikzpicture}[x=1cm, y=1cm, font=\small]

% Sample size n required to reach 80% power at alpha = 0.05 for a
% two-sample t-test, plotted against standardised effect size
% d = delta/sigma. Approx n_per_group = 16 / d^2.

% axes
\draw[->, thick] (0, 0) -- (8.5, 0)
  node[below, font=\scriptsize] {standardised effect size $d = \delta / \sigma$};
\draw[->, thick] (0, 0) -- (0, 5.5)
  node[left, font=\scriptsize, rotate=90, yshift=0.6cm, anchor=south]
  {$n$ per group};

% x-axis ticks for d = 0.1, 0.2, ..., 0.8
\foreach \d/\xc in {0.1/1, 0.2/2, 0.4/4, 0.6/6, 0.8/8} {
  \draw (\xc, 0) -- (\xc, -0.1)
    node[below, font=\scriptsize] {\d};
}
% y-axis ticks: 10, 100, 1000 (log scale)
\foreach \nval/\yc in {10/1, 100/3, 1000/5} {
  \draw (0, \yc) -- (-0.1, \yc)
    node[left, font=\scriptsize] {\nval};
}

% n = 16 / d^2 on log10 y axis: y = log10(16/d^2) - 1, mapped to plot
% units (log10(10) -> y=1, log10(1000) -> y=5, so y = 2*(log10(n)-1)
% + 1 means y plot = log10(n)*2 - 1)
% We sample d at fine increments and plot.
\draw[thick, engblue, smooth]
  plot[domain=0.10:0.85, samples=80]
  (\x*10, {2*(ln(16/(\x*\x))/ln(10)) - 1});

% Annotate three key points
\node[circle, fill=engred, inner sep=1.5pt, label={right:{$d=0.2,\ n\!\approx\!400$}}]
  at (2, {2*(ln(16/(0.2*0.2))/ln(10)) - 1}) {};
\node[circle, fill=engred, inner sep=1.5pt, label={right:{$d=0.5,\ n\!\approx\!64$}}]
  at (5, {2*(ln(16/(0.5*0.5))/ln(10)) - 1}) {};
\node[circle, fill=engred, inner sep=1.5pt, label={right:{$d=0.8,\ n\!\approx\!25$}}]
  at (8, {2*(ln(16/(0.8*0.8))/ln(10)) - 1}) {};

\end{tikzpicture}
\caption[Sample size scales as the inverse square of effect size]{For
a two-sample $t$-test at $\alpha = 0.05$ and power $0.80$, the
required sample size per group is approximately $n \approx 16/d^{2}$
where $d = \delta/\sigma$ is the standardised effect size (Cohen's
$d$). Halving the effect size quadruples the required sample. The
plot shows the rule-of-thumb curve on a logarithmic $n$ axis; an
effect of $d = 0.2$ (small) demands roughly $400$ per group, $d =
0.5$ (medium) about $64$, $d = 0.8$ (large) about $25$. Studies
chasing small effects with small samples are the statistical
mechanism behind much of the failure-to-replicate pattern of
section 8.8.}
\label{fig:vol01:ch08:effect-size-vs-n}
\end{figure}


\paragraph{Worked example: clinical-trial recruitment.} A clinical
team plans a two-arm trial comparing a new sensor against the
standard of care. The smallest difference in mean response time
that would change clinical practice is $d = 0.3$ (small to medium).
The rule of thumb returns $n \approx 16/(0.3)^{2} \approx 178$ per
arm; rounding up for the Welch correction gives $n \approx 190$ per
arm, or $380$ total. Recruitment costs the project approximately
\$2{,}400 per enrolled patient (current as of 2025), so detecting
the $d = 0.3$ effect at $\alpha = 0.05$ and power $0.80$ costs
approximately $380 \cdot \$2{,}400 \approx \$912{,}000$ in
recruitment alone. Reducing the target to $d = 0.2$ raises the
sample to $n \approx 400$ per arm and the recruitment cost to
roughly $\$1.9$ million. The numerical case for fielding the largest
detectable effect that still matters clinically is the
inverse-square scaling, not the analyst's preference.

\begin{estimation}
\textbf{Question.} A quality-engineering team screens $200$
manufacturing process parameters against a single output metric,
running an independent $t$-test for each at $\alpha = 0.05$ with no
multiple-comparison correction. The team plans to flag every
parameter with $p < 0.05$ for follow-up investigation. Assuming the
null is true for every parameter (no real effects in the screen),
how many false positives should the team expect to flag, and what is
the probability of flagging at least one?

\smallskip
\textit{Estimate, before reading on.}

\smallskip
The expected number of false positives is $200 \cdot 0.05 = 10$.
The probability of zero false positives across $200$ independent
tests is $(1 - 0.05)^{200} = 0.95^{200} \approx 3.5 \times
10^{-5}$, so the probability of at least one false positive is
above $0.99997$.
Even if half of the $200$ parameters carry small real effects, the
expected false-positive count is still around $5$, and the engineer
who follows up every flag without correction wastes follow-up
budget on noise.

The remedies match the problem. Bonferroni correction reduces the
per-test significance level to $\alpha/200 = 2.5 \times 10^{-4}$,
which keeps the family-wise error rate at $0.05$ but is
conservative when many real effects are present. Benjamini-Hochberg
control of the false discovery rate
\cite{method:ch08-benjamini-hochberg-1995} keeps the expected
fraction of false positives among the flagged tests below a target
rate (say $0.10$), at the cost of allowing some false positives in
absolute number; it is the default in high-throughput engineering
and biological screening.
\end{estimation}

\subsection{Post-hoc multiple-comparisons procedures}

When the engineer runs an ANOVA on $k$ supplier batches and
rejects the omnibus $F$-test, the next question is which pairs
of batches differ. The naive $\binom{k}{2}$ pairwise
$t$-tests inflate the family-wise error rate as in the
estimation block above. The working procedures differ in what
they protect.

\paragraph{Tukey's honest-significant-difference (HSD).} For
balanced designs with $k$ groups and a common variance, Tukey
HSD uses the studentised range distribution: the critical
value is $q_{\alpha, k, \nu}/\sqrt{2}$ where $\nu$ is the
within-group degrees of freedom. The interval for
$\bar{x}_{i} - \bar{x}_{j}$ is
$\bar{x}_{i} - \bar{x}_{j} \pm q_{\alpha, k, \nu} \cdot
s_{\text{pooled}}/\sqrt{n}$. Worked example: comparing four
supplier batches with $n = 20$ each, $s_{\text{pooled}} = 0.014
\,\si{\milli\meter}$, $\nu = 76$, at $\alpha = 0.05$,
$q_{0.05, 4, 76} = 3.74$. The HSD margin is
$3.74 \times 0.014 /\sqrt{20} = 0.0117\,\si{\milli\meter}$;
any pair differing by more than this is flagged. In R:
\texttt{TukeyHSD(aov(diameter $\sim$ supplier))}; in Python
with \texttt{statsmodels}:
\texttt{pairwise\_tukeyhsd(data, groups, alpha=0.05)}; in SAS:
\texttt{proc glm; class supplier; model diameter = supplier;
means supplier / tukey;}. Tukey HSD is exact for balanced
designs and approximate for unbalanced; it is the default
post-hoc procedure in industrial quality control.

\paragraph{Scheffé's method.} When the engineer wants
intervals for arbitrary contrasts (not just pairwise) or has
not pre-specified which contrasts to test, Scheffé's
adjustment is the conservative choice. The interval width
scales with $\sqrt{(k - 1)F_{\alpha, k-1, \nu}}$ rather than
$q_{\alpha, k, \nu}/\sqrt{2}$. For $k = 4$ at
$\alpha = 0.05$, $\nu = 76$: $\sqrt{3 \times 2.72} = 2.86$, so
the Scheffé margin is $2.86 \times 0.014 \times \sqrt{2/20}
= 0.0127\,\si{\milli\meter}$, slightly wider than Tukey HSD.
For pairwise-only comparisons in a balanced design, Tukey is
tighter; Scheffé wins when the contrast set is large or
unknown in advance.

\paragraph{Dunnett's test.} For comparing $k - 1$ treatment
batches against a single control (rather than all pairs),
Dunnett's procedure uses a multivariate-$t$ critical value
that accounts for the correlation among the $k - 1$ comparisons
sharing the control mean. For $k = 4$ (control plus 3
treatments) at $\alpha = 0.05$, $\nu = 76$: $t_{D, 0.025, 3,
76} \approx 2.41$, compared to Bonferroni's
$t_{0.025/3, 76} \approx 2.45$ and Tukey's $q/\sqrt{2}
\approx 2.65$. Dunnett wins when the question is one-against-
control. In R:
\texttt{glht(model, linfct = mcp(supplier = "Dunnett"))} from
the \texttt{multcomp} package; in SAS:
\texttt{means supplier / dunnett('control');}.

\paragraph{Choice matrix.} For $k$ groups, balanced design,
within-group variance assumption: Tukey HSD for all pairwise,
Dunnett for many-vs-one, Scheffé for arbitrary contrasts,
Bonferroni when independence cannot be assumed or when the
number of comparisons is small. For unbalanced designs or
heteroscedastic data, Games-Howell (Tukey-Kramer extension
with Welch correction) is the working choice; in R:
\texttt{userfriendlyscience::posthocTGH()} or
\texttt{rstatix::games\_howell\_test()}. The ASA
\cite{method:ch08-wasserstein-asa-2016} recommends reporting both the
adjusted $p$-values and the unadjusted effect sizes with
confidence intervals: the engineer needs the effect size to
judge practical significance, and the adjusted $p$-value to
guard against multiple-test inflation.

\subsection{Why power matters}

A study run with insufficient power has two characteristic
failures.

A negative result is uninformative: failure to reject the null could
mean the null is true or the test could not have detected the effect.
The reader who reports ``no significant difference'' from an
under-powered test has reported almost nothing.

A positive result is overstated: when an under-powered study does
reach significance, the estimated effect is biased upward (the only
way to reach $p < \alpha$ with low power is for the random sample to
land on the high tail of the effect's true distribution), and the
result is often unreplicable. This is the statistical mechanism
behind much of what section 8.8 discusses.

\subsection{Pre-registration}

The discipline that makes a sample-size calculation honest is
\engterm{pre-registration}: the engineer specifies, before collecting
data, the hypothesis, the test, the significance level, the target
power, the effect size of interest, the sample size, and the analysis
procedure. The pre-registration is recorded in a notebook, an
internal repository, or a public registry. After data collection,
the analysis follows the pre-registration; deviations are documented
and explained.

Pre-registration reduces the failure modes section 8.7 develops:
selecting analyses to favour the result, redefining hypotheses after
seeing the data, stopping data collection at convenient sample sizes,
and reporting only the analyses that worked. It works only when
deviations are recorded and reported. The framework is independent
of frequentist or Bayesian methods; it is a discipline about what
the engineer commits to before the data are in hand.

\paragraph{Working pre-registration platforms.}
A pre-registration is a timestamped, public or private, immutable
record of the analysis plan before data collection. The working
platforms, current as of 2025:

\begin{itemize}[leftmargin=*]
  \item \textit{Open Science Framework (OSF, \texttt{osf.io}).}
    The de facto standard for cross-discipline pre-registration,
    hosted by the Center for Open Science at the University of
    Virginia \cite{web:ch08t3b-osf-2025}. Free; supports
    public, embargoed, and private registrations; integrates with
    GitHub, Dropbox, and Google Drive for data hosting. The OSF
    pre-registration form has a structured template (hypotheses,
    sample size, design, analysis plan) that the engineer
    completes in roughly fifteen minutes.
  \item \textit{AsPredicted (\texttt{aspredicted.org}).} A
    short-form pre-registration platform developed at the Wharton
    School \cite{web:ch08t3b-aspredicted-2025}. Free; nine
    structured questions; produces a single-page PDF with a
    timestamp. Lighter than OSF for routine industrial work; the
    nine questions cover hypothesis, design, sample size,
    analyses, and exclusion rules.
  \item \textit{ClinicalTrials.gov} for medical-device and
    pharmaceutical trials \cite{web:ch08t3b-clinicaltrials-gov-2025}.
    Mandatory under U.S.\ federal law for FDA-regulated trials of
    drugs and devices via the Food and Drug Administration
    Amendments Act of 2007 (FDAAA, section 801)
    \cite{law:ch08t3b-fdaaa-2007}. The registry records study
    protocol, primary outcomes, sample size, and analysis plan
    before subject enrolment; the World Health Organization runs
    a parallel International Clinical Trials Registry Platform
    (ICTRP).
  \item \textit{American Economic Association RCT Registry}
    \cite{web:ch08t3b-aea-rct-registry-2025} for economic field
    experiments. Free, voluntary; structured template; the
    registry is the working pre-registration record for most
    development-economics field RCTs since approximately 2013.
  \item \textit{Engineering analog.} No discipline-specific
    public registry exists for industrial engineering experiments,
    current as of 2025. The working substitute is internal:
    a signed, dated protocol stored in the project's version-
    control system (Git, SVN, or a document-management system
    with audit log) before data collection begins. The cryptographic
    timestamp of the commit and the audit trail of the repository
    serve the same evidentiary purpose as a public registry. ASTM
    E2935-13 \cite{std:ch08t3b-astm-e2935-2013} gives the working
    standard for experimental protocols in materials testing; the
    pre-registration discipline maps onto its protocol-recording
    requirements.
\end{itemize}

\paragraph{A working pre-registration template for an engineering
experiment.}
The template below is the OSF and AsPredicted structure adapted
for an engineering experiment. The engineer fills it before data
collection and stores the completed document in the project
repository:

\begin{enumerate}[leftmargin=*]
  \item \textit{Title and team.} One-line title; principal
    investigator and team roles; date of pre-registration.
  \item \textit{Hypothesis.} The null hypothesis $H_{0}$ as a
    quantitative statement about a parameter. The alternative
    $H_{1}$ as the complement or a specific direction. Example:
    ``$H_{0}$: the mean bond strength of supplier B equals
    supplier A's; $H_{1}$: supplier B's mean exceeds supplier A's
    by at least $0.05\,\si{\mega\pascal}$.''
  \item \textit{Primary outcome measure.} The single quantitative
    measurement that determines the test result, with units,
    measurement instrument, calibration date, and uncertainty.
    Example: ``shear bond strength of lap-joint specimen,
    measured per ASTM D1002-10 on an MTS 810 servohydraulic
    frame, calibrated 2025-03-14, expanded uncertainty
    $U = 0.02\,\si{\mega\pascal}$ at $k = 2$.''
  \item \textit{Secondary outcomes.} Any additional measurements
    of interest. These are recorded but do not carry the
    confirmatory weight of the primary outcome; secondary
    analyses are exploratory.
  \item \textit{Sample size and power.} The sample size $n$, the
    effect size $\delta$, the significance level $\alpha$, and
    the target power $1 - \beta$. The justification for $\delta$
    in operational units. The power calculation that produces
    $n$ from the other three.
  \item \textit{Exclusion criteria.} The rules that determine
    which observations are included in the analysis. Example:
    ``specimens with visible voids on optical inspection are
    excluded; specimens that slip in the grips are repeated
    once and the repeat measurement is used; specimens where
    the frame load-cell warning fires are repeated and the
    repeat measurement is used.''
  \item \textit{Analysis plan.} The specific test (one-sample
    $t$, two-sample Welch $t$, Mann-Whitney, ANOVA, etc.),
    the transformations applied to the raw measurements (none,
    log, Box-Cox), the covariates included as adjustment, and
    the post-hoc procedures if any (Tukey HSD, Bonferroni).
  \item \textit{Decision rule.} What the engineer will conclude
    from each possible outcome. Example: ``if $p < 0.05$ and
    the lower bound of the $95 \,\%$ CI on the mean difference
    exceeds $0.02\,\si{\mega\pascal}$, supplier B is qualified;
    otherwise the decision is deferred to a second-batch test.''
  \item \textit{Deviations.} A blank section that the engineer
    fills after data collection, listing any deviations from the
    pre-registered plan and the reason for each. Deviations
    convert the affected analysis from confirmatory to
    exploratory; the original pre-registration is preserved.
\end{enumerate}

The completed template is typically a single document of one to
three pages. The commit hash and timestamp of the repository
version serve as the cryptographic record. The engineer who is
asked, two years later, ``did you decide that test before or after
you saw the data?'' can answer with a timestamped log entry.

\paragraph{The CONSORT-style equivalent for engineering reports.}
The Consolidated Standards of Reporting Trials (CONSORT) statement
\cite{paper:ch08t3b-consort-2010} codifies the reporting
requirements for randomised controlled trials in medical research.
The 25-item checklist requires: trial design (including the
pre-registration record), randomisation procedure, blinding,
participant flow (the CONSORT flow diagram counting screened,
enrolled, allocated, analysed, and excluded subjects), primary
and secondary outcomes, statistical methods, all results
(positive and negative), and a discussion of limitations.

The engineering analog is not yet codified into a single
checklist with the authority of CONSORT, but the substance is
identical. The working requirements for an engineering
experimental report:

\begin{itemize}[leftmargin=*]
  \item \textit{All measurements taken.} The report lists every
    specimen, sensor reading, or trial. Specimens that were
    discarded are listed with the reason. The aggregate count of
    discarded specimens is reported; a high discard rate is a
    diagnostic about the experimental setup, not a property to
    be hidden in the body text.
  \item \textit{Everything that didn't work.} Pilot studies that
    failed to produce useful data, sensor channels that lost
    calibration midway through the campaign, experimental
    configurations that were tried and abandoned: all of these
    are reported. Selective omission produces the same
    publication-bias mechanism that drove the replication crisis
    in psychology.
  \item \textit{The pre-registration record.} The report cites
    the pre-registration's timestamp and location (repository
    commit hash, OSF DOI, internal protocol number) and
    reproduces the pre-registered analysis plan verbatim.
    Deviations from the plan are listed in the body of the
    report, not in a supplementary appendix.
  \item \textit{Effect sizes with confidence intervals.} Every
    reported $p$-value carries the corresponding effect size and
    its $95 \,\%$ CI. The four-number minimum from the effect-
    size subsection of section 8.3 applies.
  \item \textit{The IMRaD structure.} Introduction, Methods,
    Results, and Discussion remain the standard sections.
    Methods contains the pre-registered plan and the deviations.
    Results contains all measurements. Discussion contains the
    engineering interpretation, limitations, and the conditions
    under which the result would not generalise.
\end{itemize}

A working report that satisfies these requirements is
roughly twice the length of a journal article that does not.
The space cost is part of the engineering-reporting discipline;
the alternative is a literature that does not reproduce.

\section{P-hacking, garden of forking paths, and what to refuse}\pathtag{core}

A statistically significant result has confirmatory weight only
when the analysis that produced it was specified before the data
were seen \cite[ch.~10]{text:wasserman2004}. The same framework
fails when the analysis plan is chosen after the data are visible.
The American Statistical Association's 2016 statement on $p$-values
\cite{method:ch08-wasserstein-asa-2016} restates the principle in
six points that should be in every engineer's working memory: a
$p$-value measures the compatibility of the data with a specified
null, it does not measure the probability that the null is true,
it does not measure the size of the effect, and scientific
conclusions should not be based on whether a $p$-value passes a
specific threshold. Ioannidis's earlier analysis
\cite{method:ch08-ioannidis-2005} derives the structural
consequence: under modest assumptions about prior plausibility,
power, and analytical flexibility, the majority of published
research findings in many fields are expected to be false. The
mechanisms section 8.7 names are how that result is generated in
practice.

\Cref{fig:vol01:ch08:p-value-distributions} makes the inference
geometry visible. Under the null, $p$ is uniform on $[0, 1]$; under
a well-powered alternative, $p$ piles up near zero. The empirical
$p$-value distribution of a research literature is therefore a
diagnostic: a sharp spike just under $0.05$ with a flat tail above
is the signature of selective reporting.

\begin{figure}[ht]
\centering
\begin{tikzpicture}[x=1cm, y=1cm, font=\small]

% Two-panel: p-value histogram under the null (flat) and under a
% well-powered alternative (right-skewed, mass at small p). Bins on
% [0, 1] in 0.05-wide buckets, ten bins shown.

% ---- Panel A: under H0 (uniform) ----
\begin{scope}[shift={(0,0)}]
  \draw[->, thick] (-0.1, 0) -- (5.5, 0)
    node[below, font=\scriptsize] {$p$-value};
  \draw[->, thick] (0, -0.1) -- (0, 3.4)
    node[left, font=\scriptsize] {frequency};
  \foreach \k in {0, 1, 2, 3, 4, 5, 6, 7, 8, 9} {
    \pgfmathsetmacro\xpos{\k*0.5}
    \pgfmathsetmacro\height{1.4 + 0.20*rand}
    \draw[fill=engblue!50, draw=engblue]
      (\xpos, 0) rectangle (\xpos+0.5, \height);
  }
  \draw[dashed, engred, thick] (0, 1.4) -- (5.0, 1.4);
  \node[font=\scriptsize, engred, anchor=west] at (3.2, 1.6)
    {expected (uniform)};
  \node[anchor=south, font=\small\bfseries] at (2.5, 3.5)
    {(a) Under the null};
  \foreach \k/\lab in {0/0, 5/0.5, 10/1.0} {
    \pgfmathsetmacro\xpos{\k*0.5}
    \node[anchor=north, font=\scriptsize] at (\xpos, -0.05) {\lab};
  }
\end{scope}

% ---- Panel B: under a well-powered alternative (mass at small p) ----
\begin{scope}[shift={(7.0,0)}]
  \draw[->, thick] (-0.1, 0) -- (5.5, 0)
    node[below, font=\scriptsize] {$p$-value};
  \draw[->, thick] (0, -0.1) -- (0, 3.4)
    node[left, font=\scriptsize] {frequency};
  \foreach \k/\h in {0/3.1, 1/1.6, 2/1.0, 3/0.7, 4/0.55, 5/0.45,
                     6/0.40, 7/0.35, 8/0.30, 9/0.25} {
    \pgfmathsetmacro\xpos{\k*0.5}
    \draw[fill=engred!50, draw=engred]
      (\xpos, 0) rectangle (\xpos+0.5, \h);
  }
  \node[anchor=south, font=\small\bfseries] at (2.5, 3.5)
    {(b) Under the alternative};
  \foreach \k/\lab in {0/0, 5/0.5, 10/1.0} {
    \pgfmathsetmacro\xpos{\k*0.5}
    \node[anchor=north, font=\scriptsize] at (\xpos, -0.05) {\lab};
  }
\end{scope}

\end{tikzpicture}
\caption[Distribution of $p$-values under null and alternative]{The
$p$-value's own distribution is the fastest way to see what a
$p$-value is. Under the null, $p$ is uniform on $[0, 1]$ (panel a):
every twenty repetitions of a true-null experiment produces, on
average, one $p < 0.05$ by chance. Under a well-powered alternative
(panel b), the distribution piles up near zero; the height of the
leftmost bin is the power of the test. A literature whose $p$-value
distribution shows a sharp spike just below $0.05$ but flat above
is exhibiting the signature of publication bias and selective
reporting.}
\label{fig:vol01:ch08:p-value-distributions}
\end{figure}


\subsection{P-hacking}

\engterm{P-hacking} is the practice of running multiple analyses on
the same dataset and reporting the one that crosses a significance
threshold \cite{method:simmons-nelson-simonsohn-2011}. The mechanism
is statistical: if the null is true, a test at $\alpha = 0.05$
produces a false positive once in twenty runs. An analyst who tries
twenty different specifications has, on average, one significant
result by chance. Reporting that one without disclosing the other
nineteen is a misleading confirmatory claim.

The same mechanism applies to subgroup analyses. Splitting a sample
by gender, age, region, time of day, condition, and so on, then
running the test in each subgroup, multiplies the number of tests
and the false-positive rate. As the number of subgroups grows, the
probability of at least one false positive approaches one even when
the null holds in every population.

The remedy is multiple-comparison correction (Bonferroni: divide
$\alpha$ by the number of tests; Benjamini-Hochberg: control the
false discovery rate; permutation methods for empirical correction).
The deeper remedy is to avoid generating the comparisons in the first
place: pre-register the analyses; report all results, not the
significant subset; and treat exploratory subgroup analyses as
hypothesis-generating rather than hypothesis-testing.

\subsection{Garden of forking paths}

The \engterm{garden of forking paths} is p-hacking's quieter cousin
\cite{method:gelman-loken-forking-paths}. The analyst does not run
twenty tests; they run one test, but the specification of that test
was determined after seeing the data. At
each branch in the analysis (which transformation, which subset,
which covariate, which test, which correction) the analyst chose the
specification that produced a clean result. The published analysis
is one path; the unwalked paths are invisible.

The forking-paths failure is harder to detect than direct p-hacking
because no count of multiple tests appears in the report. The
analyst genuinely ran only one test. They simply chose, branch by
branch, the specification that worked. The statistical correction
has to be against the entire decision tree of possible analyses,
which is rarely available.

The remedy, again, is pre-registration: write down the analysis
plan, including the transformations, the subsets, the covariates,
and the test, before the data are seen. Branches taken after the
data are seen are flagged as exploratory in the published account.

\subsection{HARKing}

\engterm{HARKing} is hypothesising after the results are known
\cite{method:kerr-harking-1998}: the analyst examines the data,
finds a pattern, and writes the paper as if the pattern was the
prediction. The hypothesis-test machinery
assumes the hypothesis was specified before the data; using the data
to specify the hypothesis breaks the inference, even when the
arithmetic of the test is correct.

The remedy is, once more, pre-registration. A clear separation
between the confirmatory analysis (the pre-registered one) and the
exploratory analysis (everything else, labelled honestly) preserves
the inference for the confirmatory part and recovers honest
hypothesis generation for the exploratory part.

\subsection{Stopping rules}

A test that allows the analyst to stop data collection when the
result becomes significant inflates the false-positive rate. The
mechanism is the same: with each additional sample, the test is
re-run; if any of the running tests crosses the threshold, the
analyst stops. The probability of crossing the threshold at some
point during data collection is much higher than $\alpha$.

Sequential testing methods (Wald's sequential probability ratio
test, group sequential designs, alpha-spending functions) handle
optional stopping correctly. They are standard in clinical trials
and increasingly used in industrial A/B testing. An analyst who
peeks at running results without using such methods is creating an
under-disclosed multiple-testing problem.

\subsection{What an engineer should refuse}

The list of practices an engineer should refuse to perform on their
own data, and refuse to accept from a colleague:

\begin{itemize}[leftmargin=*]
  \item Running multiple tests and reporting the significant one
    without correction or disclosure.
  \item Splitting the sample into subgroups after seeing the data
    and reporting subgroup-level significance.
  \item Choosing the analysis specification (transformation, subset,
    covariate set) after seeing the data.
  \item Stopping data collection when the result becomes
    significant, without using a sequential method.
  \item Reporting the hypothesis as the prediction when the data
    suggested it.
  \item Excluding outliers without a pre-specified rule and reporting
    the result as if the rule had been pre-specified.
  \item Reporting a single $p$-value as evidence of effect without
    reporting the effect size and its confidence interval.
  \item Claiming replication based on a directionally consistent
    result that is itself not significant.
\end{itemize}

The engineer who refuses these practices loses some short-term
results and gains a body of work that survives scrutiny. The cost
is real; the alternative is reports that fall apart when an
independent group tries to use the numbers.

\subsection{The discipline, restated}

\begin{principle}[Pre-registered analysis]
The hypothesis, the test, the significance level, the target power,
the effect size of interest, and the analysis specification are
chosen before the data are seen and recorded in a form that another
engineer can audit. The reported analysis is the pre-registered one.
Departures from the pre-registration are flagged as exploratory and
do not carry confirmatory weight. The reader who cannot find a
pre-registration in a report should treat the reported $p$-values
as exploratory until the pre-registration is produced.
\end{principle}

\section{Failure: replication crises across engineering disciplines}\pathtag{core}

Engineering decisions rest on published numbers. When a body of
published numbers fails to replicate, the engineering decisions
built on it become unreliable. The pattern, current as of 2026, has
appeared in multiple disciplines and is one of the major statistical
failure modes a working engineer must recognize. We discuss the
pattern in general terms; the registry does not yet contain
replication-crisis cases and the chapter avoids naming specific
scientific cases as named accidents.

\subsection{The pattern}

The 2015 Reproducibility Project: Psychology attempted replications
of 100 studies from three psychology journals. Thirty-six percent
of replications produced statistically significant results, and
replication effect sizes were much smaller than original effect
sizes \cite{method:open-science-collaboration-2015}. The 2021
Reproducibility Project: Cancer Biology repeated 50 experiments
from 23 papers and found smaller replication effects and mixed
replication success by multiple criteria
\cite{method:errington-cancer-biology-2021}.
The mechanism is statistical: the published literature is enriched
for false positives because journals selectively publish significant
results, the original studies were under-powered to detect the
true effects with confidence, and analytical flexibility (the
practices section 8.7 names) inflated the apparent effects.

The result, in any field that operates this way, is a literature in
which the reported effect sizes are biased upward, the reported
$p$-values are biased toward significance, and the failure-to-
replicate rate is high. An engineer who reads such a literature and
takes the published numbers at face value designs to numbers that
overstate the strength of the effects.

\subsection{A working measurement-physics case: OPERA's 2011 result}

A measurement-physics replication failure that the working
engineer should know is the OPERA collaboration's 2011 report
of superluminal muon neutrinos. The OPERA detector at Gran
Sasso, Italy, measured the time-of-flight of muon neutrinos
produced by the CERN SPS at $730\,\si{\kilo\meter}$ baseline and
reported an arrival $60.7 \pm 6.9 \,\si{\nano\second}$ earlier
than the light-speed prediction, corresponding to a fractional
velocity excess
$(v - c)/c = (2.48 \pm 0.28) \times 10^{-5}$
\cite{paper:ch08t3p-opera-2011}. The paper passed internal
collaboration review with a stated $6\sigma$ significance and
was posted to arXiv on $22$ September $2011$; the result
attracted intense scrutiny because it contradicted special
relativity.

\paragraph{The two errors found.} An independent OPERA
investigation, completed in February 2012, identified two
opposite-sign instrumental errors. A loose optical-fibre
connector between the GPS receiver and the master clock added
$73.2 \pm 1.5\,\si{\nano\second}$ to the apparent arrival time;
a clock signal-distribution oscillator running fast subtracted
$15\,\si{\nano\second}$ \cite{paper:ch08t3p-opera-2012}. The
net systematic effect explained the entire anomaly; the
corrected measurement was consistent with $v = c$ within the
revised uncertainty budget. OPERA retracted the result and
published the correction in the same journal in July 2012.

\paragraph{What the working engineer takes from the case.} The
$6\sigma$ statistical significance was honest, in the sense that
the data analysis was sound. The failure was in the
\emph{uncertainty budget}: the team had not characterised the
fibre-optic delay variation across temperature and time, and
had not characterised the clock signal-distribution drift. The
GPS-based time link had two systematic sources of error that
were not in the uncertainty model. The lesson generalises:
when a measurement produces a result that contradicts a
well-established prediction, the most likely explanation is
a systematic source not in the uncertainty budget. The OPERA
collaboration's discipline in seeking and finding the errors,
once the result was challenged, is the discipline the
engineering reader should emulate; the failure to seek them
before publication is the failure mode this section names.

\subsection{An engineering case: the carbon-nanotube tensile-strength claim}

A nearer-to-engineering replication failure is the
carbon-nanotube tensile-strength claim trajectory of
$2000$--$2015$. Early single-tube measurements at Rice
University in $1996$ reported individual multi-walled
nanotube tensile strengths of $63 \pm 16 \,\si{\giga\pascal}$
\cite{paper:ch08t3p-yu-rice-2000}, and individual single-
walled-nanotube measurements at IBM and elsewhere reported
values from $20$ to $150 \,\si{\giga\pascal}$. The literature
quickly stabilised on ``nanotube tensile strength
$\sim 100 \,\si{\giga\pascal}$, $100\times$ steel,'' a number
that motivated billions of dollars of nanotube research,
patents, and industrial applications across two decades.

When systematic replication using purified, length-classified
nanotube samples and standardised mounting procedures came in
through 2010--2015, the strength values declined. The 2014
Brock et al.\ campaign \cite{paper:ch08t3p-brock-cnt-2014}
measured a population of length-classified single-walled
nanotubes at $30 \pm 12 \,\si{\giga\pascal}$, a third to a
half of the early figure. The Yakobson group's 2018 modelling
study \cite{paper:ch08t3p-yakobson-cnt-2018} identified the
likely culprit: early measurements suffered from sample
selection bias (the survived-mounting samples were the
strongest of a strength-distributed population), gauge-length
artefacts (shorter samples test stronger due to defect-density
scaling), and grip-induced damage that systematically
under-reported failure stresses near grip ends. Engineering
applications that designed to the headline $100 \,\si{\giga
\pascal}$ figure had no field instrument that could deliver
it. Engineering applications that designed to the meta-
analytically-corrected $30 \,\si{\giga\pascal}$ figure had a
realistic target that field-tested samples could meet.

The case is a quieter cousin of the OPERA story. Both are
about the difference between an effect that the literature
reports and the effect that an independent measurement will
confirm. The engineering response is the same: when a number
will enter load-bearing design, the design uncertainty allows
for the possibility that the original measurement was
optimistic.

\subsection{Engineering exposure}

The mechanism reaches engineering through three doors.

The first is direct: engineering applications that build on
scientific literature inherit the literature's reproducibility
problems. Materials properties from a single study, drug effects
from a single trial, biomarker thresholds from a single biostatistical
model, sensor calibrations from a single test campaign are all
exposed.

The second is indirect: engineering practice itself uses
statistical methods, and the same analytical flexibility that
produces unreplicable scientific results produces unreplicable
engineering results. A failure-rate estimate based on a single
field study, a process-capability index based on a single sample, a
reliability prediction based on a single accelerated test inherits
the same exposure.

The third is institutional: the same incentive structures that
produce publication bias in science (positive results are
publishable, replications are not) appear in engineering with
slight modification. A reliability report that finds an effect is
publishable, presentable, billable; a study that finds no effect is
filed.

\subsection{What an engineer does about it}

Several practices, when adopted, harden engineering against the
replication-crisis pattern.

\engterm{Pre-registration} of hypotheses, sample sizes, and analysis
plans, as section 8.6 and 8.7 specified. The discipline is the same
whether the work is published or internal.

\engterm{Multiple independent replications} before a result enters
load-bearing use. A single study's effect size is an estimate, often
selected from a literature that favours positive results.
Replications and meta-analyses often reduce the effect-size
estimate. The working rule is to size the engineering decision
against a conservative defensible effect size and to state the
source of that estimate.

\engterm{Open data and reproducible analysis}. The data, the code,
the analysis pipeline, and the random seeds are recorded so that an
independent group can rerun the analysis and obtain the same numbers.
Engineering laboratories increasingly use version-controlled analysis
notebooks for this purpose.

\engterm{Effect-size reporting alongside $p$-values}. A $p$-value
without an effect size is a fragment. An effect size with a
confidence interval is the engineering report; the $p$-value is a
secondary indicator.

\engterm{Skepticism toward single-study results}. A single study
rarely deserves to be decisive when the result will enter
load-bearing design. The engineer should ask for independent
replication, a preregistered analysis, or a conservative uncertainty
allowance before treating the number as design input.

\subsection{The principle}

\begin{principle}[Reproducibility is part of the result]
A reported result is incomplete without the data, the analysis
specification, and a record of the steps that an independent
engineer could rerun to obtain the same numbers. A finding that
cannot be reproduced is, for engineering purposes, a hypothesis,
not a result. The discipline that distinguishes the two is the
discipline that prevents an engineering decision from being
silently built on a literature that overstates its effects.
\end{principle}

\section{Historical sketch: from Fisher to the ASA statement on
\texorpdfstring{$p$-values}{p-values}}\pathtag{mastery}

The statistical machinery the chapter develops is itself a
recent intellectual product, and the engineer who reads
classic papers or older textbooks should know which framework
each writer is working in.

\paragraph{Fisher (1925).} Ronald A. Fisher's \emph{Statistical
Methods for Research Workers} \cite{hist:fisher1925}
introduced significance testing as a research-discovery aid:
compute a test statistic, compare its value to the null
distribution, and report a $p$-value as evidence against the
null. Fisher's $p$ was a continuous measure of evidence
strength; the $\alpha = 0.05$ threshold appeared as a practical
convenience for routine agricultural and biological work at
Rothamsted Experimental Station, not as a binary
accept/reject decision rule. Fisher's framework was inductive:
the test was a step in the cycle of hypothesis formation.

\paragraph{Neyman and Pearson (1933).} Jerzy Neyman and Egon
Pearson's joint paper \cite{hist:neyman-pearson-1933}
reframed the test as a decision procedure. Two competing
hypotheses ($H_{0}$ and $H_{1}$) are specified in advance;
the procedure controls the long-run frequencies of Type I
($\alpha$) and Type II ($\beta$) errors. The framework is
deductive and operational: the engineer commits to a decision
rule before the data arrive and accepts the long-run error
rates as the price of consistency. Power, sample-size
calculations, and pre-registration all descend from this
framework, not from Fisher's.

\paragraph{The Fisher-Neyman-Pearson hybrid (1940s-2010s).}
Most working statistics, in engineering and elsewhere, is a
hybrid that pulls $p$-values from Fisher and accept/reject
language from Neyman-Pearson. The hybrid is internally
inconsistent: a $p$-value is not a long-run error rate, and
treating $p < 0.05$ as a binary decision rule misuses Fisher's
continuous measure of evidence. The inconsistency was visible
throughout the twentieth century but rarely confronted; the
textbook standard treated the hybrid as one method and the
working engineer never had to choose between the frameworks.

\paragraph{Tukey (1956-1980).} John W. Tukey's program of
exploratory data analysis
\cite{text:ch08t3p-tukey-eda-1977} reset the balance toward
descriptive techniques (boxplot, stem-and-leaf, robust
estimators) and against the rigid Fisher-Neyman-Pearson
machinery. Tukey's discipline-of-the-eye approach influenced
the engineering culture's preference for visual diagnostics
and effect-size estimates over $p$-values; the boxplot in
§8.1 and the residual plots in §8.4 are Tukey's
contributions.

\paragraph{Ioannidis (2005).} John P. A. Ioannidis's
\emph{Why most published research findings are false}
\cite{method:ch08-ioannidis-2005} argued, on prior-probability
grounds, that a substantial fraction of significance-tested
findings in low-prior-probability fields (drug discovery, gene
association, small-sample preclinical work) are false
positives. The argument generalised the power-and-base-rate
trap to a structural critique of the literature. The paper is
the most-cited article in the PLOS Medicine archive
($> 6{,}000$ citations as of 2024).

\paragraph{The ASA statement on $p$-values (2016).} The
American Statistical Association's first official statement
on $p$-values \cite{method:ch08-wasserstein-asa-2016} crystallised
the working consensus: a $p$-value does not measure the
probability that the studied hypothesis is true, does not
measure the probability that the data were produced by random
chance alone, and a $p$-value or statistical significance does
not measure the size of an effect or the importance of a
result. The statement was followed in 2019 by a special issue
of \emph{The American Statistician}
\cite{paper:ch08t3-wasserstein-2019} that asked
working statisticians to move beyond the $p < 0.05$ threshold;
the editorial argued that retiring the threshold would
``promote a research culture that can embrace uncertainty,
be thoughtful, open, and modest.''

\paragraph{Engineering implication.} The chapter's working
recommendation, that effect sizes with confidence intervals
should travel with $p$-values and that pre-registration should
accompany formal tests, is the post-2016 consensus position.
The reader who encounters older engineering literature
treating $p < 0.05$ as decisive is reading from the
Fisher-Neyman-Pearson hybrid era; the reader who encounters
modern preregistration and effect-size discipline is reading
from the post-ASA-statement consensus. Both will appear in
the working literature for decades.

\section{End-to-end case study: a semiconductor wafer-yield
study}\pathtag{standard}

The chapter's tools, deployed alone, are isolated tests of
specific quantities. A working engineer assembles them into a
single analysis chain that takes raw measurements through
distribution checks, hypothesis tests, multiple-comparison
adjustments, capability indices, and prediction intervals. The
case below is one such chain. The numerical values are realistic
for a contemporary fab and are drawn from the SEMATECH yield-
analysis literature \cite{web:ch08t3b-sematech-yield-2009} and
the Montgomery quality-control text
\cite[ch.~11]{text:ch08-montgomery-spc-2019}.

\paragraph{The scenario.} A $200\,\si{\milli\meter}$ silicon wafer
line produces $300\,\si{\milli\meter}$-equivalent area dies at a
nominal yield of $0.92$ good die per wafer (92 \,\%). The line
runs $10$ batches per shift, $25$ wafers per batch, for a total
of $250$ wafers in the case study. The primary measured outcome
is wafer-level functional yield, defined as the fraction of
on-wafer die that pass the post-fabrication functional test. A
batch is a single lot of wafers processed together through the
critical lithography step.

The engineering question, decomposed into the chapter's tools:

\begin{enumerate}[leftmargin=*]
  \item Is there a statistically detectable batch-to-batch
    difference in yield? (One-way ANOVA on batches; §8.3.)
  \item If so, which batches differ from which? (Tukey HSD post-
    hoc; §8.6.)
  \item Is the within-batch yield variation in statistical
    control? (Shewhart $\bar{X}$ and $R$ charts; §8.4.)
  \item Does the long-term process meet the customer's
    $C_{pk} \geq 1.33$ requirement? (Process capability
    analysis; §8.4.)
  \item What yield should the planning function expect for the
    next batch? (Prediction interval; §8.4.)
  \item What yield should the customer expect across a year of
    production at the current process state? (Tolerance
    interval; §8.4.)
\end{enumerate}

\paragraph{Step 1: distributional check.} The $250$ wafer-yield
observations have $\bar{y} = 0.918$, $s = 0.024$, range
$[0.847, 0.961]$. A Shapiro-Wilk normality test returns
$W = 0.992$, $p = 0.18$; the normality assumption is not
rejected at $\alpha = 0.05$. A Q-Q plot shows the points falling
close to the diagonal line with mild lower-tail compression,
consistent with the bounded nature of the yield variable
(yields cannot exceed $1$, so the upper tail is naturally
truncated). The Gaussian model is acceptable for the analyses
that follow; for production work near the $0.99$ yield ceiling,
the analyst would transform via $\log(y/(1-y))$ (the logit) or
use the beta distribution directly.

\paragraph{Step 2: one-way ANOVA on batches.} The $10$ batches
have batch means $\bar{y}_{1}, \ldots, \bar{y}_{10}$ and a common
within-batch sample size $n_{i} = 25$. The ANOVA table:

\begin{center}
\begin{tabular}{lccccc}
\hline
Source & df & SS & MS & $F$ & $p$ \\
\hline
Batch (between) & 9 & 0.0581 & 0.00645 & 17.4 & $< 0.001$ \\
Within & 240 & 0.0890 & 0.000371 & & \\
Total & 249 & 0.1471 & & & \\
\hline
\end{tabular}
\end{center}

\noindent
The omnibus $F = 17.4$ with $p < 0.001$ rejects the equal-means
null. The variance-explained estimate is
$\eta^{2} = 0.0581/0.1471 = 0.395$: roughly $40 \,\%$ of total
variance is attributable to batch membership. The biased-corrected
omega-squared is
$\omega^{2} = (0.0581 - 9 \cdot 0.000371)/(0.1471 + 0.000371)
\approx 0.371$, only slightly smaller; the bias correction is
modest at $n = 250$.

\paragraph{Step 3: Tukey HSD post-hoc.} With $k = 10$ batches and
$\nu = 240$ within-group degrees of freedom, the studentised range
critical value at $\alpha = 0.05$ is $q_{0.05, 10, 240} \approx
4.47$. The HSD margin is

\[
q_{0.05, 10, 240} \cdot \frac{s_{\text{pooled}}}{\sqrt{n}}
= 4.47 \cdot \frac{\sqrt{0.000371}}{\sqrt{25}}
= 4.47 \cdot \frac{0.0193}{5}
\approx 0.0172.
\]

\noindent
Any pair of batches whose mean yields differ by more than
$0.0172$ is flagged as significantly different at the family-
wise $\alpha = 0.05$. The Tukey HSD output identifies three
batches (Batch 3, Batch 7, Batch 9) whose mean yields are
significantly below the grand mean, and two batches (Batch 1,
Batch 5) significantly above. The engineering follow-up is to
investigate the process history of the low-yield batches: the
lithography step's exposure time, the resist thickness, the
develop time, and any unscheduled tool-maintenance events
between batches.

\paragraph{Step 4: control chart on within-batch variation.}
A Shewhart $\bar{X}$ chart of the batch means, with control
limits at $\bar{\bar{y}} \pm 3 s_{\text{pooled}}/\sqrt{n_{i}}$,
places the limits at $0.918 \pm 3 \cdot 0.0193/5 = 0.918 \pm
0.0116$, or $[0.906, 0.930]$. Two batches (Batch 3 at $0.882$,
Batch 9 at $0.891$) fall below the lower control limit and one
(Batch 1 at $0.940$) above the upper. The Western Electric
zone tests (one of: a point beyond $3\sigma$; two of three
points beyond $2\sigma$ on the same side; four of five points
beyond $1\sigma$ on the same side; eight points in a row on the
same side of the centreline) flag additional patterns. The
out-of-control batches are the same ones the Tukey HSD
identified; the two analyses give consistent diagnostics.

The accompanying $R$ chart of within-batch range with limits at
$D_{3} \bar{R}$ and $D_{4} \bar{R}$ (for $n = 25$, $D_{3} =
0.503$, $D_{4} = 1.497$) shows all batches in control on
within-batch range; the assignable-cause signal is in the batch
means, not the within-batch dispersion.

\paragraph{Step 5: process capability.} The customer
specification is a lower yield limit (LYL) of $0.85$ and no
upper limit (yield cannot exceed $1$, so an upper specification
is meaningless). The capability index for a one-sided
specification is

\[
C_{pk, \text{lower}} = \frac{\bar{\bar{y}} - \text{LYL}}{3 s_{
\text{within}}}
= \frac{0.918 - 0.850}{3 \cdot 0.0193}
= \frac{0.068}{0.058}
\approx 1.17,
\]

\noindent
which falls below the customer's $C_{pk} \geq 1.33$ requirement.
The performance index using the long-term overall standard
deviation $s = 0.024$ is

\[
P_{pk} = \frac{0.918 - 0.850}{3 \cdot 0.024}
= \frac{0.068}{0.072}
\approx 0.94,
\]

\noindent
which is below even the $C_{pk} \geq 1.00$ minimal acceptability
threshold. The gap between $C_{pk}$ and $P_{pk}$ ($1.17$ vs
$0.94$) is the headroom available if the engineer can eliminate
the batch-to-batch drift identified in steps 3 and 4.

\paragraph{Step 6: prediction interval for the next batch.}
The next batch will have a mean yield distributed approximately
$\mathcal{N}(\mu, \sigma^{2}_{\text{between}} + \sigma^{2}_{\text{
within}}/25)$ where $\sigma^{2}_{\text{between}} \approx
(\text{MS}_{\text{batch}} - \text{MS}_{\text{within}})/25 \approx
2.43 \times 10^{-4}$ and $\sigma^{2}_{\text{within}}/25 \approx
1.48 \times 10^{-5}$. The total variance is approximately
$2.58 \times 10^{-4}$, standard deviation $0.0161$. The $95 \,\%$
prediction interval for the next batch's mean yield is

\[
0.918 \pm t_{0.025, 9} \cdot 0.0161 \cdot \sqrt{1 + 1/10}
\approx 0.918 \pm 2.262 \cdot 0.0169
\approx 0.918 \pm 0.038,
\]

\noindent
or $[0.880, 0.956]$. The planning function should budget for
this range; reporting only the point estimate $0.918$ would
understate the operational scatter the planning function will
face.

\paragraph{Step 7: tolerance interval for annual production.}
A two-sided $95 \,\%$ confidence, $99 \,\%$ coverage tolerance
interval on wafer-level yield, using the Howe 1969 procedure
introduced earlier in this chapter \cite{paper:ch08t3p-howe-1969}
at $n = 250$, gives

\[
\bar{y} \pm k(n = 250, p = 0.99, \gamma = 0.05) \cdot s
\approx 0.918 \pm 2.745 \cdot 0.024
\approx 0.918 \pm 0.066,
\]

\noindent
or $[0.852, 0.984]$. With $95 \,\%$ confidence, $99 \,\%$ of
wafer-level yields across the year will fall in this band.
Note that the lower bound $0.852$ is barely above the customer's
$0.85$ specification limit, confirming the capability concern of
step 5: the process is operating with little margin on the
specification.

\paragraph{Integrated engineering action.} The seven steps
produce a single recommendation. The process has detectable
batch-to-batch drift contributing about $40 \,\%$ of total
variance; the Tukey HSD and the $\bar{X}$ chart agree on the
out-of-control batches; the capability indices flag inadequate
margin on the lower specification; the prediction and tolerance
intervals quantify the operational consequences. The
engineering follow-up is sequenced: (i) investigate the assignable
causes for the out-of-control batches (lithography process
parameters, tool maintenance log); (ii) reduce batch-to-batch
variance by tightening the upstream input controls; (iii)
re-run the capability analysis after the corrective actions are
verified; (iv) update the prediction interval and report the
new bounds to the planning function and the customer.

The case study uses every named tool in the chapter and produces
a single integrated recommendation. The tools are inseparable in
practice; a wafer-yield report that runs only the ANOVA without
the post-hoc, or only the capability index without the prediction
interval, is incomplete in a way the customer's audit will
detect.

\section{Project}\pathtag{core}

\begin{project}[Thirty-trial experiment of the reader's choice]
The reader designs and runs a thirty-trial experiment, analyses it
with the methods of this chapter, and writes a fifteen-hundred-word
reflection on what would change if a colleague repeated the work.
The project is the chapter's deliverable and the volume's notebook
entry on statistics.

\textbf{Track}. Analysis with optional household component (hybrid).
The analysis track uses paper, pencil, a spreadsheet or a
programming environment, and an existing dataset (the reader's own
or a public one). The household track adds simple instruments for
data collection.
\textbf{Hazard class}: none for analysis; low if the experiment uses
household instruments.
\textbf{Tools}. Paper, pencil, calculator, a spreadsheet or
programming environment. Optionally, household instruments such as a
kitchen scale, a thermometer, a stopwatch, a tape measure, or a
smartphone with a sensor app.

\textbf{Choosing the experiment}. The reader chooses an experiment
that produces a single quantitative measurement per trial and can
produce thirty trials within a reasonable budget of time. Examples:

\begin{enumerate}[leftmargin=*]
  \item \engterm{Cooking-time variability}. The reader times how long
    it takes to bring a fixed quantity of water to a rolling boil on
    the same stove, repeated thirty times across at least two weeks.
    The hypothesis: cooking time has a population mean within
    $\pm 30 \,\si{\second}$ of a target value, or two stoves give
    the same population mean.
  \item \engterm{Daily commute time}. The reader times their commute
    over thirty working days. The hypothesis: weekday and Friday
    commutes have the same population mean, or the average is within
    a stated tolerance of a published municipal estimate.
  \item \engterm{Smartphone step-counter accuracy}. The reader
    counts their steps manually over thirty walks of fixed length,
    compares to the smartphone reading, and tests whether the
    smartphone is biased.
  \item \engterm{Household appliance energy use}. The reader records
    energy use for thirty cycles of a fixed appliance task (kettle
    boil of fixed water mass, dishwasher cycle, washing machine
    cycle) and tests whether one cycle's mean differs from another's.
  \item Any equivalent experiment that produces a single
    quantitative measurement per trial and has a meaningful
    null hypothesis the reader can pre-register.
\end{enumerate}

\textbf{Pre-registration}. Before collecting any data, the reader
records, in a notebook or repository file:

\begin{itemize}[leftmargin=*]
  \item The hypothesis, written as a null and an alternative.
  \item The test (one-sample $t$, two-sample $t$, chi-square, or
    other) and its significance level $\alpha$.
  \item The target power $1 - \beta$ and the smallest effect size
    that matters for the engineering question, if the reader is
    taking the standard path through section 8.6. For the core
    path, the project uses the fixed sample size $n = 30$ and the
    reader states what effect size that sample can plausibly detect.
  \item The sample size $n$. Core path: $n = 30$ as fixed by the
    project. Standard path: $n$ computed from the rule of thumb in
    section 8.6 if a $t$-test is planned, or from a power
    calculation appropriate to the chosen test.
  \item The analysis specification: transformations applied (if
    any), the rule for excluding outliers (if any), and the
    summary statistics to be reported.
\end{itemize}

The pre-registration is dated and stored before any measurement is
made.

\textbf{Data collection}. The reader collects thirty trials,
records each measurement with date, time, conditions, and any
relevant context. Outlier candidates are noted but not excluded
without applying the pre-registered rule.

\textbf{Analysis}. The reader produces:

\begin{itemize}[leftmargin=*]
  \item Summary statistics: mean, median, standard deviation, IQR,
    range, and the count of any flagged outliers.
  \item A histogram and a boxplot of the data.
  \item A normal-probability plot or a Q-Q plot to check the
    distributional assumption of the chosen test.
  \item The pre-registered hypothesis test, with the test statistic,
    the $p$-value, the effect size, and the engineering conclusion.
  \item A confidence interval for the relevant parameter at
    $95 \,\%$ and a prediction interval for the next observation
    at $95 \,\%$.
  \item If an effect-size confidence interval is required (standard
    path), provide the formula or a software recipe in the project
    notes. Core readers may report a point estimate of the effect
    size with a qualitative discussion of its uncertainty.
\end{itemize}

\textbf{Reflection (1500 words)}. The reader writes a 1500-word
narrative addressing:

\begin{itemize}[leftmargin=*]
  \item What the engineering decision the experiment was designed
    to inform actually was, and how the result changes that
    decision.
  \item Whether the chosen sample size was adequate for the effect
    size that mattered, and whether the result would have been
    different at twice or half the sample size.
  \item Where the data violated the assumptions of the chosen test
    (non-Gaussian shape, dependence between trials, drift over
    the data-collection period) and what alternative analysis
    would have been appropriate.
  \item What a colleague repeating the experiment with their own
    instruments and conditions might find, and what the reader's
    confidence in the result would change if the colleague's
    result disagreed.
  \item Whether the reader's pre-registration anticipated the
    branches in the analysis where forking-paths flexibility could
    have entered, and how the report would have differed if the
    pre-registration had been written after the data instead of
    before.
\end{itemize}

\textbf{Deliverable}. The pre-registration document, the data
record, the analysis output, and the reflection. The general
editor's reference experiment and analysis are published online; the
reader is encouraged to compare \emph{after} completing their own.

\textbf{Time}. Approximately one hour of pre-registration,
collection time depending on the chosen experiment (often spread
across two to four weeks at one or two trials per day), three to
four hours of analysis, and three to five hours of reflection
writing. Total approximately six to ten hours of work, distributed
across the chosen collection period.

\textbf{Phases}. The project decomposes into four phases with
explicit checkpoint outputs:

\begin{itemize}[leftmargin=*]
  \item \textit{Phase 1: design and pre-registration} (1 to 2
    hours). The reader writes the pre-registration document, names
    the test, justifies $\alpha$, $1 - \beta$, $\delta$, and $n$,
    and dates the document. Checkpoint: the pre-registration is
    archived in a notebook or repository commit before any data
    are collected.
  \item \textit{Phase 2: data collection} (2 to 4 weeks at one or
    two trials per day). The reader records each measurement with
    date, time, conditions, and an outlier flag if the
    pre-registered rule triggers. Checkpoint: the raw data file
    contains thirty trials and any deviations from the
    pre-registered protocol are footnoted.
  \item \textit{Phase 3: analysis} (3 to 4 hours). The reader
    produces the summary statistics, histogram, boxplot, Q-Q plot,
    pre-registered test, confidence interval, and prediction
    interval, in that order. Checkpoint: every analysis output is
    reproducible from the data file and a single analysis script
    or spreadsheet.
  \item \textit{Phase 4: reflection and audit} (3 to 5 hours). The
    reader writes the fifteen-hundred-word reflection, then
    rereads the pre-registration alongside the analysis and notes
    every branch where the data could have nudged the analysis
    differently. Checkpoint: the reflection cites at least three
    such branches and states how each was resolved against the
    pre-registration.
\end{itemize}

\textbf{Success criteria}. The project is judged complete when
each of the following is true:

\begin{itemize}[leftmargin=*]
  \item The pre-registration was archived before any measurement
    and has a verifiable timestamp.
  \item The thirty trials are recorded with conditions, and any
    outlier flags follow the pre-registered rule.
  \item The analysis runs end to end on the recorded data and a
    colleague who reads the deliverable can reproduce the
    summary statistics, the test statistic, and the two intervals
    without further instructions.
  \item The reflection names a specific engineering decision the
    experiment was designed to inform and states how the result
    changes that decision.
  \item The reflection identifies at least one assumption of the
    chosen test that the data partially violate, names the
    alternative analysis, and states what the alternative would
    have changed.
\end{itemize}

\textbf{Common failure modes to avoid}. The pre-registration is
written after data collection begins (failure mode of section 8.7);
the outlier rule is invented after seeing the data; the chosen
test ignores a dependence between consecutive trials (commutes on
adjacent days are not independent); only the mean is reported and
the confidence interval is interpreted as a prediction interval.
Each of these failure modes is enumerated in the reflection so the
reader is forced to confront whether their own analysis is exposed
to it.

\textbf{A worked model report.} The block below is a redacted
form of a project deliverable from a working engineer who chose
the cooking-time-variability track. The redaction strips
identifying detail; the structure and the numbers are as
submitted. The reader can use it as a template for their own
deliverable.

\smallskip
\textit{Pre-registration} (timestamped 2024-11-04 17:32 UTC,
commit hash \texttt{a3f12c8}, internal Git repository).
\emph{Hypothesis}: $H_{0}$: the population mean time to bring
$1\,\si{\litre}$ of tap water at $15\,\degC$ to a rolling boil on
the kitchen induction hob, at power setting $9$ of $10$, equals
$240\,\si{\second}$. $H_{1}$: not equal.
\emph{Test}: one-sample two-sided $t$-test at $\alpha = 0.05$.
\emph{Sample size}: $n = 30$ trials, one per evening, $2024
\text{-}11\text{-}05$ through $2024\text{-}12\text{-}04$.
\emph{Effect size of interest}: $\delta = 10\,\si{\second}$
($4 \,\%$ of nominal); rule of thumb at $\sigma \approx
8\,\si{\second}$ (prior estimate from three pilot trials) gives
$n \geq 16 \cdot 64/100 \approx 11$, so $n = 30$ is more than
adequate.
\emph{Outlier rule}: exclude trials where (i) the lid is opened
during heating, or (ii) the room temperature differs from
$22 \pm 2 \,\degC$. Excluded trials are repeated the following
evening. \emph{Reporting}: mean, median, standard deviation,
range, IQR, histogram, boxplot, Q-Q plot, test statistic,
two-sided $p$-value, $95 \,\%$ CI on the mean, $95 \,\%$
prediction interval for the next trial.

\smallskip
\textit{Data}. The thirty measurements (in seconds):
$233, 238, 245, 241, 236, 247, 239, 242, 244, 235, 240, 246,
243, 237, 241, 248, 239, 244, 242, 238, 245, 240, 243, 246,
241, 239, 244, 247, 242, 245$. Two trials were excluded under
the outlier rule (one lid-opening, one room temperature
$25 \,\degC$) and repeated the following evening; the repeats
are included.

\smallskip
\textit{Summary statistics}.
$\bar{x} = 241.8\,\si{\second}$,
median $= 241.5\,\si{\second}$,
$s = 3.71\,\si{\second}$,
range $= [233, 248]$,
IQR $= [239, 244] = 5\,\si{\second}$.
The Q-Q plot shows close-to-linear behaviour with mild upper-tail
weight; Shapiro-Wilk $W = 0.972$, $p = 0.59$; the Gaussian
assumption is acceptable.

\smallskip
\textit{Test}. $t = (241.8 - 240)/(3.71/\sqrt{30}) = 2.66$,
$\text{df} = 29$, two-sided $p = 0.013$. Reject $H_{0}$ at
$\alpha = 0.05$. The mean boil time is significantly above the
$240\,\si{\second}$ target.

\smallskip
\textit{Effect size and intervals}.
Cohen's $d_{\text{one-sample}} = (241.8 - 240)/3.71 = 0.49$, a
medium effect by Cohen's benchmark; $95 \,\%$ CI on $d$ via the
non-central $t$ method is approximately $[0.11, 0.86]$. The
$95 \,\%$ CI on the mean is $241.8 \pm 2.045 \cdot 3.71/\sqrt{30}
= 241.8 \pm 1.39 = [240.4, 243.2]\,\si{\second}$. The $95 \,\%$
prediction interval for the next trial is $241.8 \pm 2.045 \cdot
3.71 \cdot \sqrt{1 + 1/30} = 241.8 \pm 7.71 = [234.1, 249.5]\,
\si{\second}$. The customer (here, the engineer's spouse, who
plans the dinner timing around the boil time) reads the
prediction interval, not the confidence interval, when deciding
when to start the pasta water.

\smallskip
\textit{Engineering conclusion}. The boil time is statistically
above the $240\,\si{\second}$ nominal by an average of $1.8\,
\si{\second}$. The effect size is medium and the CI on the mean
excludes the nominal. The next trial will, with $95 \,\%$
probability, take between $234$ and $250\,\si{\second}$. The
engineering decision (when to start heating water for a
$19{:}00$ dinner) is to set the start time at $18{:}55$, with
two seconds of margin for the upper end of the prediction
interval. The hob's nominal $240\,\si{\second}$ time is mildly
optimistic; the engineer updates the household planning
spreadsheet accordingly.

\smallskip
\textit{What a colleague repeating the work might find}. A
colleague using a different hob, a different pot, or a
different water temperature would see a different mean. The
$\pm 1.4\,\si{\second}$ CI on the mean is much narrower than
the operational variability the colleague will experience; the
result transfers only as a method, not as a number. A more
useful replication would specify the hob model, the pot mass
and material, and the inlet water temperature with calibrated
instruments.

\smallskip
\textit{Reflection on forking paths}. The pre-registration
specified the outlier rule before data collection; two trials
were excluded under it. The reflection identifies three
branches where the analyst could have nudged the analysis: (i)
the choice of the $t$-test over a Wilcoxon signed-rank test
(taken: the Q-Q plot supported Gaussian); (ii) the choice of
the one-sample test over a Welch test against a separate
calibration sample (taken: no calibration sample available);
(iii) the choice of $\alpha = 0.05$ over $\alpha = 0.01$
(taken: pre-registered). The branches are recorded in the
project deliverable so a reader can audit which choices the
data could have influenced.

\smallskip
\textit{Word count}. The full reflection runs $1{,}487$ words
including the pre-registration text and the data summary.

\textbf{Common failure modes in submitted projects.} Across the
hundred or so submissions the general editor has reviewed since
the project was introduced, four failure modes recur:

\begin{itemize}[leftmargin=*]
  \item \textit{Missing pre-registration}. The deliverable
    reports an analysis but no pre-registration is in the
    repository before the data file's first commit. The
    analysis is treated as exploratory in the grading; the
    reflection must explicitly identify which choices were
    made after the data were visible.
  \item \textit{$p$-value reported without effect size}. The
    deliverable lists the test statistic, the $p$-value, and
    a conclusion without computing Cohen's $d$ or its
    equivalent. The reader cannot judge whether the result is
    operationally meaningful. The fix: report the four-number
    minimum (effect size, CI on the effect size, $p$, $n$) for
    every comparison.
  \item \textit{Confidence interval interpreted as prediction
    interval}. The deliverable reports the $\pm 1.4\,
    \si{\second}$ CI on the mean and recommends starting the
    pasta water at $18{:}59$ (two seconds before nominal). The
    reader misunderstands the CI as bounding the next trial.
    The fix: explicitly compute both intervals and state in
    plain English what each one covers.
  \item \textit{Limitations omitted}. The reflection asserts the
    result without acknowledging the conditions under which it
    fails to transfer (different hob, different pot, different
    inlet temperature). A defensible deliverable states the
    boundary conditions and how to detect when they are
    violated.
\end{itemize}

The grading criterion is honest reporting of the four failure
modes in the reflection, not their absence from the analysis.
A deliverable that catches its own failure modes and explains
them is graded above one that asserts a clean analysis.
\end{project}

\section{Exercises}

The exercises test calculation, derivation, estimation, simulation,
design, diagnosis, failure analysis, and judgment.

\subsection*{Calculation}\pathtag{core}

\begin{exercise}
Compute the sample mean, sample median, sample standard deviation
(with the $n - 1$ denominator), range, and IQR of the dataset
$\{12.1, 12.4, 12.0, 12.5, 12.3, 12.2, 12.7, 12.1, 12.4, 12.3\}$
using the quartile convention stated in the solution notes. Report
each statistic to one more decimal place than the data.
\end{exercise}

\paragraph{Solution sketch.}
Sum is $123.0$, so $\bar{x} = 12.30$. Sorted:
$12.0, 12.1, 12.1, 12.2, 12.3, 12.3, 12.4, 12.4, 12.5, 12.7$;
median is $(12.3 + 12.3)/2 = 12.30$. Deviations from $\bar{x}$ are
$-0.2, +0.1, -0.3, +0.2, 0.0, -0.1, +0.4, -0.2, +0.1, 0.0$; sum of
squares is $0.40$, so $s^{2} = 0.40/9 \approx 0.0444$ and $s
\approx 0.211$. Range is $12.7 - 12.0 = 0.70$. With the linear
percentile convention (R's type 7), $Q_{1} \approx 12.125$ and
$Q_{3} \approx 12.400$, so IQR $\approx 0.275$. A Tukey-hinges
convention gives slightly larger IQR; state the convention.

\begin{exercise}
A sample of $n = 16$ resistors has $\bar{R} = 998 \,\si{\ohm}$ and
$s = 12 \,\si{\ohm}$. Compute the $95 \,\%$ confidence interval for
the mean resistance, using the appropriate $t$ critical value.
Report bounds to the same precision as $\bar{R}$.
\end{exercise}

\paragraph{Solution.}
\textit{Setup.} The sample contains $n = 16$ independent resistance
readings drawn from a population with unknown mean $\mu$ and unknown
standard deviation. The chapter's working model treats the parent
distribution as approximately Gaussian; the $95 \,\%$ two-sided
confidence interval for $\mu$ uses Student's $t$-distribution with
$n - 1 = 15$ degrees of freedom.

\textit{Pivot and critical value.} The pivot is
$T = (\bar{R} - \mu)/(s/\sqrt{n}) \sim t_{15}$. The two-sided
$95 \,\%$ critical value is the $0.975$ quantile,
$t_{0.975, 15} \approx 2.131$. A reader without a table can
remember the working values $t_{0.975, 5} \approx 2.78$,
$t_{0.975, 10} \approx 2.23$, $t_{0.975, 15} \approx 2.13$,
$t_{0.975, 30} \approx 2.04$, all approaching $z_{0.975} = 1.96$.

\textit{Computation.} Standard error of the mean
$s/\sqrt{n} = 12/4 = 3.0 \,\si{\ohm}$. Half-width
$h = 2.131 \cdot 3.0 = 6.394 \,\si{\ohm}$, which rounds to
$6.4 \,\si{\ohm}$ at the working precision of the data. The
interval is
\[
998 \pm 6.4 \,\si{\ohm} = [991.6, 1004.4] \,\si{\ohm},
\]
or, rounded to the unit precision of the data,
$[992, 1004] \,\si{\ohm}$.

\textit{Sanity check.} Using the normal approximation
$z_{0.975} = 1.96$ would give half-width
$1.96 \cdot 3.0 = 5.88 \,\si{\ohm}$. The $t$-correction at $n = 16$
adds about $0.5 \,\si{\ohm}$, or roughly $9 \,\%$. The order of
magnitude of the half-width is correct: the standard uncertainty of
the mean is $\sim 0.3 \,\%$ of $\bar{R}$, and a $95 \,\%$ interval
roughly doubles that.

\textit{Common errors.} (i) Forgetting the Bessel correction and
treating $s^{2}$ as a population variance; the working sample
formula already includes it. (ii) Using $z_{0.975}$ at $n = 16$,
which under-covers by enough to matter for tolerance-stack reports.
(iii) Quoting bounds to more decimals than the data justify; the
half-width is known only to within the $t$-distribution rounding,
itself $\approx 1 \,\%$ for $n = 15$.

\textit{Reporting.} \emph{$\bar{R} = 998 \,\si{\ohm}$, $95 \,\%$ CI
$[992, 1004] \,\si{\ohm}$, $n = 16$, Student-$t$ with $15 \,$df.}
The reader of the report can audit every term.

\begin{exercise}
For the resistor sample of the previous exercise, compute the
$95 \,\%$ prediction interval for the next resistor's resistance.
Compare the width to the confidence interval and explain why the
prediction interval is wider.
\end{exercise}

\paragraph{Solution.}
\textit{Setup.} The quantity of interest is the resistance of the
next single resistor drawn from the same population, given the
sample of $n = 16$ with $\bar{R} = 998 \,\si{\ohm}$ and
$s = 12 \,\si{\ohm}$. A prediction interval, not a confidence
interval, is the right object.

\textit{Model.} Treating the parent as Gaussian, the next observation
$R_{n+1}$ is independent of the sample and has variance $\sigma^{2}$;
the sample mean $\bar{R}$ has variance $\sigma^{2}/n$. Their
difference $R_{n+1} - \bar{R}$ has variance $\sigma^{2}(1 + 1/n)$,
so the pivot $(R_{n+1} - \bar{R})/(s \sqrt{1 + 1/n})$ follows
$t_{n-1}$.

\textit{Computation.} The $95 \,\%$ prediction half-width is
\[
h_{\text{PI}} = t_{0.975, 15} \cdot s \cdot \sqrt{1 + \tfrac{1}{n}}
= 2.131 \cdot 12 \cdot \sqrt{17/16}
\approx 2.131 \cdot 12 \cdot 1.0308
\approx 26.36 \,\si{\ohm},
\]
which rounds to $26 \,\si{\ohm}$. The $95 \,\%$ prediction interval
is $[972, 1024] \,\si{\ohm}$.

\textit{Comparison.} The confidence interval from the previous
exercise had half-width $\approx 6.4 \,\si{\ohm}$. The prediction
interval is $26 \,\si{\ohm}/6.4 \,\si{\ohm} \approx 4.1$ times
wider. The ratio is $\sqrt{1 + 1/n}/(1/\sqrt{n})
= \sqrt{n + 1} \approx \sqrt{17} \approx 4.12$ at $n = 16$, which
matches.

\textit{Sanity check on the limits.} As $n \to \infty$, the
prediction half-width tends to $z_{0.975} \cdot s \approx
1.96 \cdot 12 \approx 23.5 \,\si{\ohm}$, while the confidence
half-width tends to zero. The prediction interval cannot shrink
below the population's intrinsic spread.

\textit{Why it is wider.} The confidence interval expresses
uncertainty about a single fixed parameter (the population mean).
The prediction interval includes that uncertainty \emph{plus} the
inherent variability of a single new draw. The two error sources
add in quadrature, so the prediction interval's width is dominated
by the population $\sigma$ once $n$ is moderate.

\textit{Reporting.} \emph{Next-unit resistance $\in [972, 1024]
\,\si{\ohm}$ with $95 \,\%$ coverage, based on $n = 16$, Gaussian
parent assumption, Student-$t$ with $15 \,$df.} A customer who
needs an acceptance bound on every unit should use this interval,
not the CI.

\begin{exercise}
A diagnostic test has sensitivity $90 \,\%$, specificity $95 \,\%$,
and is applied to a population with prior defect rate $5 \,\%$.
Compute the probability that a unit testing positive is actually
defective. Compute the probability after a second independent
positive.
\end{exercise}

\paragraph{Solution.}
\textit{Setup.} A unit drawn from a population with defect prior
$\Pr(H) = 0.05$ is tested by an instrument whose conditional rates
are sensitivity $\Pr(D \mid H) = 0.90$ and specificity
$\Pr(\bar{D} \mid \bar{H}) = 0.95$ (so false-positive rate
$0.05$). The first positive yields $\Pr(H \mid D_{1})$. A second
\emph{conditionally independent} positive (given the unit's true
state) yields $\Pr(H \mid D_{1} \cap D_{2})$.

\textit{Population calculation for the first test.} Imagine
$N = 10{,}000$ units. Defective count $0.05 \cdot 10{,}000 = 500$;
true positives $0.90 \cdot 500 = 450$. Non-defective count
$9{,}500$; false positives $0.05 \cdot 9{,}500 = 475$. Total
positives $450 + 475 = 925$. Posterior
\[
\Pr(H \mid D_{1}) = \frac{450}{925} \approx 0.486.
\]

\textit{Bayes update for the second test.} Use the posterior after
$D_{1}$ as the new prior. Conditional independence means the second
test's sensitivity and specificity carry forward unchanged. Then
\[
\Pr(H \mid D_{1} \cap D_{2})
= \frac{0.90 \cdot 0.486}{0.90 \cdot 0.486 + 0.05 \cdot 0.514}
= \frac{0.4374}{0.4631}
\approx 0.945.
\]

\textit{Contingency table check (second test).} Of $10{,}000$
units, restrict to those that tested positive once: $925$ units, of
which $450$ are defective and $475$ are good. Re-test all $925$
independently. Among the $450$ defective, the second test flags
$0.90 \cdot 450 = 405$. Among the $475$ good, the second test
falsely flags $0.05 \cdot 475 \approx 24$. Total double-positives
$405 + 24 = 429$. Posterior $405/429 \approx 0.944$, matching the
sequential calculation to within rounding.

\textit{Sanity check on the prior shift.} The Bayes factor for a
positive test is $\Pr(D \mid H)/\Pr(D \mid \bar{H})
= 0.90/0.05 = 18$. The odds form of Bayes says
\[
\text{odds}_{\text{post}} = \text{odds}_{\text{prior}} \times
\text{Bayes factor}.
\]
Prior odds $0.05/0.95 \approx 0.0526$; after one positive
$0.0526 \cdot 18 = 0.947$, giving posterior $0.947/1.947 \approx
0.486$. After a second positive $0.947 \cdot 18 = 17.05$, giving
posterior $17.05/18.05 \approx 0.945$. Both rounds match.

\textit{Operational meaning.} The first positive moves belief from
$5 \,\%$ to about $49 \,\%$, which is large enough to investigate
but not large enough to scrap the unit. A second independent
positive raises it to about $95 \,\%$, decision-relevant for almost
any cost ratio. This is the principle of two-step screening in
industrial quality control.

\textit{Common errors.} (i) Multiplying the two posteriors
($0.486^{2}$) instead of updating sequentially. (ii) Treating the
second test as if it had the same prior as the first. (iii)
Forgetting that conditional independence is a substantive
assumption; correlated errors (a shared firmware bug, a common
calibration drift) reduce the second test's informativeness.

\textit{Reporting.} \emph{Single positive: posterior defect
probability $\approx 0.49$. Double positive (conditionally
independent): $\approx 0.94$. Disposition rule: investigate on
single, scrap on double.}

\begin{exercise}
A Poisson process has mean rate $\lambda = 3$ events per minute.
Compute the probability of zero events in a one-minute window, the
probability of at least five events, and the standard deviation of
the count.
\end{exercise}

\paragraph{Solution.}
\textit{Setup.} The count $K$ of Poisson events in a one-minute
window with rate $\lambda = 3$ has PMF
$\Pr(K = k) = \lambda^{k} e^{-\lambda}/k!$ for
$k = 0, 1, 2, \ldots$. We want $\Pr(K = 0)$, $\Pr(K \geq 5)$, and
$\Var(K)$.

\textit{Computation, $\Pr(K = 0)$.} Setting $k = 0$ gives
$\Pr(K = 0) = e^{-3} \approx 0.0498$. A rough sanity check: at the
mean rate $3$ per minute, the distribution should have non-trivial
mass at zero (about one minute in twenty is event-free), which
matches.

\textit{Computation, $\Pr(K \geq 5)$.} The cumulative through
$k = 4$ is
\[
\Pr(K \leq 4) = e^{-3} \sum_{k=0}^{4} \frac{3^{k}}{k!}
= e^{-3} \cdot \left(1 + 3 + \tfrac{9}{2} + \tfrac{27}{6}
  + \tfrac{81}{24}\right).
\]
The bracket evaluates to $1 + 3 + 4.5 + 4.5 + 3.375 = 16.375$, so
$\Pr(K \leq 4) \approx 0.0498 \cdot 16.375 \approx 0.815$. Hence
$\Pr(K \geq 5) \approx 1 - 0.815 = 0.185$.

\textit{Variance.} The Poisson satisfies
$\E[K] = \Var(K) = \lambda$, so
$\sigma = \sqrt{\lambda} = \sqrt{3} \approx 1.73$. The derivation
via the moment-generating function is in the next subsection's
exercise.

\textit{Sanity check via the normal approximation.} For
$\lambda = 3$ the Poisson is moderately skewed, but the normal
approximation
$K \approx \mathcal{N}(\lambda, \lambda)$ with a continuity
correction estimates
\[
\Pr(K \geq 5) \approx \Pr\!\left(Z \geq \frac{4.5 - 3}{\sqrt{3}}\right)
= \Pr(Z \geq 0.866) \approx 0.193,
\]
within $4 \,\%$ of the exact $0.185$. The approximation improves
quickly: at $\lambda = 10$ it matches the exact Poisson to three
significant figures, which suffices for engineering purposes.

\textit{Common errors.} (i) Forgetting the continuity correction
when comparing Poisson to normal; the bare normal at $z = (5 - 3)
/\sqrt{3} = 1.155$ gives $0.124$, off by a third. (ii) Computing
$\Pr(K \geq 5)$ as $\Pr(K > 5) = \Pr(K \geq 6)$. (iii) Writing the
variance as $\lambda^{2}$ by analogy with the squared standard
deviation of a Gaussian.

\textit{Reporting.} \emph{For Poisson rate $\lambda = 3$ per
minute: $\Pr(K = 0) \approx 0.050$, $\Pr(K \geq 5) \approx 0.185$,
$\sigma = \sqrt{3} \approx 1.73$.} The reporting precision matches
three significant figures, the relevant scale for queueing and
inspection design.

\begin{exercise}
A Weibull lifetime distribution has $\beta = 2.0$ and $\eta = 1000
\,\si{\hour}$. Compute the probability that a unit survives to
$500 \,\si{\hour}$, the probability that it survives to $1500 \,
\si{\hour}$, and the time at which the cumulative failure
probability is $50 \,\%$.
\end{exercise}

\paragraph{Solution.}
\textit{Setup.} The Weibull survival function is
$S(t) = \exp(-(t/\eta)^{\beta})$ for $t \geq 0$. With $\beta = 2.0$
and $\eta = 1000 \,\text{h}$, we want $S(500)$, $S(1500)$, and the
median lifetime $t_{m}$ satisfying $S(t_{m}) = 0.5$.

\textit{Computation.} At $t = 500 \,\text{h}$,
$(t/\eta)^{\beta} = (0.5)^{2} = 0.25$, so $S(500) = e^{-0.25}
\approx 0.779$. About $78 \,\%$ of units survive to $500 \,\text{h}$,
or equivalently $22 \,\%$ have failed.

At $t = 1500 \,\text{h}$, $(1.5)^{2} = 2.25$, so $S(1500) = e^{-2.25}
\approx 0.105$. About $10.5 \,\%$ survive past $1500 \,\text{h}$.

The median solves $\exp(-(t_{m}/\eta)^{\beta}) = 0.5$, so
$(t_{m}/\eta)^{2} = \ln 2 \approx 0.693$, giving
$t_{m} = \eta \sqrt{\ln 2} = 1000 \cdot 0.8326 \approx 833 \,\text{h}$.

\textit{Sanity check on shape.} The Weibull at $\beta = 2$ is
right-skewed. The order of $t_{0.1}, t_{m}, \eta, t_{0.9}$ should be
strictly increasing. Solving $S(t_{p}) = 1 - p$ for $p = 0.1$:
$t_{0.1} = \eta\sqrt{-\ln 0.9} = 1000\sqrt{0.105}
\approx 325 \,\text{h}$. And $t_{0.9} = \eta\sqrt{-\ln 0.1}
= 1000\sqrt{2.303} \approx 1518 \,\text{h}$. The order $325, 833,
1000, 1518$ is monotone, confirming the right skew.

\textit{Why the median is not $\eta$.} The scale parameter
$\eta$ is the time at which $1 - e^{-1} \approx 63.2 \,\%$ of units
have failed, by direct evaluation $S(\eta) = e^{-1}$. Only for the
exponential ($\beta = 1$) does the median happen to coincide with
a clean fraction of $\eta$; for any other $\beta$ the median sits
at $\eta(\ln 2)^{1/\beta}$, here at $833 \,\text{h}$ rather than at
$1000 \,\text{h}$. A reliability report that quotes $\eta$ as a
median lifetime mis-states by about $20 \,\%$ at $\beta = 2$.

\textit{Engineering use.} The numbers position the design between
the $B_{10}$ life and the $B_{90}$ life: $B_{10} \approx 325
\,\text{h}$, median $\approx 833 \,\text{h}$, $B_{90} \approx 1518
\,\text{h}$. A maintenance policy that replaces at $B_{10}$
replaces about ten times more often than one that replaces at the
median; the cost-benefit calculation runs against the failure cost
of a $90 \,\%$-survivor at intervention time.

\textit{Common errors.} (i) Treating $\eta$ as a median.
(ii) Forgetting that $\beta$ governs the shape and $\eta$ governs
the scale; raising $\eta$ by a factor of two doubles every quantile,
whereas raising $\beta$ at fixed $\eta$ compresses the failure
distribution. (iii) Reporting the survival probability without the
parameters, which gives the reader no way to compute neighbouring
quantiles.

\textit{Reporting.} \emph{Weibull$(\beta = 2.0, \eta = 1000
\,\text{h})$: $S(500) \approx 0.78$, $S(1500) \approx 0.11$, median
$\approx 833 \,\text{h}$, $B_{10} \approx 325 \,\text{h}$,
$B_{90} \approx 1518 \,\text{h}$.}

\subsection*{Derivation}\pathtag{standard}

\begin{exercise}
Show that the sample mean $\bar{x}$ has expectation equal to the
population mean $\mu$ when the sample is drawn independently from a
distribution with mean $\mu$. Identify a sampling design under which
this identity fails.
\end{exercise}

\paragraph{Solution sketch.}
By linearity of expectation, $\E[\bar{x}] = (1/n) \sum_{i=1}^{n}
\E[X_{i}] = (1/n) \cdot n\mu = \mu$. Independence is not required
for unbiasedness; identical marginals are. The identity fails when
the sampling design is biased: a self-selected sample (a survey
that respondents may decline), size-biased sampling (units chosen
in proportion to a covariate that correlates with $X$), or
truncation (only units above a threshold are recorded). In each
case $\E[X_{i}] \neq \mu$ for the recorded units, so $\E[\bar{x}]
\neq \mu$ even though the arithmetic of $\bar{x}$ is unchanged.

\begin{exercise}
Show that for a Gaussian sample, the standard uncertainty of the
mean is $\sigma/\sqrt{n}$, then derive the corresponding $95 \,\%$
confidence interval expression. State explicitly which two
quantities the interval estimates as the sample size grows.
\end{exercise}

\paragraph{Solution.}
\textit{Setup.} A sample of $n$ independent observations
$X_{1}, \ldots, X_{n}$ is drawn from $\mathcal{N}(\mu, \sigma^{2})$
with both $\mu$ and $\sigma^{2}$ unknown. The estimator of $\mu$ is
the sample mean $\bar{X} = (1/n) \sum X_{i}$. The estimator of
$\sigma^{2}$ is the Bessel-corrected sample variance
$S^{2} = (1/(n-1)) \sum (X_{i} - \bar{X})^{2}$.

\textit{Variance of the sample mean.} Independence gives
\[
\Var(\bar{X}) = \Var\!\left(\frac{1}{n} \sum_{i=1}^{n} X_{i}\right)
= \frac{1}{n^{2}} \sum_{i=1}^{n} \Var(X_{i})
= \frac{n \sigma^{2}}{n^{2}}
= \frac{\sigma^{2}}{n}.
\]
The standard error of the mean is therefore $\sigma/\sqrt{n}$. When
$\sigma$ is replaced by its estimator $S$, the working standard
error becomes $S/\sqrt{n}$.

\textit{Choice of pivot.} For known $\sigma$ the standardised mean
$(\bar{X} - \mu)/(\sigma/\sqrt{n}) \sim \mathcal{N}(0, 1)$ is the
natural pivot. With $\sigma$ unknown, substituting $S$ inflates the
tails because $S$ is itself a random estimator; the correct pivot
is the Student-$t$ pivot
\[
T = \frac{\bar{X} - \mu}{S/\sqrt{n}} \sim t_{n-1},
\]
under the Gaussian-parent assumption. Gosset's $t$-distribution
with $n - 1$ degrees of freedom captures the extra uncertainty
from estimating $\sigma$.

\textit{Confidence interval.} A two-sided $95 \,\%$ interval comes
from inverting the pivot: $\Pr(|T| \leq t_{0.975, n-1}) = 0.95$
gives
\[
\bar{X} - t_{0.975, n-1} \cdot \frac{S}{\sqrt{n}} \leq \mu
\leq \bar{X} + t_{0.975, n-1} \cdot \frac{S}{\sqrt{n}}.
\]
The realised interval is
$\bar{x} \pm t_{0.975, n-1} \cdot s/\sqrt{n}$, with the
procedural-coverage interpretation: in repeated sampling, $95 \,\%$
of such intervals contain $\mu$.

\textit{Asymptotic behaviour.} As $n \to \infty$,
$t_{0.975, n-1} \to z_{0.975} = 1.96$, the standard error
$\sigma/\sqrt{n} \to 0$, and the CI shrinks around $\mu$. The
half-width is $O(1/\sqrt{n})$: to halve it costs four times the
data; to reduce it tenfold costs one hundred times. By contrast,
the prediction interval's half-width
$t_{0.975, n-1} \cdot s \cdot \sqrt{1 + 1/n}$ tends to
$z_{0.975} \cdot \sigma$, the intrinsic population spread. The two
intervals estimate different things: the CI is for $\mu$, the PI
is for the next $X_{n+1}$.

\textit{Engineering takeaway.} The $1/\sqrt{n}$ wall is the central
constraint of statistical precision; pursuing arbitrarily tight
confidence intervals at fixed $\sigma$ is an exercise in
diminishing returns. A working report quotes
$\bar{x} \pm h$, $n$, and $s$, so a reader can re-derive the pivot
and audit the assumptions.

\begin{exercise}
Starting from the variance formula $\Var(X + Y) = \Var(X) + \Var(Y)
+ 2\Cov(X, Y)$ for two random variables, show that the standard
deviation of the sum of $n$ uncorrelated variables with equal
variance $\sigma^{2}$ is $\sigma\sqrt{n}$. Then show that the
standard deviation of the mean is $\sigma/\sqrt{n}$.
\end{exercise}

\paragraph{Solution sketch.}
Uncorrelated implies $\Cov(X_{i}, X_{j}) = 0$ for $i \neq j$.
Iterating the two-variable identity gives $\Var(\sum_{i=1}^{n}
X_{i}) = \sum \Var(X_{i}) = n\sigma^{2}$, so the standard deviation
of the sum is $\sigma\sqrt{n}$. The sample mean is the sum divided
by $n$; $\Var(\bar{x}) = (1/n^{2}) \cdot n\sigma^{2} =
\sigma^{2}/n$, so the standard deviation of the mean is
$\sigma/\sqrt{n}$. The $\sqrt{n}$ on the sum and the $1/\sqrt{n}$
on the mean are the same $\sqrt{n}$ wearing different hats.

\begin{exercise}
For a Poisson distribution with parameter $\lambda$, show that the
mean and variance both equal $\lambda$. Use the
probability-generating function, the moment-generating function,
or the direct sum form. If using the moment-generating function,
state the derivative step explicitly.
\end{exercise}

\paragraph{Solution sketch.}
Using the MGF: $M(t) = \E[e^{tK}] = \sum_{k \geq 0} e^{tk}
\lambda^{k} e^{-\lambda}/k! = e^{-\lambda} \sum (\lambda e^{t})^{k}
/k! = e^{\lambda(e^{t} - 1)}$. First derivative: $M'(t) = \lambda
e^{t} \cdot M(t)$; at $t = 0$, $M'(0) = \lambda$, so $\E[K] =
\lambda$. Second derivative: $M''(t) = (\lambda e^{t})^{2} M(t) +
\lambda e^{t} M(t)$; at $t = 0$, $M''(0) = \lambda^{2} + \lambda$,
so $\E[K^{2}] = \lambda^{2} + \lambda$ and $\Var(K) = \E[K^{2}] -
(\E[K])^{2} = \lambda$.

\begin{exercise}
Show that for a two-sample $t$-test with equal sample sizes and
equal population variances $\sigma^{2}$, the standard error of the
difference of means is $\sigma\sqrt{2/n}$. Use this to motivate
the rule of thumb $n \approx 16 \sigma^{2}/\delta^{2}$ for $\alpha
= 0.05$ and power $0.80$.
\end{exercise}

\paragraph{Solution.}
\textit{Setup.} Two independent Gaussian samples have means
$\bar{X}_{1}, \bar{X}_{2}$ from populations with common variance
$\sigma^{2}$ and equal size $n$ each. The test statistic is the
difference of means, $D = \bar{X}_{1} - \bar{X}_{2}$. The
hypotheses are $H_{0}: \mu_{1} = \mu_{2}$ and
$H_{1}: \mu_{1} \neq \mu_{2}$ at significance $\alpha$ and target
power $1 - \beta$ at true difference $\delta$.

\textit{Variance of the difference.} Independence gives
\[
\Var(D) = \Var(\bar{X}_{1}) + \Var(\bar{X}_{2})
= \frac{\sigma^{2}}{n} + \frac{\sigma^{2}}{n}
= \frac{2 \sigma^{2}}{n}.
\]
The standard error of the difference is therefore
$\sigma \sqrt{2/n}$.

\textit{Critical value.} Under $H_{0}$ the standardised difference
$Z = D/(\sigma \sqrt{2/n}) \sim \mathcal{N}(0, 1)$. The two-sided
test rejects when $|Z| > z_{\alpha/2}$, equivalently when
$|D| > z_{\alpha/2} \cdot \sigma \sqrt{2/n}$.

\textit{Power equation.} Under $H_{1}$ with true difference
$\delta > 0$, the standardised difference has mean
$\delta/(\sigma \sqrt{2/n})$ and unit variance. The probability of
rejecting on the upper tail (the dominant tail when $\delta > 0$)
is
\[
1 - \beta \approx \Pr\!\left(Z > z_{\alpha/2} - \frac{\delta}
{\sigma \sqrt{2/n}}\right).
\]
Setting this equal to $1 - \beta$ requires
$z_{\alpha/2} - \delta/(\sigma\sqrt{2/n}) = -z_{\beta}$, so
\[
\delta = (z_{\alpha/2} + z_{\beta}) \cdot \sigma \sqrt{2/n}.
\]

\textit{Solve for $n$.} Squaring and rearranging,
\[
n = \frac{2 \sigma^{2} (z_{\alpha/2} + z_{\beta})^{2}}{\delta^{2}}.
\]
At $\alpha = 0.05$ (two-sided): $z_{\alpha/2} \approx 1.96$. At
power $0.80$: $z_{\beta} \approx 0.84$. Sum: $1.96 + 0.84 = 2.80$.
Square: $2.80^{2} = 7.84$. Double: $2 \cdot 7.84 = 15.68 \approx 16$.
The rule of thumb is therefore
\[
n \approx \frac{16 \sigma^{2}}{\delta^{2}}.
\]

\textit{Sanity check.} For $\sigma = \delta$ (effect equals one
standard deviation) the formula gives $n \approx 16$ per group, a
familiar working number. For $\sigma = 2\delta$ ($d = 0.5$,
medium effect) it gives $n \approx 64$ per group, matching the
software output for the same parameters.

\textit{When the rule of thumb fails.} The formula assumes Gaussian
parents, equal variances, and a two-sided test at the stated
$\alpha$ and power. Welch's $t$-test (unequal variances) changes
the denominator slightly; one-sided tests halve the constant in
$z_{\alpha/2}^{2}$ to $z_{\alpha}^{2}$; non-Gaussian parents may
require nonparametric methods with larger samples. The exercises
below extend the derivation to the unequal-variance case.

\textit{Engineering use.} The rule of thumb is a planning tool, not
a final calculation. The discipline is to choose $\delta$ from the
engineering question (the smallest effect that matters) rather
than from the sample size that is convenient to collect; choosing
$\delta$ to make $n$ tractable inverts the discipline.

\subsection*{Estimation}\pathtag{core}

\begin{exercise}
Estimate, by an order-of-magnitude calculation, the sample size
needed to detect a $5 \,\%$ change in the mean of a process with
relative standard deviation $10 \,\%$, at $\alpha = 0.05$ and
power $0.80$. State your assumptions and compare to the rule-of-
thumb formula.
\end{exercise}

\paragraph{Solution sketch.}
Standardised effect $d = \delta/\sigma = 0.05/0.10 = 0.5$. Rule of
thumb $n \approx 16/d^{2} = 16/0.25 = 64$ per group. Assumptions:
Gaussian populations, two-sample two-sided $t$-test, equal sample
sizes, equal variances. Exact software value for these settings is
approximately $n \approx 64$ per group; rule of thumb is within
one or two units.

\begin{exercise}
Estimate the probability that the mean of $25$ independent samples
from a Gaussian distribution falls within one standard error of the
true mean. Compare to the same probability for $100$ samples and
explain the relationship to the central limit theorem.
\end{exercise}

\paragraph{Solution sketch.}
$\bar{x} \sim \mathcal{N}(\mu, \sigma^{2}/n)$ exactly for a
Gaussian parent. $\Pr(|\bar{x} - \mu| \leq \sigma/\sqrt{n}) =
\Pr(|Z| \leq 1) \approx 0.683$. The answer is the same for $n =
25$ and $n = 100$ when expressed in standard-error units; the
absolute width of the interval shrinks as $1/\sqrt{n}$ even though
the coverage probability does not. The CLT says this story holds
approximately for non-Gaussian parents at large $n$; here the
parent is Gaussian, so the result is exact.

\begin{exercise}
A factory produces $10{,}000$ units per day with a defect rate
that the engineer estimates is between $0.5 \,\%$ and $1.5 \,\%$.
Estimate the daily defect count and its uncertainty range. State
the distributional assumption you used.
\end{exercise}

\paragraph{Solution.}
\textit{Setup.} Daily production $N = 10{,}000$ units, defect rate
$p \in [0.005, 0.015]$ (a $3 \times$ uncertainty range). We want
the central estimate and uncertainty of the daily defect count $K$.

\textit{Distributional model.} Treating each unit's defect as
independent Bernoulli$(p)$, the daily count
$K \sim \text{Binomial}(N, p)$. For $N$ large and $p$ small, the
Poisson approximation $K \approx \text{Poisson}(\lambda)$ with
$\lambda = Np$ is excellent: the working rule is $N \geq 50$ and
$Np \leq 10$ for the bare Poisson, and both bounds are loose at our
parameters.

\textit{Central estimate.} At $p = 0.010$ (mid-range), $\lambda =
N p = 100$ defects per day. The Poisson standard deviation is
$\sqrt{\lambda} = 10$ defects.

\textit{Two uncertainty sources.} Two sources contribute. The
\emph{rate uncertainty} from $p \in [0.005, 0.015]$ moves $\lambda$
between $50$ and $150$, a swing of $\pm 50 \,\%$. The
\emph{sampling uncertainty} at fixed $p$ has standard deviation
$\sqrt{\lambda} \approx 10$ at $\lambda = 100$, a swing of
$\pm 10 \,\%$. The rate uncertainty dominates by a factor of five.

\textit{Combined range.} Combining naively, the $\pm 2\sigma$ range
at $\lambda = 100$ is roughly $[80, 120]$; widening for the rate
swing yields $[50 - 14, 150 + 25] \approx [40, 175]$ at the
$\pm 2\sigma$ level. A defensible quote for the daily count is
$[50, 150]$ for the central tendency and roughly $[40, 175]$ if
both sources are bracketed.

\textit{Sanity check on independence.} Independence is the
load-bearing assumption. A shift change that affects every unit in
that shift, a tool wear that drifts $p$ during the day, or a batch
of bad raw material that affects a contiguous run all violate
independence and inflate the variance. The diagnostic is the
variance-to-mean ratio of the empirical count series: a ratio
substantially greater than one is over-dispersion and signals that
the Poisson model is missing structure.

\textit{Common errors.} (i) Reporting $\sqrt{Np(1-p)}
\approx \sqrt{99} \approx 9.95$ as the binomial standard deviation
and quoting that as the full uncertainty, missing the rate
uncertainty entirely. (ii) Computing $\sqrt{p}$ at $p = 0.01$ and
treating it as a per-unit number; the per-unit standard deviation
is $\sqrt{p(1-p)} \approx 0.10$ in defect-fraction units.
(iii) Forgetting that $\lambda$ scales with $N$; a different shift
size changes both the expected count and the absolute spread.

\textit{Reporting.} \emph{Daily defect count $\approx 100$
($p = 0.01$), $\pm 2\sigma$ sampling range $[80, 120]$, central
range across plausible $p \in [0.005, 0.015]$ is $[50, 150]$;
Poisson approximation; independence assumed.}

\begin{exercise}
Estimate the half-width of a $95 \,\%$ confidence interval for a
process mean given $n = 10$ samples and a process standard
deviation $s = 0.5 \,\si{\milli\meter}$. Then estimate the same
half-width for $n = 100$. Compare the cost of each sample size to
the half-width reduction.
\end{exercise}

\paragraph{Solution sketch.}
At $n = 10$: $t_{0.975, 9} \approx 2.26$, half-width $2.26 \cdot
0.5/\sqrt{10} \approx 0.36 \,\si{\milli\meter}$. At $n = 100$:
$t_{0.975, 99} \approx 1.98$, half-width $1.98 \cdot 0.5/\sqrt{100}
\approx 0.099 \,\si{\milli\meter}$. Ten times the data buys roughly
$3.6 \times$ the precision (the $\sqrt{n}$ wall) and the small-
sample $t$-correction recedes from $\approx 2.26$ to $\approx 1.98$.
Whether $0.26 \,\si{\milli\meter}$ of additional precision justifies
ten times the data depends on what engineering decision the
interval informs.

\begin{exercise}
Estimate the probability that, in a class of thirty independent
significance tests at $\alpha = 0.05$, at least one rejects a true
null. State the assumption that allows you to compute the
probability and explain why this calculation motivates multiple-
comparison correction.
\end{exercise}

\paragraph{Solution sketch.}
Independence (the load-bearing assumption) gives
$\Pr(\text{no rejection}) = (1 - 0.05)^{30} = 0.95^{30} \approx
0.215$. Probability of at least one rejection $\approx 0.785$.
With thirty independent tests under the global null, the chance of
a false positive somewhere is nearly four in five. Bonferroni
correction lowers the per-test $\alpha$ to $0.05/30 \approx 0.0017$
and brings the family-wise error rate back to $\approx 0.05$. The
calculation generalises to subgroup analyses, exploratory plots,
and dashboard alarms: the cost of multiplicity is paid in
false-positive rate.

\subsection*{Simulation}\pathtag{mastery}

\begin{exercise}
In a programming environment of the reader's choice, simulate
$10^{4}$ samples of size $n = 30$ from a Gaussian distribution with
$\mu = 0$ and $\sigma = 1$. For each sample compute the sample mean.
Plot the empirical distribution of the sample means and compare to
the theoretical Gaussian with mean $0$ and standard deviation
$1/\sqrt{30}$.
\end{exercise}

\paragraph{Solution sketch.}
Theoretical sd of the mean is $1/\sqrt{30} \approx 0.183$.
Empirical sd of $10^{4}$ means should land within a few percent of
$0.183$; the histogram of means looks Gaussian to the eye even at
$n = 30$ because the parent is already Gaussian (each sample mean
is exactly Gaussian for any $n$). The reference companion script
\repopath{code/clt\_skewed.py} runs the equivalent experiment for
an exponential parent; compare the two runs to see what the CLT
adds versus what it does not.

\begin{exercise}
Simulate $10^{4}$ samples of size $n = 30$ from a heavily skewed
distribution (e.g., exponential with rate $1$, or log-normal with
$\mu = 0$ and $\sigma = 1$). Plot the distribution of sample means
and compare to the central limit theorem prediction. Identify the
sample size at which the means become approximately Gaussian.
\end{exercise}

\paragraph{Solution sketch.}
For exponential$(1)$ the parent has mean $1$, sd $1$, skewness $2$.
At $n = 30$ the sample-mean sd matches $1/\sqrt{30} \approx 0.183$
to within Monte Carlo noise, while the empirical skewness of the
sample means is approximately $2/\sqrt{n} \approx 0.37$, still
visible on a Q-Q plot. Approximate Gaussian shape (skewness below
$0.1$) appears around $n \approx 200$ to $400$ for the exponential
parent; lognormal$(0, 1)$ takes considerably longer because its
skewness is $6.18$. The companion script
\repopath{code/clt\_skewed.py} tabulates the convergence.

\begin{exercise}
Simulate an engineering trial in which the true effect size is
zero. Run a two-sample $t$-test on $1000$ simulated trials at
$\alpha = 0.05$ and verify that the false-positive rate is
approximately $0.05$. Then implement p-hacking by, in each trial,
running tests on three post hoc subgroups such as test rig, shift,
and batch, and reporting the smallest $p$-value. Show that the
empirical false-positive rate for the reported $p$-value is much
higher than $0.05$.
\end{exercise}

\paragraph{Solution sketch.}
The single-test FPR converges to $\approx 0.05$ within
$\sqrt{1000} \cdot \sqrt{0.05 \cdot 0.95}/1000 \approx 0.007$
binomial uncertainty. Reporting the smallest of three subgroup
tests inflates the empirical FPR to roughly $0.10$ to $0.14$
(slightly below $1 - (1 - 0.05)^{3} = 0.143$ because the subgroup
tests share data and are not fully independent). The companion
script \repopath{code/p\_hacking\_demo.py} runs the five-path
version with outlier-trim, subgroup-A, subgroup-B, and
log-transform branches added; its empirical FPR for the
minimum-$p$ strategy is in the range $0.15$ to $0.22$.

\begin{exercise}
Simulate the Bayes-rule defect-test calculation of section 8.5 by
generating $10^{6}$ units with a $1 \,\%$ defect prior, applying a
test with $95 \,\%$ sensitivity and $98 \,\%$ specificity, and
empirically computing the posterior probability of defect given a
positive test. Compare to the analytical $0.324$. Then repeat with
a $10 \,\%$ prior and explain how the posterior depends on the
prior at fixed test characteristics.
\end{exercise}

\paragraph{Solution sketch.}
Expected counts: defective $10^{4}$, of which $9{,}500$ are flagged;
good $9.9 \times 10^{5}$, of which $1.98 \times 10^{4}$ are flagged
falsely. Posterior $9{,}500/29{,}300 \approx 0.324$, matching the
analytical answer to three decimal places (Monte Carlo error is
$\approx 0.003$ at $10^{6}$ trials). At $10 \,\%$ prior: $10^{5}$
defective of which $95{,}000$ flagged, $9 \times 10^{5}$ good of
which $1.8 \times 10^{4}$ flagged. Posterior $95{,}000/113{,}000
\approx 0.841$. Lifting the prior from $1 \,\%$ to $10 \,\%$ moves
the posterior from $0.32$ to $0.84$ at fixed test characteristics:
the posterior depends on the prior approximately linearly when the
prior is small and saturates near $1$ when the prior is moderate.

\subsection*{Design}\pathtag{standard}

\begin{exercise}
Design an experimental protocol to determine whether a new
manufacturing line produces units with the same mean weight as the
existing line. State the hypothesis, the test, the significance
level, the target power, the smallest effect size that would
trigger an engineering response, and the sample size per line.
Include the pre-registration document outline.
\end{exercise}

\paragraph{Discussion.}
A defensible response names: $H_{0}: \mu_{\text{new}} =
\mu_{\text{old}}$, $H_{1}$ two-sided; two-sample Welch $t$-test;
$\alpha = 0.05$; power $0.80$ at smallest engineering-meaningful
difference (a fraction of process specification tolerance). Sample
size from the rule of thumb $n \approx 16(\sigma/\delta)^{2}$ per
line; record $\sigma$ from historical line data. The
pre-registration outline includes: hypothesis, test, $\alpha$,
power, $\delta$, $n$, outlier rule, sampling scheme (random across
shifts), and analysis script committed before sampling. Common
weakness: stating $\delta$ that is convenient for $n$ rather than
the actual engineering tolerance.

\begin{exercise}
Design an A/B test for a software-deployment decision, assuming
the metric is a continuous per-user response with known pilot
standard deviation. State the null and alternative, the smallest
effect size that matters, the sample-size calculation using the
two-sample rule of thumb, the duration of the test, and a
pre-specified stopping rule. If the metric remains daily active
users as a count or proportion, move this exercise to mastery and
provide the appropriate model.
\end{exercise}

\paragraph{Discussion.}
A defensible response specifies the metric (e.g.\ per-session
response time), the unit of randomisation (per user, not per
request), $H_{0}$: no mean difference, $\alpha = 0.05$ with a
sequential correction (alpha-spending or group-sequential) since
the test is monitored. Sample size from the two-sample rule of
thumb at the chosen $\delta$. The stopping rule is the alpha-
spending or group-sequential boundary, not a peek-and-stop on
nominal $p < 0.05$. A weak response chooses a Bayesian framing
without specifying the prior or chooses a fixed sample size with
no stopping rule but allows engineers to look at the dashboard,
which is the failure mode of section 8.7.

\begin{exercise}
Design a reliability test for a new component whose target lifetime
distribution is Weibull with $\beta = 2$ and unknown scale $\eta$.
The engineer needs evidence that the design meets a $10{,}000 \,\si
{\hour}$ $B_{10}$ life (the time at which $10 \,\%$ of units have
failed). State the required relationship between $B_{10}$ and
$\eta$, the test duration, censoring plan, and analysis procedure.
Mark this exercise mastery unless the chapter adds the
reliability-test machinery.
\end{exercise}

\paragraph{Discussion.}
$B_{10}$ satisfies $1 - \exp(-(B_{10}/\eta)^{2}) = 0.1$, so $B_{10}
= \eta\sqrt{-\ln(0.9)} \approx 0.3246\,\eta$. To meet $B_{10} \geq
10{,}000 \,\text{h}$ requires $\eta \geq 30{,}800 \,\text{h}$. A
defensible test plan uses type-I censoring at, say, $5{,}000
\,\text{h}$ on $n = 30$ units, estimating $(\beta, \eta)$ from the
censored sample by maximum likelihood (a standard reliability-
software function) and producing a confidence interval for $B_{10}$.
The submission passes if the lower confidence bound on $B_{10}$
exceeds $10{,}000 \,\text{h}$. Weak responses propose
running-until-all-fail (cost-prohibitive) or testing all units to
$10{,}000 \,\text{h}$ and counting failures (statistically wasteful
of censored information).

\subsection*{Diagnosis and reverse engineering}\pathtag{core}

\begin{exercise}
A colleague reports that a process change produced a statistically
significant improvement, $p = 0.04$, with $n = 12$ before and
$n = 10$ after. Identify the information missing from this
report that an engineering decision would require. Propose at
least four follow-up questions.
\end{exercise}

\paragraph{Solution.}
\textit{Setup.} The reported information is: $p = 0.04$ from a
two-sample test, $n_{1} = 12$ before, $n_{2} = 10$ after. The
implicit claim is that the process change produced a real
improvement. Our task is to enumerate what is missing and what to
ask.

\textit{What is missing from the report.}
\begin{enumerate}[leftmargin=*]
  \item \textit{Effect size and its confidence interval.} A
    $p$-value without an effect size cannot inform an engineering
    decision. A $p = 0.04$ at the small $n$ here could correspond
    to a large effect (engineering-relevant) or a borderline one
    (likely a chance fluctuation in a noisy process).
  \item \textit{Pre-registration status.} Was the test specified
    before data collection? If the test was selected after seeing
    the data (forking paths, section 8.7), the nominal $p$ is not
    the correct false-positive rate.
  \item \textit{Power.} What is the power at the
    engineering-meaningful effect size? At $n_{1} + n_{2} = 22$
    the test has limited power for moderate effects; if it just
    crossed threshold, the winner's-curse mechanism makes the
    estimated effect biased upward.
  \item \textit{Units and operational meaning.} What is the
    measured variable, in what units, and what counts as
    ``improvement''? An ``improvement'' of $0.5 \,\%$ may be below
    the noise floor of the downstream process.
  \item \textit{Outlier rule.} Were any observations excluded? If
    so, by what pre-specified rule?
  \item \textit{Distributional check.} Q-Q plot, Shapiro-Wilk, or
    equivalent. A two-sample $t$-test on heavily skewed data is
    not the right tool.
  \item \textit{Equal-variance assumption.} If the standard $t$
    rather than Welch's was used, was the equal-variance
    assumption checked?
  \item \textit{Independence.} Were the $22$ observations drawn
    from the same machine, same shift, same batch? Pseudo-
    replication inflates apparent significance.
\end{enumerate}

\textit{Four follow-up questions (priority-ordered).}
\begin{enumerate}[leftmargin=*]
  \item ``What is the effect size and its $95 \,\%$ confidence
    interval, in the operational units?''
  \item ``Was the test pre-specified before data collection? If
    not, how many other analyses or specifications were considered
    before this one was reported?''
  \item ``What was the power at the smallest engineering-meaningful
    effect, given the $n_{1} = 12, n_{2} = 10$ sample?''
  \item ``Were the before and after samples drawn from comparable
    operating conditions (same shift, same operator, same
    upstream feed)?''
\end{enumerate}

\textit{What a defensible follow-up report contains.} The original
$p$-value, the effect-size estimate with $95 \,\%$ CI, a Q-Q plot
of each sample, a pre-registration if one exists, a statement of
the alternative analyses considered, the power calculation, and a
plain-language statement of how the engineering decision changes
at the lower bound of the effect's CI. Without these, the report
is a hypothesis, not a finding.

\textit{Common errors.} (i) Accepting the $p = 0.04$ at face value
because it crosses $0.05$. (ii) Treating the colleague's report as
adversarial; the discipline is to ask for the missing information
respectfully and treat the response as informative. (iii) Asking
for ``more data'' without specifying which decision the additional
data would inform.

\begin{exercise}
A published study reports a $30 \,\%$ improvement in a sensor
metric with $p < 0.001$ and $n = 8$. State why the small sample
size is concerning even with a small $p$-value, and identify the
specific failure mechanism that small underpowered studies
exhibit when they reach significance.
\end{exercise}

\paragraph{Solution.}
\textit{Setup.} A published study reports a $30 \,\%$ improvement in
a sensor metric, $p < 0.001$, sample size $n = 8$. The intuition
that a small $p$ guarantees a real result needs interrogating in
the context of the small sample.

\textit{Why small $n$ is concerning.} The two-sample rule of thumb
$n \approx 16/d^{2}$ at $\alpha = 0.05$ and power $0.80$ gives the
detectable standardised effect size at $n = 8$ per group:
$d^{2} \approx 16/8 = 2$, so $d \approx 1.4$. To reach $p < 0.001$
at $n = 8$ requires an even larger observed standardised effect.
A truly large effect at the population level produces such results
routinely; a moderate true effect at the population level produces
them only when the sample lands on the upper tail of the
sampling distribution.

\textit{The mechanism: winner's curse.} The conditional distribution
of the estimated effect, given that it cleared a high threshold,
is biased upward. Algebraically, $\E[\hat{d} \mid \hat{d} > c]
> \E[\hat{d}]$ for any threshold $c$. The lower the power at the
true effect, the more aggressive the upward bias. At $n = 8$ and
true effect $d \approx 0.5$ (medium), the power is roughly $0.10$;
the conditional expectation of $\hat{d}$ given $\hat{d}$ reaches
the $p < 0.001$ threshold is several times the true $0.5$.

\textit{Predicted replication behaviour.} The replication at larger
$n$ regresses the estimator toward the true effect. The published
$30 \,\%$ effect from $n = 8$ should be expected to shrink, often
by a factor of $2 \times$ to $3 \times$, in any well-powered
replication. The published $p < 0.001$ is not protective: it
indexes only the sample's location on the noise distribution, not
the calibration of the point estimate.

\textit{What the engineer does.} Treat the $30 \,\%$ result as a
hypothesis for replication, not a calibrated effect size. Size any
downstream design margin against a conservative lower bound on the
effect, not against the point estimate. If the design margin
depends sensitively on the lower bound, commission an independent
replication before committing to the design. The same discipline
applies to vendor specifications, regulatory submissions, and
literature-derived constants used as engineering inputs.

\textit{Sanity check.} The Open Science Collaboration replication
attempts (section 8.8) found replication effect sizes about half
the originals across $100$ psychology studies
\cite{method:open-science-collaboration-2015}. The mechanism is
identical: small-$n$ originals, high publication threshold, upward
bias on selected results. Engineering disciplines face the same
mechanism whenever literature constants come from few-sample
studies.

\textit{Common errors.} (i) Treating the small $p$ as protective.
(ii) Computing the post-hoc power at the reported effect (which is
itself the biased estimator), producing a circular guarantee.
(iii) Quoting a single-study effect as a calibrated number in a
specification, rather than as a hypothesis.

\begin{exercise}
A factory's defect-rate dashboard shows that one of $20$ inspection
metrics crossed a $p < 0.05$ threshold this week. The dashboard
has been running for a year. State whether the result is evidence
of a real change and what additional information would resolve
the question.
\end{exercise}

\paragraph{Solution sketch.}
Twenty metrics times fifty-two weeks gives roughly $10^{3}$
threshold-crossing opportunities under the null; the expected
count of $p < 0.05$ flags is about $50$ per year by chance alone.
A single weekly crossing, especially on a metric that has not
been independently flagged by process knowledge, is consistent
with chance. Resolving the question requires: the false-positive
rate of the dashboard at the chosen threshold (Bonferroni or FDR
corrected), the metric's prior plausibility of change (was there a
process intervention this week?), and whether the metric crosses
again in a follow-up week. Acting on every flag wastes follow-up
budget; ignoring all flags misses real change. A pre-specified
two-week confirmation rule is a practical compromise.

\begin{exercise}
A laboratory reports a confidence interval $[1.20, 1.25]$ for a
process mean and notes that the next ten observations all fell
within this interval. Identify the conceptual error in interpreting
this as a property of the confidence interval and state which
interval should have been used.
\end{exercise}

\paragraph{Solution.}
\textit{Setup.} A laboratory reports a $95 \,\%$ confidence
interval $[1.20, 1.25]$ for the population mean and observes that
ten consecutive new readings fall inside. The colleague treats
this as confirming the CI ``works'' for future observations.

\textit{The conceptual error.} The CI is a statement about a
parameter (the population mean $\mu$), constructed from a single
sample by a procedure with $95 \,\%$ long-run coverage of $\mu$.
The probability that an individual new observation falls in the CI
is a different quantity. Under a Gaussian parent, the half-width
of the CI is $t_{\alpha/2} \cdot s/\sqrt{n}$, while the half-width
of the equivalent PI is $t_{\alpha/2} \cdot s \cdot \sqrt{1 + 1/n}$.
The ratio is $\sqrt{n + 1}$. The CI is narrower by exactly the
factor that makes it inadequate for next-observation coverage.

\textit{Quantitative computation.} If the CI has half-width
$0.025$ ($[1.20, 1.25]$ centred on $1.225$), and $n$ is moderate
(say $n = 50$), then $s/\sqrt{n} = 0.025/2.01 \approx 0.0124$, so
$s \approx 0.088$. The corresponding PI half-width is
$\approx 2.01 \cdot 0.088 \cdot \sqrt{1.02} \approx 0.179$, more
than seven times the CI half-width. The ``correct'' PI is
$[1.046, 1.404]$, not $[1.20, 1.25]$.

\textit{Probability that one observation falls in the CI.} Under
Gaussian assumptions, $\Pr(X \in [\bar{x} - h_{\text{CI}},
\bar{x} + h_{\text{CI}}])$ is approximately
$\Pr(|Z| \leq h_{\text{CI}}/\sigma)$ with $\sigma$ the population
standard deviation. At $h_{\text{CI}}/\sigma \approx 0.025/0.088
\approx 0.28$, this probability is roughly $0.22$, far below the
nominal $0.95$. About one observation in five should fall inside,
not ten in ten.

\textit{Why ten in ten was observed.} Three possibilities:
(i) coincidence (the ten-observation run was an unusually quiet
stretch; this happens with non-trivial probability if the
population is autocorrelated), (ii) the true population spread is
much smaller than $s$ (early estimates of $s$ are themselves noisy,
especially at small $n$), or (iii) the underlying process has
drifted into a low-spread state since the CI was computed (the
CI says nothing about non-stationarity).

\textit{Correct tool.} A $95 \,\%$ \emph{prediction interval} for
the next observation, here $[1.046, 1.404]$, is the right object.
For setting acceptance bounds that should contain a stated
fraction of all production units with stated confidence, a
\emph{tolerance interval} \cite[ch.~7]{text:devore2015} is the
correct object.

\textit{Engineering takeaway.} Three distinct intervals serve three
distinct questions. CI: where is the parameter? PI: where will the
next observation fall? TI: where will (say) $99 \,\%$ of all
production units fall? Mixing them is the most common reporting
error in industrial statistics and is the failure mode behind
specifications that customers read as guarantees.

\textit{Common errors.} (i) Citing the CI as if it bounded
production scatter. (ii) Treating the small-sample $s$ as an
authoritative estimate of $\sigma$, especially when $n < 30$.
(iii) Inferring something about the process from the lucky
observation that ten readings landed inside the CI.

\subsection*{Failure analysis}\pathtag{core}

\begin{exercise}
A reliability engineer reports a Weibull fit to a $20$-unit
lifetime sample with $\beta = 0.7$. The component is known to fail
by fatigue. Argue, in 200 to 300 words, whether the fit is
trustworthy and what physical or statistical reason might produce
$\beta < 1$ in a wear-out process.
\end{exercise}

\paragraph{Solution sketch.}
Fatigue is a wear-out mechanism: the hazard rate increases with
time, which predicts $\beta > 1$ (often $\beta \in [2, 4]$). A
fitted $\beta = 0.7$ contradicts the physics. Plausible
explanations: (i) the sample mixes two failure modes (early
infant-mortality failures from manufacturing defects together
with later fatigue failures), producing a bimodal hazard that the
single-Weibull fit averages into a decreasing rate; (ii) the
sample is heavily censored at the right and the fit sees mostly
early failures; (iii) the $20$-unit sample is too small to
estimate $\beta$ reliably (the standard error of $\beta$ at $n =
20$ is large enough that the $95 \,\%$ CI may include $\beta = 2$);
(iv) the fitting code has a sign or convention error. Next steps:
inspect the lifetime distribution for bimodality, partition the
sample by failure mode, increase the sample size, and report the
confidence interval on $\beta$ alongside the point estimate. A
$\beta$ that disagrees with the physics is a diagnostic that the
fit, the sample, or the failure-mode model is wrong.

\begin{exercise}
A widely-cited engineering result, originally published with $n =
15$ and $p = 0.03$, has been used as design input by multiple
teams. A replication attempt with $n = 100$ finds an effect that
is one-third the size of the original and not statistically
significant at $\alpha = 0.05$. Argue, in 200 to 300 words, what
the design teams should do, and identify the specific statistical
mechanism that produced the disagreement.
\end{exercise}

\paragraph{Solution sketch.}
The mechanism is the winner's curse of underpowered original
studies. At $n = 15$, only a sample that lands on the upper tail
of the effect's true distribution reaches $p = 0.03$; the
reported effect is therefore biased upward. The $n = 100$
replication, with much higher power, estimates closer to the true
(smaller) effect. The original report and the replication are not
contradictory; the original is the noisy upper-tail draw and the
replication is the calibrated estimate. Design teams should: (i)
re-derive any safety margin or capacity claim using the replication
effect size, not the original; (ii) if the replication's effect
size is engineering-relevant at its lower confidence bound,
preserve the original design; (iii) if not, redesign or commission
a second independent replication before reverting. The general
discipline is to size design margins against a defensible lower
bound on the effect, not against the published point estimate.

\begin{exercise}
A regulatory submission cites a single study to support a safety
margin. The study reports a one-sided $t$-test at $p = 0.04$ with
$n = 12$. Argue, in 200 to 300 words, whether the regulator should
accept the submission, and identify the discipline (pre-
registration, replication, effect-size reporting) that would
strengthen the submission.
\end{exercise}

\paragraph{Solution sketch.}
A regulator should not accept a single-study, $n = 12$, one-sided
$p = 0.04$ submission for a safety claim. The combination of low
$n$, one-sided test (which is half as strict as a two-sided test
at the same nominal $\alpha$), and a $p$ just below threshold is
the signature pattern of underpowered studies that selectively
report the favourable direction. The submission lacks: a
pre-registration that fixes hypothesis, test, and analysis ahead
of time; independent replication; effect-size reporting with a
confidence interval; and a discussion of the safety margin's
sensitivity to the lower bound of the effect's CI. Strengthened
submission: a pre-registered confirmatory study with $n$ sized for
$80 \,\%$ power at the safety-critical effect; an effect-size
estimate with a confidence interval whose lower bound exceeds the
safety threshold; ideally an independent replication. The
regulator's job in the present submission is to refuse on
statistical grounds and request the strengthened form.

\subsection*{Judgment}\pathtag{standard}

\begin{exercise}
A senior colleague reports a $p = 0.049$ result and proposes
publishing it as evidence of a real effect. The data were collected
from three different test rigs and the analysis dropped two
outliers. Write a one-paragraph response that takes the colleague
seriously, identifies the missing information, and proposes a
minimum acceptable form for the report.
\end{exercise}

\paragraph{Discussion.}
A defensible response acknowledges the colleague's effort,
identifies the structural concerns (three rigs is a multilevel
sampling structure that requires a mixed-effects analysis or an
explicit per-rig sensitivity analysis; dropped outliers require a
pre-specified rule and a reported analysis with the outliers
included as a sensitivity check), and proposes a minimum
acceptable form: pre-registered hypothesis (or, failing that, an
honest exploratory label), the effect size with a confidence
interval, an analysis with and without the dropped outliers,
per-rig effect estimates, and a discussion of whether the
engineering decision changes at the lower bound of the effect's CI.
The tone respects the colleague's professional standing while
naming the specific deficiencies. A response that simply says
``$p = 0.049$ is not enough'' misses the chance to give the
colleague a concrete path to a defensible report.

\begin{exercise}
An engineer is asked to analyse a dataset that has already been
collected. They notice that one analysis specification gives
$p = 0.04$ and another gives $p = 0.12$. Argue, in 200 to 300
words, which $p$-value the engineer should report and what
disclosures the report must contain.
\end{exercise}

\paragraph{Discussion.}
The engineer should report both, with full disclosure of which was
the pre-specified analysis (if any) and which was the alternative
specification. If the pre-specified analysis is unavailable
(retrospective dataset), the report must declare the analysis as
exploratory and present a sensitivity table covering the
alternative specifications the analyst considered. Selecting and
reporting only the lower $p$-value without disclosing the
alternative is the garden-of-forking-paths failure of section 8.7
in its most direct form. Selecting and reporting only the higher
$p$-value is equally distorting. A defensible report contains: the
specifications considered, the criteria used to choose among them,
the $p$-values from each, the effect sizes from each, and a
plain-language statement that the analysis was exploratory. The
engineering decision should then be sized against the more
conservative result, with any confirmatory claim deferred to a
pre-registered follow-up.

\begin{exercise}
Argue, in one paragraph, whether engineering training should
emphasise frequentist or Bayesian methods. Take a position and
defend it with a specific example from your own experience or a
documented case.
\end{exercise}

\paragraph{Discussion.}
A defensible response takes a position rather than hedging. One
defensible position: training should emphasise frequentist methods
first (Bessel correction, $t$-tests, confidence intervals,
$p$-values, power calculations) because these are the language of
specification documents, regulatory submissions, and process-
control software; Bayesian methods are then introduced as a more
flexible framework for sequential updating, sensor fusion, and
small-sample decisions where prior information is genuine. Another
defensible position: start with Bayes because the conditional-
probability reasoning is more transparent and the
hypothesis-testing machinery can be taught later as a specialised
decision rule. Either position is acceptable if it engages with
the concrete example (a specific calibration history, a specific
A/B test, a specific reliability prediction) and acknowledges the
operational facts about which frameworks are present in the
engineer's working environment.

\begin{exercise}
A vendor reports a confidence interval for a sensor's accuracy and
the customer interprets it as the range within which the next
reading will fall. Write a 200-word note to the vendor proposing a
report format that prevents the customer's interpretation while
preserving the inferential content the vendor wants to communicate.
\end{exercise}

\paragraph{Discussion.}
A defensible note: explain the confusion (confidence intervals are
for the parameter, not the next observation), recommend that the
vendor report both a confidence interval for the calibration
parameter \emph{and} a prediction interval (or a tolerance
interval, if specifying acceptance bounds) for the next reading,
suggest that the report state in plain language what each interval
covers (``$95\,\%$ of long-run averages'' vs ``$95\,\%$ of next
readings''), and propose a worked numerical example in the
documentation. The note should be respectful of the vendor's
expertise while making the customer's misinterpretation a design
constraint on the report, not a customer error to be corrected.

\begin{exercise}
A reader's pre-registered hypothesis test fails to reject the null
($p = 0.42$) at the planned sample size. The reader is tempted to
add ten more samples and re-test. Argue, in 200 to 300 words,
whether the reader should add the samples, what the addition does
to the false-positive rate, and what the principled alternative is.
\end{exercise}

\paragraph{Discussion.}
Adding samples and re-testing without a sequential method inflates
the false-positive rate above the pre-registered $\alpha$. The
peek-and-extend strategy is a special case of optional-stopping
bias: every additional test gives the random fluctuation another
chance to cross the threshold. Principled alternatives: (i) accept
the result as a non-rejection and update the engineering decision
accordingly; (ii) commission a fresh pre-registered study at a
larger sample size if the engineering question still warrants
investigation; (iii) use a group-sequential design or
alpha-spending function from the start if optional stopping was
anticipated. Honest reporting: if the reader proceeds with the
extension, the result is labelled exploratory in the report, and
the original pre-registration is preserved as the primary
analysis. The reader's instinct to extend the sample is
understandable; the discipline is to recognise that the
extension's $p$-value cannot carry the confirmatory weight the
pre-registration was designed to support.
